% Lectura y comentario : Internet y las NTIC ; traduction : expression de la préférence et de l'habitude — Espagnol Tle A — Cours signé OnBuch+
% Programme officiel MINESEC. Compiler avec : tectonic E28-internet-ntic-preferencia-habito.tex
\newcommand{\DOCMATIERE}{Espagnol}
\newcommand{\DOCNIVEAU}{Tle A}
\newcommand{\DOCMODULE}{Módulo 5 : Médias, communication et culture de la paix}
\newcommand{\DOCLECON}{E28}
\newcommand{\DOCTITRE}{Lectura y comentario : Internet y las NTIC ; traduction : expression de la préférence et de l'habitude}
% OnBuch+ — préambule « cours signés » ENRICHI (compilé avec tectonic / XeLaTeX)
% Système visuel « Pop » d'OnBuch+ : fond crème (jamais blanc), bordures encre
% épaisses, ombres « dures » décalées, titres Archivo Black en capitales.
% Version MATHÉMATIQUES : graphes (pgfplots/tikz), boîtes pédagogiques
% (définition, propriété, méthode, attention, le savais-tu, exercice résolu,
% exercice, TP, bilan), page de garde et signature OnBuch+ ancrée en bas.
%
% Macros attendues dans main.tex AVANT \input{preamble} :
%   \DOCMATIERE  \DOCNIVEAU  \DOCMODULE  \DOCLECON (numéro)  \DOCTITRE
\documentclass[11pt]{article}

\usepackage{fontspec}
\usepackage[spanish,french]{babel} % français = langue principale ; l'espagnol via \esp{..} / espagnol
\usepackage[a4paper,margin=2cm,top=3.4cm,bottom=2.8cm,headheight=1.4cm,footskip=1.3cm]{geometry}
\usepackage{amsmath,amssymb,amsfonts}
\usepackage{mathtools} % \xmapsto, \coloneqq... souvent utilisés par les modèles
\usepackage{cancel} % simplifications barrées (bilans de forces)
% \abs{z} (module, valeur absolue) et \norm{u} (norme), utilisés par les modèles
\providecommand{\abs}{}\let\abs\relax\DeclarePairedDelimiter{\abs}{\lvert}{\rvert}
\providecommand{\norm}{}\let\norm\relax\DeclarePairedDelimiter{\norm}{\lVert}{\rVert}
\usepackage[table]{xcolor} % [table] : fournit aussi \rowcolors (alternance de lignes)
\usepackage{pagecolor}
\usepackage{tikz}
\usetikzlibrary{arrows.meta,positioning,calc,shapes.geometric,shapes.misc,shapes.arrows,decorations.pathmorphing,decorations.pathreplacing,decorations.markings,patterns,shadows,mindmap,trees,fit,backgrounds,matrix,intersections,angles,quotes}
\usepackage{pgfplots}
\usepackage[version=4]{mhchem} % équations nucléaires \ce{^{235}_{92}U}
\usepackage[european,siunitx]{circuitikz} % schémas électriques
\pgfplotsset{compat=1.18}
\usepgfplotslibrary{fillbetween,groupplots} % aires entre courbes (calcul intégral)
\usepackage{enumitem}
\usepackage{fancyhdr}
% [most] charge toutes les bibliothèques tcolorbox (listings, minted, raster,
% documentation...) et gonfle inutilement la mémoire TeX sur un document long
% (~300 boîtes) : on ne charge que ce qui est réellement utilisé.
\usepackage[skins,breakable]{tcolorbox}
\usepackage{graphicx}
\usepackage{colortbl}
\usepackage{booktabs}
\usepackage{tabularx}
\usepackage{diagbox} % case d'en-tête diagonale des tableaux à double entrée
% Colonnes extensibles centrée / gauche / droite (C, L, R), souvent utilisées par les modèles
\newcolumntype{C}{>{\centering\arraybackslash}X}
\newcolumntype{L}{>{\raggedright\arraybackslash}X}
\newcolumntype{R}{>{\raggedleft\arraybackslash}X}
\usepackage{array}
\usepackage{multicol}
\usepackage{siunitx}
% parse-numbers=false : accepte aussi \SI{2,5 \times 10^{-3}}{...}
\sisetup{locale=FR,output-decimal-marker={,},group-minimum-digits=4,parse-numbers=false}
\DeclareSIUnit\ohm{\ensuremath{\symup{\Omega}}} % U+2126 absent de la police
\DeclareSIUnit\division{div} % réglages d'oscilloscope (ms/div, V/div)
\DeclareSIUnit\tr{tr} % tours (vitesse de rotation, ex. \SI{1500}{\tr\per\minute})
\DeclareSIUnit\molar{M} % concentration molaire (ex. \SI{3}{\molar}, \SI{1}{\micro\molar})
\DeclareSIUnit\an{an} % année (le modèle confond parfois avec \year, un compteur TeX) — ex. \SI{5}{\kilo\meter\per\an}
\DeclareSIUnit\FCFA{FCFA} % franc CFA (ex. \SI{500}{\FCFA\per\kilo\gram})
\DeclareSIUnit\degreeBrix{\textdegree Brix} % teneur en sucre d'un sirop/jus (réfractométrie)
\DeclareSIUnit\month{mois} % le modèle utilise parfois \month, un compteur TeX, comme unité de durée
\usepackage{qrcode}
% \degree : commande fréquemment utilisée par les modèles pour un angle ou une
% température en dehors de \SI{}{} (non fournie par les paquets chargés).
\providecommand{\degree}{^{\circ}}
% \qed (fin de démonstration), utilisé par les modèles sans amsthm
\providecommand{\qed}{\hfill$\square$}
% \barre : double barre des valeurs interdites dans un tableau de variations (tkz-tab)
\providecommand{\barre}{\ensuremath{\|}}
% environnement proof (amsthm non chargé) utilisé par les modèles
\ifdefined\proof\else\newenvironment{proof}[1][Démonstration]{\par\noindent\textit{#1.}\ }{\qed\par}\fi
\usepackage{lastpage}
\usepackage{hyperref}
\hypersetup{hidelinks}

% ============================================================
% Palette « Pop » OnBuch+ — mêmes hex que la classe P (plus_ui.dart)
% ============================================================
\definecolor{poporange}{HTML}{F59321}
\definecolor{popdark}{HTML}{E07A0C}
\definecolor{popcream}{HTML}{FFF6E8}
\definecolor{popink}{HTML}{1C1714}
\definecolor{poppurple}{HTML}{7A5AE0}
\definecolor{popgreen}{HTML}{2E9E6A}
\definecolor{poppink}{HTML}{E0457B}
\definecolor{popblue}{HTML}{2F7DE1}
\definecolor{popgold}{HTML}{C98A12}
\definecolor{popmuted}{HTML}{8A7F76}
\definecolor{popline}{HTML}{EADFD2}
\definecolor{popgray}{HTML}{8A7F76}
\colorlet{popgrey}{popgray}
% Alias tolérants (le modèle nomme parfois les couleurs différemment)
\colorlet{popwhite}{white}
\colorlet{popblack}{popink}
\colorlet{popred}{poppink}
\colorlet{popyellow}{popgold}
% Teintes claires (fonds de tableaux/encadrés) — le modèle nomme parfois
% les couleurs avec un suffixe L (ex. popblueL) sans les avoir définies.
\colorlet{poporangeL}{poporange!20}
\colorlet{popdarkL}{popdark!20}
\colorlet{poppurpleL}{poppurple!20}
\colorlet{popgreenL}{popgreen!20}
\colorlet{poppinkL}{poppink!20}
\colorlet{popblueL}{popblue!20}
\colorlet{popgoldL}{popgold!20}
\colorlet{popredL}{popred!20}
\colorlet{popgrayL}{popgray!20}
% Teintes claires pour fonds de tableaux / zones
\colorlet{popblueL}{popblue!10!white}
\colorlet{poporangeL}{poporange!14!white}
\colorlet{popgreenL}{popgreen!12!white}
\colorlet{poppurpleL}{poppurple!12!white}
\colorlet{poppinkL}{poppink!12!white}
\colorlet{popgoldL}{popgold!14!white}

\pagecolor{popcream}

% ============================================================
% Typographie
% ============================================================
\defaultfontfeatures{Ligatures=TeX}
\setmainfont{PlusJakartaSans-Medium.ttf}[Path=../../../fonts/,BoldFont=PlusJakartaSans-Bold.ttf]
% Maths en sans-serif (Fira Math) avec chiffres et lettres droites de Plus
% Jakarta, pour que les formules se fondent dans le texte.
\usepackage[math-style=ISO,bold-style=ISO]{unicode-math}
\setmathfont{FiraMath-Regular.otf}[Path=../../../fonts/]
\setmathfont{PlusJakartaSans-Medium.ttf}[Path=../../../fonts/,range={up/num,up/{Latin,latin},\mathup}]
\usepackage{icomma} % virgule décimale sans espace en mode math
% Caractères mathématiques Unicode absents des polices de texte (Plus Jakarta,
% Archivo Black) : rendus par leur équivalent mathématique, en texte comme en
% formule — sinon ils s'affichent en carré vide (ex. « Arithmétique dans ℤ »).
\usepackage{newunicodechar}
\newunicodechar{ℝ}{\ensuremath{\mathbb{R}}}
\newunicodechar{ℤ}{\ensuremath{\mathbb{Z}}}
\newunicodechar{ℂ}{\ensuremath{\mathbb{C}}}
\newunicodechar{ℕ}{\ensuremath{\mathbb{N}}}
\newunicodechar{ℚ}{\ensuremath{\mathbb{Q}}}
\newunicodechar{𝔻}{\ensuremath{\mathbb{D}}}
\newunicodechar{α}{\ensuremath{\alpha}}
\newunicodechar{β}{\ensuremath{\beta}}
\newunicodechar{γ}{\ensuremath{\gamma}}
\newunicodechar{δ}{\ensuremath{\delta}}
\newunicodechar{ε}{\ensuremath{\varepsilon}}
\newunicodechar{θ}{\ensuremath{\theta}}
\newunicodechar{λ}{\ensuremath{\lambda}}
\newunicodechar{μ}{\ensuremath{\mu}}
\newunicodechar{σ}{\ensuremath{\sigma}}
\newunicodechar{τ}{\ensuremath{\tau}}
\newunicodechar{φ}{\ensuremath{\varphi}}
\newunicodechar{ψ}{\ensuremath{\psi}}
\newunicodechar{ω}{\ensuremath{\omega}}
\newunicodechar{Σ}{\ensuremath{\Sigma}}
% majuscules produites par \MakeUppercase dans les titres (α-aminés -> Α)
\newunicodechar{Α}{\ensuremath{\alpha}}
\newunicodechar{Β}{\ensuremath{\beta}}
\newunicodechar{Γ}{\ensuremath{\Gamma}}
\newunicodechar{Δ}{\ensuremath{\Delta}}
\newunicodechar{Ε}{\ensuremath{\varepsilon}}
\newunicodechar{Θ}{\ensuremath{\Theta}}
\newunicodechar{Λ}{\ensuremath{\Lambda}}
\newunicodechar{Μ}{\ensuremath{\mu}}
\newunicodechar{Τ}{\ensuremath{\tau}}
\newunicodechar{Φ}{\ensuremath{\Phi}}
\newunicodechar{Ψ}{\ensuremath{\Psi}}
\newunicodechar{Ω}{\ensuremath{\Omega}}
\newunicodechar{↦}{\ensuremath{\mapsto}}
\newunicodechar{∈}{\ensuremath{\in}}
\newunicodechar{∉}{\ensuremath{\notin}}
\newunicodechar{∧}{\ensuremath{\wedge}}
\newunicodechar{⇔}{\ensuremath{\Leftrightarrow}}
\newunicodechar{⟺}{\ensuremath{\Longleftrightarrow}}
\newunicodechar{⇒}{\ensuremath{\Rightarrow}}
\newunicodechar{⟹}{\ensuremath{\Longrightarrow}}
\newunicodechar{∘}{\ensuremath{\circ}}
\newunicodechar{∪}{\ensuremath{\cup}}
\newunicodechar{∩}{\ensuremath{\cap}}
\newunicodechar{≡}{\ensuremath{\equiv}}
\newunicodechar{⊂}{\ensuremath{\subset}}
\newunicodechar{⊥}{\ensuremath{\perp}}
\newunicodechar{∃}{\ensuremath{\exists}}
\newunicodechar{∀}{\ensuremath{\forall}}
\newunicodechar{∛}{\ensuremath{\sqrt[3]{\,}}}
\newunicodechar{∜}{\ensuremath{\sqrt[4]{\,}}}
\newunicodechar{⃗}{\ensuremath{{}^{\to}}}
\newunicodechar{ⁿ}{\ifmmode^{n}\else\textsuperscript{\ensuremath{n}}\fi}
\newunicodechar{ˣ}{\ifmmode^{x}\else\textsuperscript{\ensuremath{x}}\fi}
\newunicodechar{ᵇ}{\ifmmode^{b}\else\textsuperscript{\ensuremath{b}}\fi}
\newunicodechar{ʳ}{\ifmmode^{r}\else\textsuperscript{\ensuremath{r}}\fi}
\newunicodechar{ᵘ}{\ifmmode^{u}\else\textsuperscript{\ensuremath{u}}\fi}
\newunicodechar{ʸ}{\ifmmode^{y}\else\textsuperscript{\ensuremath{y}}\fi}
\newunicodechar{ᵃ}{\ifmmode^{a}\else\textsuperscript{\ensuremath{a}}\fi}
\newunicodechar{ᵉ}{\ifmmode^{e}\else\textsuperscript{\ensuremath{e}}\fi}
\newunicodechar{ᵏ}{\ifmmode^{k}\else\textsuperscript{\ensuremath{k}}\fi}
\newunicodechar{ᵖ}{\ifmmode^{p}\else\textsuperscript{\ensuremath{p}}\fi}
\newunicodechar{ᴺ}{\ifmmode^{N}\else\textsuperscript{\ensuremath{N}}\fi}
\newunicodechar{ᶜ}{\ifmmode^{c}\else\textsuperscript{\ensuremath{c}}\fi}
\newunicodechar{ʰ}{\ifmmode^{h}\else\textsuperscript{\ensuremath{h}}\fi}
\newunicodechar{ᵠ}{\ifmmode^{q}\else\textsuperscript{\ensuremath{q}}\fi}
\newunicodechar{ₙ}{\ifmmode_{n}\else\textsubscript{\ensuremath{n}}\fi}
\newunicodechar{ₐ}{\ifmmode_{a}\else\textsubscript{\ensuremath{a}}\fi}
\newunicodechar{ᵢ}{\ifmmode_{i}\else\textsubscript{\ensuremath{i}}\fi}
\newunicodechar{ᵣ}{\ifmmode_{r}\else\textsubscript{\ensuremath{r}}\fi}
\newunicodechar{ₖ}{\ifmmode_{k}\else\textsubscript{\ensuremath{k}}\fi}
\newunicodechar{ᵦ}{\ifmmode_{\beta}\else\textsubscript{\ensuremath{\beta}}\fi}
\newunicodechar{ₓ}{\ifmmode_{x}\else\textsubscript{\ensuremath{x}}\fi}
\newfontfamily\popdisplay{ArchivoBlack-Regular.ttf}[Path=../../../fonts/]
\newfontfamily\poplogo{SpaceGrotesk-Bold.ttf}[Path=../../../fonts/]
\newfontfamily\popsemi{PlusJakartaSans-SemiBold.ttf}[Path=../../../fonts/]
\newfontfamily\popxbold{PlusJakartaSans-ExtraBold.ttf}[Path=../../../fonts/]
\newfontfamily\popxboldls{PlusJakartaSans-ExtraBold.ttf}[Path=../../../fonts/,LetterSpace=3.0]

\setlist[enumerate]{leftmargin=1.7em,itemsep=5pt,topsep=4pt}
\setlist[itemize]{leftmargin=1.7em,itemsep=4pt,topsep=4pt,label={\color{poporange}\textbullet}}
\setlength{\parindent}{0pt}
\setlength{\parskip}{5pt}
\renewcommand{\arraystretch}{1.25}

% pgfplots : style Pop par défaut
\pgfplotsset{
  every axis/.append style={axis line style={popink,line width=1pt},tick style={popink},
    grid style={popline,line width=0.5pt},label style={font=\small},tick label style={font=\footnotesize}},
  popcurve/.style={poporange,line width=1.8pt,no markers,smooth},
}
% Même style disponible dans un \draw TikZ simple (hors pgfplots)
\tikzset{popcurve/.style={poporange,line width=1.8pt}}
\tikzset{popgrid/.style={popline,very thin}}
\colorlet{popcurve}{poporange} % le nom du style est parfois utilisé comme couleur

% ============================================================
% Pill / squiggle
% ============================================================
\newcommand{\poppill}[3]{%
  \tikz[baseline=-0.55ex]\node[draw=popink,line width=1.2pt,rounded corners=2.6pt,%
    fill=#2,inner xsep=6.5pt,inner ysep=3pt]{%
    \popxboldls\fontsize{7}{7}\selectfont\color{#3}\MakeUppercase{#1}};%
}
\newcommand{\popsquiggle}[1]{%
  \tikz[baseline=-0.4ex]{
    \draw[#1,line width=2.6pt,line cap=round]
      plot [smooth] coordinates {(0,0.09) (0.34,0.01) (0.68,0.21) (1.02,0.01) (1.36,0.10)};
  }%
}
% Mot-clé surligné (marqueur orange sous le texte)
\newcommand{\cle}[1]{{\popxbold\color{popdark}#1}}

% ============================================================
% En-tête / pied
% ============================================================
\pagestyle{fancy}
\fancyhf{}
\renewcommand{\headrulewidth}{0pt}
\renewcommand{\footrulewidth}{0pt}
\lhead{\raisebox{-0.15ex}{\poplogo\fontsize{13}{13}\selectfont\color{popink}OnBuch\color{poporange}+}}
\rhead{\poppill{\DOCMATIERE\ \textbullet\ \DOCNIVEAU\ \textbullet\ Leçon \DOCLECON}{white}{popink}}
\lfoot{\popxbold\fontsize{7.6}{7.6}\selectfont\color{popmuted}\MakeUppercase{Cours signé OnBuch+}\ \textbullet\ {\popsemi\DOCTITRE}}
\rfoot{\tikz[baseline=-0.6ex]\node[rounded corners=8.5pt,fill=popink,minimum height=17pt,minimum width=17pt,inner xsep=5pt,inner sep=0]{\popxbold\fontsize{8}{8}\selectfont\color{popcream}\thepage\,/\,\pageref*{LastPage}};}

% ============================================================
% Titres
% ============================================================
\newcommand{\chaptitre}[2]{%
  \begin{tcolorbox}[enhanced,colback=poporange,colframe=popink,boxrule=2.6pt,arc=11pt,
      drop shadow=popink,left=15pt,right=15pt,top=12pt,bottom=11pt,before skip=3pt,after skip=18pt]
    \poppill{#2}{popink}{popcream}\par\vspace{9pt}
    {\raggedright\hyphenpenalty=10000\exhyphenpenalty=10000\popdisplay\fontsize{22}{24}\selectfont\color{popcream}\MakeUppercase{#1}\par}
  \end{tcolorbox}
}
% Section numérotée : pastille orange avec le numéro + titre Archivo
\newcounter{popsec}
\newcommand{\coursec}[1]{%
  \stepcounter{popsec}\setcounter{popsub}{0}%
  \par\vspace{16pt}\noindent
  \begin{minipage}{\linewidth}
  \tikz[baseline=-0.7ex]\node[draw=popink,line width=1.6pt,fill=poporange,rounded corners=4pt,minimum size=22pt,inner sep=0]{\popdisplay\fontsize{12}{12}\selectfont\color{popcream}\thepopsec};%
  \hspace{8pt}{\popdisplay\fontsize{15}{17}\selectfont\color{popink}\MakeUppercase{#1}}\par
  \vspace{4pt}\noindent{\color{poporange}\rule{36pt}{3.6pt}}
  \end{minipage}\par\nopagebreak\vspace{8pt}
}
\newcounter{popsub}
\newcommand{\courssub}[1]{%
  \stepcounter{popsub}%
  \par\vspace{9pt}\noindent{\popxbold\fontsize{11.5}{13}\selectfont\color{popink}{\color{poporange}\thepopsec.\thepopsub}\hspace{6pt}#1}\par\nopagebreak\vspace{4pt}
}

% ============================================================
% Boîtes pédagogiques — même recette : fond blanc, bordure épaisse colorée,
% ombre dure encre, badge plein. Titre optionnel [#1].
% ============================================================
\tcbset{popbase/.style={breakable,enhanced,colback=white,boxrule=2.3pt,arc=9pt,
  shadow={3pt}{-3pt}{0pt}{popink},left=14pt,right=14pt,top=11pt,bottom=11pt,before skip=11pt,after skip=15pt,
  fontupper=\popsemi\fontsize{10.6}{14.5}\selectfont\color{popink},
  fontlower=\popsemi\fontsize{10.6}{14.5}\selectfont\color{popink}}}
% Coche dessinée (le glyphe ✓ n'existe pas dans les polices du document)
\newcommand{\popcheck}[1]{\tikz[baseline=-0.5ex]\draw[#1,line width=1.6pt,line cap=round,line join=round](-0.13,0.0)--(-0.04,-0.09)--(0.13,0.1);}
\providecommand{\coche}{\popcheck{popgreen}}
\providecommand{\resultat}[1]{\par\smallskip\noindent\fcolorbox{popgreen}{popgreenL}{\parbox{\dimexpr\linewidth-2\fboxsep-2\fboxrule\relax}{\textbf{Résultat.} #1}}\par\smallskip}
\newcommand{\popboxhead}[3]{\poppill{#1}{#2}{white}\ifx\relax#3\relax\else\hspace{6pt}{\popxbold\fontsize{10.6}{12}\selectfont\color{popink}#3}\fi\par\vspace{7pt}}

\newtcolorbox{aretenir}[1][]{popbase,colframe=popgreen,before upper={\popboxhead{À retenir}{popgreen}{#1}}}
% Texte en espagnol (document de lecture, dialogue, poème, extrait) : cadre
% d'étude. Un titre optionnel (source, auteur) s'affiche dans l'en-tête.
\newtcolorbox{textebox}[1][]{popbase,colframe=popblue,before upper={\popboxhead{Texto}{popblue}{#1}\begin{otherlanguage}{spanish}},after upper={\end{otherlanguage}}}
% Marques ✓ / ✗ (forme correcte / fautive) souvent utilisées par les modèles
\providecommand{\cross}{\textcolor{poppink}{\ensuremath{\times}}}
\providecommand{\xmark}{\cross}\providecommand{\crossmark}{\cross}\providecommand{\cmark}{\ensuremath{\checkmark}}
% Ligne de réponse / justification à compléter (inventée par les modèles)
\providecommand{\justification}[1]{\par\noindent\textit{Justification~:} \dotfill\par}
% \textipa (package tipa non chargé) : la police ne couvre pas l'API -> simple passage
\providecommand{\textipa}[1]{#1}
% Soulignement (ulem non chargé) : \uline, \ul
\providecommand{\uline}[1]{\underline{#1}}\providecommand{\ul}[1]{\underline{#1}}
% Mot ou phrase en espagnol à l'intérieur d'un paragraphe français
\newcommand{\esp}[1]{\ifmmode\text{\textit{#1}}\else\begingroup\selectlanguage{spanish}\itshape #1\endgroup\fi}
\newtcolorbox{exemplebox}[1][]{popbase,colframe=poporange,before upper={\popboxhead{Exemple}{poporange}{#1}}}
\newtcolorbox{definition}[1][]{popbase,colframe=popblue,before upper={\popboxhead{Définition}{popblue}{#1}}}
\newtcolorbox{propriete}[1][]{popbase,colframe=poppurple,before upper={\popboxhead{Propriété}{poppurple}{#1}}}
\newtcolorbox{methode}[1][]{popbase,colframe=popgold,before upper={\popboxhead{Méthode}{popgold}{#1}}}
\newtcolorbox{attention}[1][]{popbase,colframe=poppink,before upper={\popboxhead{Attention}{poppink}{#1}}}
\newtcolorbox{experience}[1][]{popbase,colframe=popblue,colback=popblueL,before upper={\popboxhead{Expérience}{popblue}{#1}}}
% « Le savais-tu ? » : carte violette pleine (contexte, culture, Cameroun)
\newtcolorbox{savaistu}[1][]{popbase,colback=poppurple,colframe=popink,
  fontupper=\popsemi\fontsize{10.4}{14.2}\selectfont\color{white},
  before upper={\poppill{Le savais-tu\,?}{white}{poppurple}\ifx\relax#1\relax\else\hspace{6pt}{\popxbold\color{white}#1}\fi\par\vspace{7pt}}}
% Exercice résolu : énoncé en haut, \tcblower puis la solution sur fond crème
\newcounter{popexo}\newcounter{popexr}
\newtcolorbox{exoresolu}[1][]{popbase,colframe=poporange,colbacklower=popcream,
  segmentation style={popink,line width=1.2pt,dashed},
  before upper={\stepcounter{popexr}\popboxhead{Exercice résolu \thepopexr}{poporange}{#1}},
  before lower={\poppill{Solution}{popink}{popcream}\par\vspace{6pt}}}
% Exercice (non résolu en place) — la correction est dans la partie Corrigés
\newtcolorbox{exercice}[1][]{popbase,colframe=popink,
  before upper={\stepcounter{popexo}\popboxhead{Exercice \thepopexo}{popink}{#1}}}
% Corrigé
\newtcolorbox{corrige}[1][]{popbase,colframe=popgreen,colback=popgreenL,
  before upper={\popboxhead{Corrigé}{popgreen}{#1}}}
% Objectifs / prérequis / bilan
\newtcolorbox{objectifs}{popbase,colframe=popink,colback=poporangeL,
  before upper={\popboxhead{Objectifs de la leçon}{popink}{}}}
\newtcolorbox{prerequis}{popbase,colframe=popink,colback=popblueL,
  before upper={\popboxhead{Prérequis}{popblue}{}}}
\newtcolorbox{bilan}{popbase,colframe=popink,colback=white,boxrule=2.8pt,
  before upper={\popboxhead{Fiche bilan}{poporange}{L'essentiel à connaître pour l'examen}}}
% Figure encadrée avec légende
\newtcolorbox{popfigure}[1][]{enhanced,breakable=false,colback=white,colframe=popline,boxrule=1.4pt,arc=8pt,
  left=10pt,right=10pt,top=10pt,bottom=8pt,before skip=10pt,after skip=12pt,halign=center,
  lower separated=false,halign lower=center,
  fontlower=\popsemi\fontsize{9}{11.5}\selectfont\color{popmuted},
  #1}
\newcommand{\legende}[1]{\par\vspace{6pt}{\popsemi\fontsize{9}{11.5}\selectfont\color{popmuted}#1}}

% ============================================================
% Signature OnBuch+
% ============================================================
\newcommand{\poplogoinline}[1]{{\poplogo\fontsize{#1}{#1}\selectfont\color{popink}OnBuch\color{poporange}+}}
\newcommand{\signatureonbuch}{%
  \begin{tcolorbox}[enhanced,colback=popink,colframe=popink,boxrule=2.2pt,arc=10pt,
      drop shadow=poporange,left=14pt,right=14pt,top=10pt,bottom=10pt,before skip=0pt,after skip=0pt]
    \tikz[baseline=-0.7ex]\node[circle,fill=popgreen,draw=popcream,line width=1.3pt,minimum size=22pt,inner sep=0]{\popcheck{white}};%
    \hspace{9pt}\begin{minipage}[c]{0.62\linewidth}
      {\popxbold\fontsize{12}{13}\selectfont\color{popcream}Signé\ \textbullet\ {\poplogo OnBuch\color{poporange}+}}\par\vspace{2pt}
      {\raggedright\popsemi\fontsize{8}{10.5}\selectfont\color{popline}\MakeUppercase{Équipe pédagogique OnBuch+}\\\MakeUppercase{Conforme au programme officiel MINESEC}\par}
    \end{minipage}\hfill
    {\poplogo\fontsize{22}{22}\selectfont\color{popcream}OnBuch\color{poporange}+}
  \end{tcolorbox}%
}

% ============================================================
% Page de garde
% #1 titre · #2 matière · #3 niveau · #4 module · #5 n° leçon · #6 durée ·
% #7 description / objectifs (texte ou itemize)
% ============================================================
\newcommand{\pagegarde}[7]{%
  \thispagestyle{empty}
  \newgeometry{a4paper,margin=1.7cm,top=1.6cm,bottom=1.4cm}%
  \begin{tikzpicture}[remember picture,overlay]
    \fill[poporange!18] ([xshift=-3.2cm,yshift=-3.6cm]current page.north east) circle (4.6cm);
    \fill[poppurple!14] ([xshift=2.4cm,yshift=7.6cm]current page.south west) circle (3.8cm);
  \end{tikzpicture}%
  {\poplogo\fontsize{30}{30}\selectfont\color{popink}OnBuch\color{poporange}+}\hfill
  \raisebox{0.6ex}{\poppill{Cours signé \textbullet\ Premium}{popink}{popcream}}\par
  \vspace{1pt}\popsquiggle{poppurple}\par
  \vspace{8pt}
  \begin{tcolorbox}[enhanced,colback=poporange,colframe=popink,boxrule=2.8pt,arc=13pt,
      drop shadow=popink,left=18pt,right=18pt,top=12pt,bottom=13pt,after skip=10pt]
    \poppill{#2 \textbullet\ #3}{popink}{popcream}\hspace{5pt}\poppill{#4}{white}{popink}\par\vspace{8pt}
    {\popdisplay\fontsize{14}{14}\selectfont\color{popink}LEÇON #5}\par\vspace{4pt}
    {\raggedright\hyphenpenalty=10000\exhyphenpenalty=10000\popdisplay\fontsize{25}{27}\selectfont\color{popcream}\MakeUppercase{#1}\par}
    \vspace{8pt}
    {\popxbold\fontsize{9.5}{9.5}\selectfont\color{popink}\MakeUppercase{Durée conseillée : #6}}
  \end{tcolorbox}
  \begin{tcolorbox}[enhanced,colback=white,colframe=popink,boxrule=2.2pt,arc=10pt,
      drop shadow=popink,left=16pt,right=16pt,top=10pt,bottom=10pt,after skip=8pt]
    \poppill{Dans cette leçon}{poppurple}{white}\par\vspace{5pt}
    \popsemi\fontsize{9.3}{12.6}\selectfont\color{popink}#7
  \end{tcolorbox}
  \noindent\begin{minipage}[t]{0.487\linewidth}
    \begin{tcolorbox}[enhanced,colback=popgreen,colframe=popink,boxrule=2.2pt,arc=10pt,
        drop shadow=popink,left=11pt,right=11pt,top=10pt,bottom=10pt,height=3.1cm,valign=center]
      \tcbox[colback=white,colframe=popink,boxrule=1.3pt,arc=5pt,boxsep=0pt,nobeforeafter,
        left=3pt,right=3pt,top=3pt,bottom=3pt,valign=center]{\qrcode[height=1.9cm]{https://play.google.com/store/apps/details?id=com.luvvix.onbuch&hl=fr}}%
      \hspace{8pt}\begin{minipage}[c]{0.5\linewidth}
        {\popdisplay\fontsize{11}{12.5}\selectfont\color{white}\MakeUppercase{Télécharge OnBuch}}\par\vspace{4pt}
        {\raggedright\popsemi\fontsize{8.4}{11}\selectfont\color{white}Gratuit \textbullet\ Google Play\\Tuteur IA, annales, concours, résultats\par}
      \end{minipage}
    \end{tcolorbox}
  \end{minipage}\hfill
  \begin{minipage}[t]{0.487\linewidth}
    \begin{tcolorbox}[enhanced,colback=poppurple,colframe=popink,boxrule=2.2pt,arc=10pt,
        drop shadow=popink,left=11pt,right=11pt,top=10pt,bottom=10pt,height=3.1cm,valign=center]
      \tikz[baseline=-0.7ex]\node[circle,fill=white,draw=popink,line width=1.3pt,minimum size=1.6cm,inner sep=0]{\popdisplay\fontsize{24}{24}\selectfont\color{poppurple}+};%
      \hspace{8pt}\begin{minipage}[c]{0.6\linewidth}
        {\popdisplay\fontsize{11}{12.5}\selectfont\color{white}\MakeUppercase{Passe à OnBuch+}}\par\vspace{4pt}
        {\raggedright\popsemi\fontsize{8.4}{11}\selectfont\color{white}Tous les cours signés, Tuteur IA illimité, fiches PDF — onglet OnBuch+ de l'app\par}
      \end{minipage}
    \end{tcolorbox}
  \end{minipage}
  \vspace{8pt}
  \popcontenu
  \vfill
  \signatureonbuch
  \restoregeometry
  \newpage
}

% Rangée « Ce cours contient » (page de garde)
\newcommand{\poptile}[3]{% #1 couleur #2 gros texte #3 légende
  \begin{tcolorbox}[enhanced,colback=#1,colframe=popink,boxrule=2pt,arc=9pt,shadow={2.5pt}{-2.5pt}{0pt}{popink},
      left=6pt,right=6pt,top=9pt,bottom=9pt,halign=center,valign=center,height=1.75cm,width=\linewidth,nobeforeafter]
    {\popdisplay\fontsize{12.5}{13}\selectfont\color{white}\MakeUppercase{#2}}\par\vspace{3pt}
    {\centering\popsemi\fontsize{7.8}{9.5}\selectfont\color{white}#3\par}
  \end{tcolorbox}}
\newcommand{\popcontenu}{%
  \par\noindent{\popxboldls\fontsize{7.5}{7.5}\selectfont\color{popmuted}CE COURS SIGNÉ CONTIENT}\par\vspace{7pt}
  \noindent\begin{minipage}[t]{0.235\linewidth}\poptile{popblue}{Cours}{complet, illustré, pas à pas}\end{minipage}\hfill
  \begin{minipage}[t]{0.235\linewidth}\poptile{poporange}{Méthodes}{et exercices résolus}\end{minipage}\hfill
  \begin{minipage}[t]{0.235\linewidth}\poptile{poppink}{Exercices}{type examen + corrigés}\end{minipage}\hfill
  \begin{minipage}[t]{0.235\linewidth}\poptile{popgreen}{Fiche bilan}{l'essentiel à retenir}\end{minipage}\par}

% Sommaire
\newtcolorbox{sommaire}{enhanced,colback=white,colframe=popink,boxrule=2.2pt,arc=10pt,
  drop shadow=popink,left=18pt,right=18pt,top=13pt,bottom=13pt,before skip=6pt,after skip=20pt,
  before upper={\poppill{Sommaire}{poppurple}{white}\par\vspace{9pt}\popsemi\fontsize{10.6}{14.5}\selectfont\color{popink}}}

% Signature de fin de cours — poussée tout en bas de la dernière page
\newcommand{\findecours}{%
  \par\vfill\nopagebreak
  \begin{center}\popsquiggle{poporange}\end{center}\vspace{4pt}
  \signatureonbuch
}
\AtBeginDocument{\def\llbracket{\mathopen{[\mkern-3.5mu[}}\def\rrbracket{\mathclose{]\mkern-3.5mu]}}} % intervalles d'entiers (glyphe ⟦ absent de la police)
% Glyphes absents de Fira Math : redéfinis à partir de glyphes disponibles
\AtBeginDocument{%
  \def\setminus{\mathbin{\backslash}}\let\smallsetminus\setminus
  \def\vdots{\mathord{\vbox{\baselineskip3.2pt\lineskiplimit0pt\kern4pt\hbox{.}\hbox{.}\hbox{.}}}}%
  \def\ddots{\mathinner{\mkern1mu\raise6.5pt\hbox{.}\mkern2mu\raise3.6pt\hbox{.}\mkern2mu\raise0.7pt\hbox{.}\mkern1mu}}%
  \def\triangle{\mathord{\scalebox{0.9}{$\symup{\Delta}$}}}%
  \def\nabla{\mathord{\raisebox{\depth}{\scalebox{1}[-1]{$\symup{\Delta}$}}}}%
  \def\checkmark{\popcheck{popdark}}%
  \def\star{\mathbin{\ast}}\def\bigstar{\mathord{\ast}}%
  \def\diamond{\mathbin{\scalebox{0.6}{$\Diamond$}}}%
  \def\bigcup{\mathop{\mathchoice{\raisebox{-0.3em}{\scalebox{1.7}{$\cup$}}}{\scalebox{1.25}{$\cup$}}{\cup}{\cup}}\displaylimits}%
  \def\bigcap{\mathop{\mathchoice{\raisebox{-0.3em}{\scalebox{1.7}{$\cap$}}}{\scalebox{1.25}{$\cap$}}{\cap}{\cap}}\displaylimits}%
  \def\lesseqqgtr{\mathrel{\lessgtr}}\def\bigoplus{\mathop{\oplus}\displaylimits}%
}
\providecommand{\ssi}{si et seulement si} % abréviation fréquente
\AtBeginDocument{\providecommand{\wideparen}{\overparen}} % arc de cercle

%% FIGFIX-BEGIN v3 — correctifs globaux des figures (docs/onbuch-generation/tools/figfix)
\usepackage{adjustbox}\usepackage{etoolbox}
% toute figure TikZ/pgfplots plus large que la ligne est réduite à la largeur disponible
\BeforeBeginEnvironment{tikzpicture}{\begin{adjustbox}{max width=\linewidth}}
\AfterEndEnvironment{tikzpicture}{\end{adjustbox}}
% légendes pgfplots lisibles par-dessus les courbes
\pgfplotsset{every axis/.append style={legend style={fill=white,fill opacity=0.92,text opacity=1,draw=black!15,inner sep=3pt}}}
% étiquettes d'arêtes (syntaxe edge["..."]) avec fond blanc
\tikzset{every edge quotes/.append style={fill=white,inner sep=1.2pt}}
% bulles de cartes mentales : texte à la ligne dans la bulle, bulle un peu plus grande
\tikzset{every concept/.append style={text width=2.0cm,align=center,minimum size=2.3cm,inner sep=2pt}}
%% FIGFIX-END
\begin{document}

\pagegarde{Lectura y comentario : Internet y las NTIC ; traduction : expression de la préférence et de l'habitude}{Espagnol}{Tle A}{Módulo 5 : Médias, communication et culture de la paix}{E28}{2 h}{Leçon mixte du Módulo 5 : lecture et commentaire d'un texte sur Internet et les NTIC, enrichissement du lexique des technologies, puis étude grammaticale et traductive de l'expression de la préférence (preferir, gustar más) et de l'habitude (soler, acostumbrar, imparfait).
\vspace{4pt}
\begin{multicols}{2}\small
\begin{itemize}
\item À la fin de cette leçon, je sais lire et comprendre un texte argumentatif sur Internet et les NTIC
\item je sais utiliser le lexique des technologies : red social, contraseña, descargar, navegar, teléfono inteligente, brecha digital
\item je sais rédiger un commentaire structuré d'un texte sur les NTIC (introduction, développement, conclusion)
\item je sais exprimer une préférence avec preferir, me gusta más, prefiero... a...
\item je sais exprimer une habitude présente avec soler + infinitif, acostumbrar a, tener la costumbre de
\item je sais exprimer une habitude passée avec l'imparfait de l'indicatif
\end{itemize}\end{multicols}}

\begin{sommaire}
\begin{multicols}{2}
\begin{enumerate}
\item Texte support : Internet y las NTIC
\item Compréhension du texte
\item Lexique thématique des NTIC
\item Grammaire 1 : l'expression de la préférence
\item Grammaire 2 : l'expression de l'habitude
\item Méthode du commentaire et production guidée
\item Activité d'intégration
\item Méthodes et exercices résolus
\item Exercices
\item Corrigés des exercices
\item Fiche bilan
\end{enumerate}
\end{multicols}
\end{sommaire}

\begin{objectifs}
\begin{itemize}
\item À la fin de cette leçon, je sais lire et comprendre un texte argumentatif sur Internet et les NTIC
\item je sais utiliser le lexique des technologies : red social, contraseña, descargar, navegar, teléfono inteligente, brecha digital
\item je sais rédiger un commentaire structuré d'un texte sur les NTIC (introduction, développement, conclusion)
\item je sais exprimer une préférence avec preferir, me gusta más, prefiero... a...
\item je sais exprimer une habitude présente avec soler + infinitif, acostumbrar a, tener la costumbre de
\item je sais exprimer une habitude passée avec l'imparfait de l'indicatif
\item je sais conjuguer preferir et soler (verbes à diphtongaison e→ie, o→ue)
\item je sais traduire du français vers l'espagnol et inversement des phrases contenant préférence et habitude
\item je sais éviter les pièges de traduction propres aux francophones (faux amis, constructions)
\end{itemize}
\end{objectifs}

\begin{prerequis}
\begin{itemize}
\item Présent de l'indicatif des verbes réguliers et des verbes à diphtongaison (Première)
\item Imparfait de l'indicatif : formation et emplois de base (Première)
\item Verbe gustar et verbes de sentiment construits avec l'indirect (Première)
\item Vocabulaire général des médias et de la communication (début du Módulo 5)
\item Techniques de compréhension de texte : repérage des connecteurs et de la thèse (Première)
\end{itemize}
\end{prerequis}

\newpage

\coursec{Texte support : Internet y las NTIC}

À Yaoundé comme à Douala, les élèves de Terminale consultent WhatsApp, TikTok et Google avant même d'ouvrir leurs cahiers. Au Cameroun, plus de 10 millions d'internautes naviguent chaque jour, mais beaucoup de zones rurales restent déconnectées : c'est la fameuse \cle{brecha digital}. En Espagne et en Amérique latine, le débat est le même : Internet est-il un outil de liberté ou une nouvelle dépendance ? Aujourd'hui, nous allons lire un texte sur les NTIC, enrichir notre vocabulaire et apprendre à dire en espagnol ce que nous préférons et ce que nous avons l'habitude de faire en ligne.

\courssub{Lecture globale : « Internet: ¿puente o brecha? »}

\begin{textebox}[Internet: ¿puente o brecha?]
Internet: ¿puente o brecha? — Hoy en día, millones de jóvenes navegan por Internet cada día. Gracias al teléfono inteligente, se puede descargar información, chatear en las redes sociales y estudiar en línea. Sin embargo, no todos tienen acceso: la brecha digital separa a las ciudades de los pueblos. Además, muchos usuarios no protegen su contraseña y sufren robos de datos. Por eso, los expertos acostumbran a aconsejar prudencia. Internet es un puente hacia el saber, pero solamente si se usa con responsabilidad.
\end{textebox}

\begin{textebox}[Glossaire]
\textbf{navegar} : naviguer (sur le web) \\
\textbf{descargar} : télécharger \\
\textbf{chatear} : discuter en ligne (anglicisme \emph{chat}) \\
\textbf{en línea} : en ligne \\
\textbf{la brecha digital} : la fracture numérique \\
\textbf{la contraseña} : le mot de passe \\
\textbf{el robo de datos} : le vol de données \\
\textbf{el saber} : le savoir, la connaissance \\
\textbf{aconsejar} : conseiller \\
\textbf{la prudencia} : la prudence \\
\textbf{acostumbrar a} : avoir l'habitude de (suivi de l'infinitif)
\end{textebox}

\noindent
\textbf{1. Thème et thèse.} Le thème est \cle{Internet y las nuevas tecnologías de la información y la comunicación (NTIC)}. La thèse, nuancée, se trouve dans la dernière phrase : \esp{Internet es un puente hacia el saber, pero solamente si se usa con responsabilidad.} (« Internet est un pont vers le savoir, mais seulement s'il est utilisé avec responsabilité. ») L'auteur défend l'idée que l'outil est neutre ; tout dépend de l'usage et de l'accès.

\noindent
\textbf{2. Structure argumentative.}
\begin{itemize}
    \item \textbf{Arguments « pour » (avantages) :} accès à l'information (\esp{descargar información}), communication (\esp{chatear en las redes sociales}), éducation (\esp{estudiar en línea}), pont vers le savoir (\esp{puente hacia el saber}).
    \item \textbf{Arguments « contre » (dangers/limites) :} inégalité d'accès (\esp{la brecha digital separa a las ciudades de los pueblos}), insécurité (\esp{no protegen su contraseña, sufren robos de datos}), nécessité de prudence.
\end{itemize}

\courssub{Analyse détaillée : connecteurs logiques et champs lexicaux}

\begin{definition}[Connecteurs logiques du texte]
Les connecteurs (ou marqueurs de relation) organisent la progression du raisonnement. Les repérer permet de suivre la pensée de l'auteur.
\end{definition}

\begin{center}
\begin{tabularx}{\linewidth}{|X|X|X|}
\hline
\rowcolor{popblueL} \textbf{Connecteur} & \textbf{Fonction} & \textbf{Exemple du texte / Traduction} \\
\hline
\textbf{Sin embargo} & Opposition / Concession & \esp{Sin embargo, no todos tienen acceso} (Cependant, tout le monde n'a pas accès) \\
\hline
\textbf{Además} & Addition / Renforcement & \esp{Además, muchos usuarios no protegen su contraseña} (De plus, beaucoup d'utilisateurs ne protègent pas leur mot de passe) \\
\hline
\textbf{Por eso} & Conséquence / Conclusion & \esp{Por eso, los expertos acostumbran a aconsejar prudencia} (C'est pourquoi les experts ont l'habitude de conseiller la prudence) \\
\hline
\textbf{Pero} & Opposition (coordination) & \esp{es un puente..., pero solamente si...} (c'est un pont..., mais seulement si...) \\
\hline
\textbf{Gracias a} & Cause (valeur positive) & \esp{Gracias al teléfono inteligente...} (Grâce au smartphone...) \\
\hline
\end{tabularx}
\end{center}

\begin{propriete}[Champs lexicaux : repérage et classification]
Un champ lexical regroupe les mots liés à une même idée. En les classant, on voit la structure thématique du texte.
\end{propriete}

\begin{center}
\begin{tabularx}{\linewidth}{|X|X|X|X|}
\hline
\rowcolor{popblueL} \textbf{Herramientas (Outils)} & \textbf{Acciones (Actions)} & \textbf{Ventajas (Avantages)} & \textbf{Peligros / Problemas (Dangers)} \\
\hline
el teléfono inteligente & navegar & la información & la brecha digital \\
las redes sociales & descargar & el saber & el robo de datos \\
Internet & chatear & estudiar en línea & no protegen la contraseña \\
la contraseña & estudiar & un puente & sufren \\
& proteger & & \\
& aconsejar & & \\
& usar & & \\
\hline
\end{tabularx}
\end{center}

\begin{aretenir}[Observation grammaticale anticipée]
Remarquez dans le texte deux structures que nous étudierons en grammaire (sections 4 et 5) :
\begin{enumerate}
    \item \textbf{Expression de l'habitude :} \esp{los expertos \cle{acostumbran a} aconsejar prudencia} (présent de \esp{acostumbrar a} + infinitif).
    \item \textbf{Impersonnel + infinitif :} \esp{se puede descargar} (on peut télécharger), \esp{se usa} (on utilise / il s'utilise). La voix passive réflexe (\esp{se} + verbe) est très fréquente en espagnol journalistique.
\end{enumerate}
\end{aretenir}

\courssub{Contexte camerounais : la fracture numérique au quotidien}

Au Cameroun, la \cle{brecha digital} est une réalité vécue quotidiennement. Si les élèves des lycées de Yaoundé,

\coursec{Texte support : Internet y las NTIC}

\courssub{Lecture analytique et découverte du lexique}

\begin{textebox}[Texte original : Internet y las NTIC — una herramienta de doble filo]
En Yaoundé, como en Madrid o Buenos Aires, el \cle{teléfono inteligente} se ha vuelto una extensión de la mano de los jóvenes. Cada mañana, antes incluso de desayunar, miles de estudiantes \cle{navegan} por la \cle{red} para consultar sus \cle{redes sociales}, responder mensajes en WhatsApp o \cle{descargar} apuntes de clase desde la nube. Las \cle{NTIC} (Nuevas Tecnologías de la Información y la Comunicación) han roto las fronteras del aula: un alumno de Nkolbisson puede seguir un curso de la Universidad de Salamanca sin salir de su barrio.

Sin embargo, esta revolución digital tiene su \cle{brecha digital}. En las zonas rurales de la región del Este o del Norte de Camerún, la conexión es inestable o inexistente; muchos hogares no tienen computadora y comparten un solo smartphone familiar. Además, el exceso de pantalla genera riesgos: \cle{ciberacoso}, robo de \cle{contraseña}, adicción a los videojuegos y pérdida de sueño. Los psicólogos advierten que la \cle{identidad digital} de un adolescente se construye hoy entre \emph{likes} y \emph{stories}, a veces lejos de la realidad.

La solución no es prohibir, sino \cle{educar}. Padres y profesores deben enseñar el \cle{uso responsable}: verificar la fuente antes de compartir, proteger la privacidad con contraseñas robustas, limitar el tiempo de conexión y privilegiar el encuentro presencial. Internet es un \cle{mar de conocimiento}; saber nadar en él sin ahogarse es la gran competencia del siglo XXI.
\end{textebox}

\coursec{Compréhension du texte}

\courssub{Transition et rappel}

Après la lecture analytique du texte \emph{« Internet y las NTIC : una herramienta de doble filo »}, nous abordons l'étape fondamentale de la \cle{compréhension écrite} : extraire le sens global, identifier les arguments précis, maîtriser le lexique en contexte et reformuler avec ses propres mots. Au Bac (épreuve écrite de LV2), cette compétence est notée sur 6 à 8 points ; elle exige rigueur, citations exactes (guillemets espagnols \esp{« ... »}) et reformulation personnelle — jamais de copier-coller intégral.

\courssub{Questions de compréhension globale}

\begin{exemplebox}[Questions types — Bac LV2]
\begin{enumerate}
    \item \textbf{Sujet du texte :} De quoi parle le texte en une phrase ? (Sujet = thème + thèse).
    \item \textbf{Thèse de l'auteur :} Quelle est l'opinion défendue ? (Pour / contre / nuancé).
    \item \textbf{Ton du texte :} Informatif ? Alarmiste ? Engagé ? Pédagogique ? Justifiez par des indices lexicaux.
\end{enumerate}
\end{exemplebox}

\begin{exoresolu}[Réponses types — Compréhension globale]
\textbf{Énoncé :} Répondez aux trois questions globales ci-dessus en vous appuyant sur le texte.

\tcblower

\textbf{Solution détaillée :}

\begin{enumerate}
    \item \textbf{Sujet :} Le texte traite de l'impact des NTIC et d'Internet sur la jeunesse (notamment camerounaise), entre opportunités éducatives et risques sociaux, en plaidant pour une éducation numérique responsable.
    \item \textbf{Thèse :} L'auteur adopte une position \textbf{nuancée et pédagogique} : Internet est une \esp{« herramienta de doble filo »} (outil à double tranchant). Il reconnaît les avantages (\esp{« han roto las fronteras del aula »}) mais alerte sur les dangers (\esp{« ciberacoso », « adicción »}) et propose une solution éducative (\esp{« educar », « uso responsable »}).
    \item \textbf{Ton :} \textbf{Pédagogique et engagé.} Indices : verbes d'obligation/conseil (\esp{deben enseñar, hay que verificar, proteger, limitar}), métaphore finale (\esp{« mar de conocimiento ; saber nadar »}), vocabulaire didactique (\esp{competencia, educar, solución}). Ce n'est pas un ton purement informatif ni alarmiste.
\end{enumerate}
\end{exoresolu}

\courssub{Questions de détail : Vrai / Faux justifié par citation}

\begin{methode}[Méthode Bac : Vrai ou Faux justifié]
\begin{enumerate}
    \item Lisez l'affirmation et repérez le passage correspondant (mots-clés).
    \item Décidez : \textbf{VRAI} (correspond au texte) ou \textbf{FAUX} (contraire ou non dit).
    \item \textbf{Citez} le segment exact entre guillemets espagnols \esp{« ... »}.
    \item Expliquez en français \textbf{pourquoi} la citation prouve votre choix (reformulation).
    \item \textbf{Interdit :} recopier tout le paragraphe ; répondre seulement « Vrai/Faux » sans citation.
\end{enumerate}
\end{methode}

\begin{exercice}[Vrai / Faux justifié]
Déterminez si les affirmations sont VRAIES ou FAUSSES. Justifiez par une citation précise et une brève explication en français.
\begin{enumerate}
    \item Les NTIC permettent à un élève de Nkolbisson d'accéder à des cours universitaires étrangers sans voyager.
    \item La connexion internet est stable et rapide dans toutes les régions du Cameroun, y compris les zones rurales.
    \item L'auteur considère que l'interdiction totale d'Internet est la meilleure solution pour protéger les jeunes.
    \item Le vol de mot de passe est cité comme l'un des risques liés à l'excès d'écran.
    \item L'auteur compare Internet à une mer où il faut apprendre à nager pour ne pas se noyer.
\end{enumerate}
\end{exercice}

\begin{corrige}[Exercice : Vrai / Faux justifié]
\begin{enumerate}
    \item \textbf{VRAI.} Citation : \esp{« un alumno de Nkolbisson puede seguir un curso de la Universidad de Salamanca sin salir de su barrio »}. Explication : Exemple concret d'accès à distance à l'enseignement supérieur.
    \item \textbf{FAUX.} Citation : \esp{« En las zonas rurales de la región del Este o del Norte de Camerún, la conexión es inestable o inexistente »}. Explication : Le texte affirme l'instabilité ou l'absence de connexion en zone rurale.
    \item \textbf{FAUX.} Citation : \esp{« La solución no es prohibir, sino educar »}. Explication : L'auteur rejette l'interdiction (\esp{no es prohibir}) et préconise l'éducation (\esp{sino educar}).
    \item \textbf{VRAI.} Citation : \esp{« robo de contraseña »} (dans l'énumération : \esp{ciberacoso, robo de contraseña, adicción...}). Explication : Figure explicitement dans la liste des dangers.
    \item \textbf{VRAI.} Citation : \esp{« Internet es un mar de conocimiento ; saber nadar en él sin ahogarse es la gran competencia del siglo XXI »}. Explication : Métaphore finale comparant Internet à une mer et l'usage responsable à la natation.
\end{enumerate}
\end{corrige}

\courssub{Réformulation : résumer en trois phrases}

\begin{methode}[Méthode : Résumer un texte en 3 phrases (Bac)]
\begin{enumerate}
    \item Identifiez les \textbf{3 macro-idées} (Contexte/Avantages $\rightarrow$ Problème/Risques $\rightarrow$ Solution/Conclusion).
    \item Rédigez \textbf{une phrase par macro-idée} avec \textbf{vos propres mots} (pas de citation).
    \item Utilisez des connecteurs logiques (\esp{en primer lugar, sin embargo, por tanto, en conclusión}).
    \item Vérifiez : limite de mots (60–80), exactitude, correction.
\end{enumerate}
\end{methode}

\begin{exoresolu}[Résumé type en 3 phrases]
\textbf{Énoncé :} Résumez le texte en exactement trois phrases (environ 70 mots) en espagnol.

\tcblower

\textbf{Proposition de résumé (en espagnol, ~65 mots) :}

\esp{En primer lugar, las NTIC han transformado la educación permitiendo el acceso remoto al conocimiento, aunque persiste una brecha digital que margina a las zonas rurales. Sin embargo, el uso excesivo de pantallas conlleva riesgos graves como el ciberacoso, el robo de datos o la adicción. Por tanto, la solución no radica en la prohibición, sino en una educación digital responsable que enseñe a proteger la privacidad y a equilibrar vida virtual y real.}

\textbf{Traduction (vérification) :} « Premièrement, les NTIC ont transformé l'éducation en permettant l'accès à distance au savoir, bien qu'une fracture numérique persiste en marge des zones rurales. Cependant, l'usage excessif des écrans entraîne des risques graves comme le cyberharcèlement, le vol de données ou l'addiction. Par conséquent, la solution ne réside pas dans l'interdiction, mais dans une éducation numérique responsable qui enseigne à protéger la vie privée et à équilibrer vie virtuelle et réelle. »
\end{exoresolu}

\courssub{Organigramme de la structure argumentative}

\begin{popfigure}
\begin{tikzpicture}[
    node distance=1.2cm and 2.5cm,
    main/.style={draw=popblue, thick, rounded corners, fill=popblueL, text width=3.2cm, align=center, font=\small},
    arrow/.style={-Stealth, thick, popdark}
]
\node[main, align=center] (thesis){\textbf{Tesis central}\\\emph{Internet: herramienta de doble filo}};
\node[main, below left=of thesis, align=center] (arg1){\textbf{Argumento 1: Ventajas}\\\emph{Acceso educativo, romper fronteras\\Ej: alumno Nkolbisson / Salamanca}};
\node[main, below right=of thesis, align=center] (arg2){\textbf{Argumento 2: Peligros}\\\emph{Brecha digital, ciberacoso,\\adicción, robo de contraseña}};
\node[main, below=2.5cm of thesis, align=center] (concl){\textbf{Conclusión: Uso responsable}\\\emph{Educar, no prohibir\\Verificar fuentes, proteger privacidad\\Limitar tiempo, encuentro real}};

\draw[arrow] (thesis.south west) -- (arg1.north);
\draw[arrow] (thesis.south east) -- (arg2.north);
\draw[arrow] (arg1.south east) -- (concl.north west);
\draw[arrow] (arg2.south west) -- (concl.north east);
\end{tikzpicture}
\legende{Structure argumentative : Thèse $\rightarrow$ Arguments (Avantages / Dangers) $\rightarrow$ Conclusion (Usage responsable)}
\end{popfigure}

\coursec{Lexique thématique des NTIC}

\courssub{Tableau récapitulatif (vocabulaire du texte + extensions)}

\begin{center}
\begin{tabularx}{\linewidth}{|X|X|X|}\hline
\rowcolor{popblueL} \textbf{Español} & \textbf{Français} & \textbf{Note / Contexte} \\ \hline
\cle{teléfono inteligente} & smartphone / téléphone intelligent & Synonyme : \esp{móvil}, \esp{celular} (AmL) \\ \hline
\cle{navegar} & naviguer (sur le web) & \esp{navegar por la red / por internet} \\ \hline
\cle{red} & réseau (internet) & \esp{la red de redes} = internet \\ \hline
\cle{redes sociales} & réseaux sociaux & \esp{usuario de redes sociales} \\ \hline
\cle{descargar} & télécharger (vers son appareil) & $\neq$ \esp{cargar} (charger une page/batterie) \\ \hline
\cle{subir / colgar} & mettre en ligne / uploader & \esp{subir una foto a Instagram} \\ \hline
\cle{NTIC} & NTIC & Nouvelles Technologies de l'Information et de la Communication \\ \hline
\cle{brecha digital} & fracture numérique & Inégalité d'accès (urbain/rural, riche/pauvre) \\ \hline
\cle{ciberacoso} & cyberharcèlement & \esp{ciberacosador / víctima} \\ \hline
\cle{contraseña} & mot de passe & \esp{contraseña robusta / segura} \\ \hline
\cle{identidad digital} & identité numérique & Traces laissées en ligne \\ \hline
\cle{uso responsable} & usage responsable & Éducation aux médias (\esp{alfabetización mediática}) \\ \hline
\cle{verificar la fuente} & vérifier la source & Lutte contre les \emph{fake news} (\esp{noticias falsas}) \\ \hline
\cle{proteger la privacidad} & protéger la vie privée & \esp{configuración de privacidad} \\ \hline
\cle{limitar el tiempo} & limiter le temps & \esp{control parental}, \esp{tiempo de pantalla} \\ \hline
\end{tabularx}
\end{center}

\courssub{Pièges lexicaux et nuances}

\begin{attention}[Piège fréquent : \esp{descargar} vs \esp{cargar} vs \esp{subir}]
\begin{itemize}
    \item \esp{Descargar} = \textbf{télécharger} (serveur $\rightarrow$ appareil). Ex: \esp{Descargué el PDF del temario.}
    \item \esp{Cargar} = \textbf{charger} (batterie, page web, camion). Ex: \esp{Tengo que cargar el móvil, la batería está al 10\%.} / \esp{La página tarda en cargar.}
    \item \esp{Subir} / \esp{Colgar} = \textbf{mettre en ligne, uploader} (appareil $\rightarrow$ serveur). Ex: \esp{Subí las fotos a la nube.}
    \item \textbf{Erreur de francophone :} dire \esp{*cargar un archivo} pour « télécharger un fichier » ou \esp{*bajar} (calque de l'anglais \emph{download} mais moins standard que \esp{descargar}).
\end{itemize}
\end{attention}

\begin{attention}[Faux ami : \esp{red} vs \esp{réseau} (fr)]
\esp{Red} est féminin (\esp{la red}). Ne pas confondre avec \esp{el red} (n'existe pas). \esp{Red social} = réseau social (singulier) ; \esp{redes sociales} = réseaux sociaux (pluriel courant).
\end{attention}

\begin{exemplebox}[Recherche d'équivalents dans le texte (Entraînement Bac)]
Retrouvez dans le texte le mot/la expression espagnole correspondant à :
\begin{enumerate}
    \item Télécharger (verbe, §1) $\rightarrow$ \esp{descargar}
    \item Mot de passe (nom f., §2) $\rightarrow$ \esp{contraseña}
    \item Fracture numérique (nom f., §2) $\rightarrow$ \esp{brecha digital}
    \item Cyberharcèlement (nom m., §2) $\rightarrow$ \esp{ciberacoso}
    \item Identité numérique (nom f., §2) $\rightarrow$ \esp{identidad digital}
    \item Usage responsable (nom m., §3) $\rightarrow$ \esp{uso responsable}
    \item Vérifier la source (locution verbale, §3) $\rightarrow$ \esp{verificar la fuente}
\end{enumerate}
\end{exemplebox}

\coursec{Grammaire 1 : Expression de la préférence}

\courssub{Les structures principales}

\begin{propriete}[Verbe \esp{preferir} (préférer)]
\esp{Preferir} est un verbe \textbf{irrégulier} (diphthongaison \textbf{e $\rightarrow$ ie} à toutes les personnes sauf \emph{nosotros/as} et \emph{vosotros/as} ; et \textbf{e $\rightarrow$ i} à la 3\textsuperscript{e} personne du singulier et pluriel au passé simple / subjonctif présent).
\begin{center}
\begin{tabular}{|l|l|l|}\hline
\rowcolor{popblueL} \textbf{Personne} & \textbf{Présent indicatif} & \textbf{Imparfait indicatif} \\ \hline
\emph{yo} & prefier\textbf{o} & prefer\textbf{ía} \\ \hline
\emph{tú} & prefier\textbf{es} & prefer\textbf{ías} \\ \hline
\emph{él/ella/usted} & prefier\textbf{e} & prefer\textbf{ía} \\ \hline
\emph{nosotros/as} & prefer\textbf{imos} & prefer\textbf{íamos} \\ \hline
\emph{vosotros/as} & prefer\textbf{ís} & prefer\textbf{íais} \\ \hline
\emph{ellos/ellas/ustedes} & prefier\textbf{en} & prefer\textbf{ían} \\ \hline
\end{tabular}
\end{center}
\textbf{Constructions :}
\begin{itemize}
    \item \esp{Preferir + infinitif} : \esp{Prefiero estudiar en la biblioteca.}
    \item \esp{Preferir + nom} : \esp{Prefiero el té al café.} (Notez \esp{al} = \esp{a + el}).
    \item \esp{Preferir A a B} (structure comparative) : \esp{Prefiero la playa \textbf{a} la montaña.}
\end{itemize}
\end{propriete}

\begin{propriete}[Verbe \esp{gustar} + \emph{más} / \emph{mejor}]
\esp{Gustar} fonctionne à l'envers : la chose aimée est le \textbf{sujet} (accord du verbe), la personne est \textbf{COI} (\emph{me, te, le, nos, os, les}).
\begin{itemize}
    \item \esp{Me gusta \textbf{más} el fútbol \textbf{que} el baloncesto.} (Comparatif de supériorité).
    \item \esp{Me gusta \textbf{mejor} la versión original.} (Synonyme de \emph{más}, plus familier).
    \item \esp{Antes \textbf{que} + subjonctif} : \esp{Prefiero ir al cine \textbf{antes que} quedarme en casa.} (\esp{antes que} exige le subjonctif : \esp{quedarme} $\rightarrow$ \esp{quede} ? Non, \esp{antes que} + subjonctif $\rightarrow$ \esp{antes que me quede}). \textbf{Attention :} \esp{Prefiero ir al cine antes que quedarme en casa} (infinitif si même sujet) mais \esp{Prefiero que vayas tú antes que vaya yo} (subjonctif si sujet différent). Avec \esp{preferir}, on utilise souvent l'infinitif pour le même sujet : \esp{Prefiero leer antes que ver la tele.}
\end{itemize}
\end{propriete}

\begin{attention}[Pièges francophones : Préférence]
\begin{itemize}
    \item \textbf{*Prefiero que} + indicatif $\rightarrow$ \textbf{FAUX}. \esp{Preferir} + \emph{que} + \textbf{subjonctif} : \esp{Prefiero que vengas mañana.}
    \item \textbf{*Me prefiero} $\rightarrow$ \textbf{FAUX}. \esp{Preferir} n'est pas pronominal. On dit \esp{Yo prefiero...} ou \esp{A mí me gusta más...}
    \item \textbf{*Mejor que} pour comparer deux noms avec \esp{gustar} $\rightarrow$ \esp{Me gusta más A que B}. \esp{Mejor} est adjectif/adverbe, pas conjonction de comparaison.
    \item \textbf{Accord \esp{gustar}} : \esp{Me gustan \textbf{más} los perros \textbf{que} los gatos.} (Sujet pluriel $\rightarrow$ verbe pluriel).
\end{itemize}
\end{attention}

\begin{exemplebox}[Exemples traduits : Préférence]
\begin{itemize}
    \item \esp{Prefiero navegar por páginas educativas en lugar de perder el tiempo en redes sociales.} $\rightarrow$ « Je préfère naviguer sur des sites éducatifs au lieu de perdre du temps sur les réseaux sociaux. »
    \item \esp{Me gusta más leer un libro físico que estar frente a la pantalla.} $\rightarrow$ « J'aime mieux lire un livre physique qu'être devant l'écran. »
    \item \esp{Prefiero que mis hijos usen Internet con supervisión.} $\rightarrow$ « Je préfère que mes enfants utilisent Internet sous surveillance. » (Subjonctif \esp{usen}).
    \item \esp{Antes que compartir una noticia, verifico la fuente.} $\rightarrow$ « Avant de partager une nouvelle, je vérifie la source. » (Même sujet $\rightarrow$ infinitif après \esp{antes que} ? Non, \esp{antes de} + infinitif. \esp{Antes que} + subjonctif si sujet différent. Ici \esp{Antes de compartir} est correct. \esp{Antes que comparta} si sujet différent).
\end{itemize}
\end{exemplebox}

\begin{exoresolu}[Exercice d'application : Préférence]
\textbf{Énoncé :} Traduisez en espagnol.
\begin{enumerate}
    \item Je préfère utiliser un ordinateur plutôt qu'un téléphone.
    \item Nous aimons mieux lire sur papier que sur écran.
    \item Ils préfèrent que nous limitions le temps de connexion.
    \item Elle préfère les réseaux sociaux à la télévision.
\end{enumerate}

\tcblower

\textbf{Solution détaillée :}
\begin{enumerate}
    \item \esp{Prefiero usar un ordenador \textbf{en lugar de} un teléfono.} (ou \esp{antes que usar un teléfono} / \esp{a un teléfono}).
    \item \esp{Nos gusta \textbf{más} leer en papel \textbf{que} en pantalla.} (Sujet infinitif singulier $\rightarrow$ \esp{gusta}).
    \item \esp{Prefieren \textbf{que} limitemos el tiempo de conexión.} (Subjonctif présent \esp{limitemos} car sujet différent : \emph{ellos} vs \emph{nosotros}).
    \item \esp{Ella prefiere las redes sociales \textbf{a} la televisión.} (Structure \esp{preferir A a B}).
\end{enumerate}
\end{exoresolu}

\begin{exercice}[Entraînement : Préférence]
Complétez avec la forme correcte (\esp{preferir}, \esp{gustar más}, \esp{antes que}, subjonctif si nécessaire).
\begin{enumerate}
    \item Yo \_\_\_\_\_\_ (preferir) estudiar en silencio \_\_\_\_\_\_ con música.
    \item \_\_\_\_\_\_ (gustar) \_\_\_\_\_\_ (más) las clases presenciales \_\_\_\_\_\_ las virtuales.
    \item El profesor prefiere \_\_\_\_\_\_ (que / nosotros / entregar) los deberes a tiempo.
    \item \_\_\_\_\_\_ (antes que) \_\_\_\_\_\_ (subir) la foto, \_\_\_\_\_\_ (verificar) la privacidad.
\end{enumerate}
\end{exercice}

\begin{corrige}[Exercice : Préférence]
\begin{enumerate}
    \item \esp{prefiero ... a} (\esp{Prefiero estudiar en silencio a con música}).
    \item \esp{Nos gustan más ... que} (Sujet \emph{las clases} pluriel).
    \item \esp{que entreguemos} (Subjonctif présent 1ère personne pluriel : \esp{entreguemos}).
    \item \esp{Antes de subir ... verifica} (Même sujet implicite $\rightarrow$ \esp{Antes de} + infinitif ; impératif ou présent indicatif pour le conseil). Ou : \esp{Antes de que la subas, verifica...} (Sujet différent).
\end{enumerate}
\end{corrige}

\coursec{Grammaire 2 : Expression de l'habitude}

\courssub{Les quatre piliers de l'habitude}

\begin{propriete}[1. \esp{Soler} + infinitif (Habitude présente / passée)]
Verbe \textbf{irrégulier} (diphthongaison \textbf{o $\rightarrow$ ue} au présent indicatif sauf \emph{nosotros/vosotros}). N'existe qu'aux temps où l'aspect habituel a du sens (Présent, Imparfait, Conditionnel). \textbf{Pas de gérondif, pas de participe présent, pas d'impératif.}
\begin{center}
\begin{tabular}{|l|l|l|}\hline
\rowcolor{popblueL} \textbf{Personne} & \textbf{Présent (Habitude actuelle)} & \textbf{Imparfait (Habitude passée)} \\ \hline
\emph{yo} & suel\textbf{o} & sol\textbf{ía} \\ \hline
\emph{tú} & suel\textbf{es} & sol\textbf{ías} \\ \hline
\emph{él/ella/usted} & suel\textbf{e} & sol\textbf{ía} \\ \hline
\emph{nosotros/as} & sol\textbf{emos} & sol\textbf{íamos} \\ \hline
\emph{vosotros/as} & sol\textbf{eis} & sol\textbf{íais} \\ \hline
    \emph{ellos/ellas/ustedes} & suel\textbf{en} & sol\textbf{ían} \\ \hline
\end{tabular}
\end{center}
\textbf{Exemples :}
\esp{Suelo navegar por la red cada mañana.} (Habitude actuelle).
\esp{Antes solía descargar muchos juegos, ya no.} (Habitude passée révolue).
\esp{¿Sueles verificar las fuentes?} (Interrogation sur l'habitude).
\end{propriete}

\begin{propriete}[2. \esp{Acostumbrar a} + infinitif]
Verbe \textbf{régulier} (souvent pronominal \esp{acostumbrarse a} = « s'habituer à », mais \esp{acostumbrar a} = « avoir l'habitude de »).
\begin{itemize}
    \item \esp{Acostumbro a revisar el correo al despertar.} (Habitude présente).
    \item \esp{Me acostumbré a la nueva interfaz.} (Processus d'adaptation : \emph{s'habituer}, passé composé).
    \item \textbf{Attention :} \esp{Acostumbrar} (sans \emph{a}) + infinitif existe en Amérique latine (\esp{Acostumbro levantarme temprano}), mais la norme péninsulaire et l'examen exigent \textbf{\esp{acostumbrar a} + infinitif}.
\end{itemize}
\end{propriete}

\begin{propriete}[3. \esp{Tener la costumbre de} + infinitif]
Locution figée. \esp{Tener} se conjugue normalement. Registre plus formel / descriptif.
\begin{itemize}
    \item \esp{Tengo la costumbre de cambiar mis contraseñas cada mes.}
    \item \esp{Los jóvenes tienen la costumbre de comunicarse por WhatsApp.}
\end{itemize}
\end{propriete}

\begin{propriete}[4. L'Imparfait indicatif (valeur aspectuelle d'habitude passée)]
Utilisé pour décrire une \textbf{habitude passée dans un contexte narratif ou descriptif}, souvent avec des marqueurs temporels (\esp{antes, antiguamente, cuando era niño, cada día}).
\begin{itemize}
    \item \esp{Cuando era pequeño, \textbf{jugaba} en la calle, no \textbf{miraba} la pantalla.} (Habitudes répétées dans le passé).
    \item \esp{Mi abuelo \textbf{leía} el periódico cada mañana.} (Rituel passé).
\end{itemize}
\textbf{Contraste \esp{Solía} vs \esp{Imparfait} :} \esp{Solía} met l'accent sur la \textbf{disparition} de l'habitude (« j'avais l'habitude de... mais plus maintenant »). L'imparfait décrit le \textbf{cadre de vie} passé (« je faisais... »).
\end{propriete}

\begin{attention}[Pièges francophones : Habitude]
\begin{itemize}
    \item \textbf{*Suelo a + infinitif} $\rightarrow$ \textbf{FAUX}. \esp{Soler} + infinitif \textbf{SANS préposition}. \esp{Suelo navegar.}
    \item \textbf{*Acostumbrar + infinitif (sans \emph{a})} $\rightarrow$ \textbf{FAUX} en norme standard (correct en AmL seulement). Écrivez \esp{Acostumbro a navegar.}
    \item \textbf{*Tener costumbre de} (sans article) $\rightarrow$ \textbf{FAUX}. Il faut \esp{la} costumbre.
    \item \textbf{Confusion \esp{Suelo} (verbe) / \esp{Suelo} (nom = sol/parquet)}. Contexte : \esp{Suelo ir} (verbe) vs \esp{El suelo está sucio} (nom).
    \item \textbf{\esp{Ya no} + \esp{soler}} : \esp{Ya no suelo ir.} = « Je n'ai plus l'habitude d'y aller ». \esp{Ya no} nie l'habitude actuelle.
\end{itemize}
\end{attention}

\begin{exemplebox}[Exemples traduits : Habitude (Contexte NTIC)]
\begin{itemize}
    \item \esp{Suelo verificar las fuentes antes de compartir.} $\rightarrow$ « J'ai l'habitude de vérifier les sources avant de partager. » (Présent).
    \item \esp{Antes solía pasar horas en TikTok, ahora lo limito.} $\rightarrow$ « Avant, j'avais l'habitude de passer des heures sur TikTok, maintenant je limite. » (Passé révolu).
    \item \esp{Los estudiantes se acostumbran rápido a las nuevas apps.} $\rightarrow$ « Les étudiants s'habituent vite aux nouvelles applis. » (Processus \esp{acostumbrarse}).
    \item \esp{Tenemos la costumbre de hacer copias de seguridad los domingos.} $\rightarrow$ « Nous avons l'habitude de faire des sauvegardes les dimanches. »
    \item \esp{En mi pueblo, antes no \textbf{había} internet; la gente \textbf{hablaba} en la plaza.} $\rightarrow$ « Dans mon village, avant il n'y avait pas internet ; les gens parlaient sur la place. » (Imparfait descriptif).
\end{itemize}
\end{exemplebox}

\begin{exoresolu}[Exercice d'application : Habitude]
\textbf{Énoncé :} Traduisez en espagnol en choisissant la structure la plus appropriée.
\begin{enumerate}
    \item J'ai l'habitude de vérifier mes mails le matin. (Habitude actuelle)
    \item Autrefois, les jeunes jouaient dehors. (Description passée)
    \item Nous n'avons plus l'habitude d'écrire des lettres. (Habitude perdue)
    \item Il s'habitue à travailler avec une tablette. (Processus d'adaptation)
    \item Ils ont la coutume de se déconnecter le week-end. (Habitude établie, formel)
\end{enumerate}

\tcblower

\textbf{Solution détaillée :}
\begin{enumerate}
    \item \esp{Suelo revisar mis correos por la mañana.} (\esp{Soler} + inf présent).
    \item \esp{Antiguamente, los jóvenes \textbf{jugaban} fuera.} (Imparfait descriptif/narratif).
    \item \esp{Ya no solemos escribir cartas.} (\esp{Ya no} + \esp{soler} présent pour habitude perdue). Ou \esp{Antes solíamos escribir cartas, ya no.}
    \item \esp{Se está acostumbrando a trabajar con una tableta.} (\esp{Acostumbrarse} + gérondif pour processus en cours) ou \esp{Se acostumbra a trabajar...} (présent).
    \item \esp{Tienen la costumbre de desconectarse el fin de semana.} (\esp{Tener la costumbre de}).
\end{enumerate}
\end{exoresolu}

\begin{exercice}[Entraînement mixte : Préférence et Habitude]
Complétez les phrases avec la structure indiquée entre parenthèses.
\begin{enumerate}
    \item \_\_\_\_\_\_ (Preferir) yo usar contraseñas largas \_\_\_\_\_\_ cortas.
    \item \_\_\_\_\_\_ (Gustar más) \_\_\_\_\_\_ (nosotros) leer en papel \_\_\_\_\_\_ en pantalla.
    \item \_\_\_\_\_\_ (Soler) \_\_\_\_\_\_ (yo) cambiar la contraseña cada mes.
    \item \_\_\_\_\_\_ (Acostumbrar a) \_\_\_\_\_\_ (ellos) navegar sin antivirus. (Danger !)
    \item \_\_\_\_\_\_ (Tener costumbre de) \_\_\_\_\_\_ (mi abuelo) contar historias, no ver videos.
    \item \esp{Antes que} \_\_\_\_\_\_ (descargar / tú) un archivo, \_\_\_\_\_\_ (verificar / tú) la fuente.
\end{enumerate}
\end{exercice}

\begin{corrige}[Exercice : Préférence et Habitude]
\begin{enumerate}
    \item \esp{Prefiero ... a} (\esp{Prefiero usar contraseñas largas a cortas}).
    \item \esp{Nos gusta más ... que} (\esp{Nos gusta más leer en papel que en pantalla}).
    \item \esp{Suelo} (\esp{Suelo cambiar la contraseña cada mes}).
    \item \esp{Acostumbran a navegar} (\esp{Acostumbran a navegar sin antivirus}).
    \item \esp{Tiene la costumbre de contar} (\esp{Mi abuelo tiene la costumbre de contar historias...}).
    \item \esp{Antes de que descargues, verifica} (Sujet différent \emph{tú} / \emph{tú} implicite impératif $\rightarrow$ \esp{Antes de que descargues} subjonctif / \esp{verifica} impératif). Ou si même sujet : \esp{Antes de descargar, verifica.}
\end{enumerate}
\end{corrige}

\coursec{Méthode du commentaire et production guidée}

\courssub{Production écrite / orale guidée : Question personnelle}

\begin{experience}[Production orale / écrite guidée : Question personnelle]
\textbf{Consigne :} Répondez en espagnol (60–90 mots) à la question suivante, en utilisant le vocabulaire du texte \textbf{ET} les structures de préférence (section 4) et d'habitude (section 5). Soyez cohérent et personnel.
\esp{« ¿Cómo usas Internet en tu vida diaria? ¿Qué prefieres y qué hábitos tienes para un uso responsable? »}

\textbf{Pistes de réponse (modèle corrigé et enrichi) :}
\esp{« En mi vida diaria, \textbf{suelo} usar Internet sobre todo para estudiar y comunicarme con mis compañeros. \textbf{Prefiero} consultar páginas educativas \textbf{a} perder el tiempo en redes sociales sin control. \textbf{Me gusta más} leer artículos completos \textbf{que} ver solo videos cortos. Para ser responsable, \textbf{tengo la costumbre de} verificar las fuentes antes de compartir nada y \textbf{acostumbro a} cambiar mis contraseñas cada trimestre. Además, \textbf{limito} mi tiempo de pantalla antes de dormir para evitar la adicción y dormir bien. »}

\textbf{Traduction :} « Dans ma vie quotidienne, j'ai l'habitude d'utiliser Internet surtout pour étudier et communiquer avec mes camarades. Je préfère consulter des sites éducatifs plutôt que perdre du temps sur les réseaux sociaux sans contrôle. J'aime mieux lire des articles complets que regarder seulement des vidéos courtes. Pour être responsable, j'ai la coutume de vérifier les sources avant de partager quoi que ce soit et j'ai l'habitude de changer mes mots de passe chaque trimestre. De plus, je limite mon temps d'écran avant de dormir pour éviter l'addiction et bien dormir. »
\end{experience}

\courssub{Synthèse des attendus pour l'épreuve de compréhension et d'expression}

\begin{aretenir}[Check-list Bac — Compréhension \& Expression]
\begin{itemize}
    \item \textbf{Lecture active :} Surlignez mots-clés, connecteurs (\esp{sin embargo, por tanto, además}), marques d'opinion (\esp{advierten que, la solución es}).
    \item \textbf{Questions globales :} Sujet + Thèse + Ton = points faciles si structurés (Phrase nominale pour le sujet, phrase complète pour la thèse).
    \item \textbf{Vrai/Faux :} \textbf{Citation obligatoire} \esp{« ... »} + explication en français. Pas de citation = 0 point.
    \item \textbf{Vocabulaire en contexte :} Cherchez le mot \textbf{dans le texte}, respectez la catégorie grammaticale (verbe, nom, adjectif).
    \item \textbf{Question personnelle / Expression :} Écrivez \textbf{en espagnol}. Utilisez le lexique du texte \textbf{ET} les structures grammaticales ciblées (\esp{preferir, gustar más, soler, acostumbrar a, tener la costumbre de, imparfait}). Soignez l'orthographe (accents, \esp{ñ}, \esp{¿¡}).
    \item \textbf{Résumé :} 3 phrases = 3 macro-idées. Vos mots. Connecteurs logiques. Limite de mots respectée.
    \item \textbf{Traduction (Thème/Version) :} Maîtrisez les correspondances : \esp{preferir A a B} / \esp{me gusta más A que B} / \esp{suelo + inf} / \esp{solía + inf} / \esp{acostumbro a + inf} / \esp{tengo la costumbre de + inf}.
\end{itemize}
\end{aretenir}

\begin{savaistu}[Le Cameroun, la Guinée équatoriale et la \cle{brecha digital}]
Le terme \esp{brecha digital} (litt. « brèche/fissure numérique ») désigne l'inégalité d'accès aux NTIC entre zones urbaines/rurales, générations, niveaux de revenus. Au Cameroun, selon les rapports de l'ARTEC (Agence de Régulation des Télécommunications), le taux de pénétration d'Internet est \textbf{significativement plus élevé dans les grandes métropoles} (Yaoundé, Douala) que dans les régions du Nord, de l'Est ou de l'Adamaoua, où l'infrastructure fait défaut et le coût des données reste prohibitif pour beaucoup. Des initiatives comme \emph{« Un élève, un ordinateur »} ou le déploiement du \emph{Backbone} national (fibre optique) visent à la réduire.
En \textbf{Guinée équatoriale} (seul pays africain hispanophone officiel), le défi est analogue : connexion concentrée à Malabo (île de Bioko) et Bata (continent), coût élevé, débit souvent faible. L'espagnol y est la langue de l'administration, de l'éducation et des débats sur la \esp{ciudadanía digital} (citoyenneté numérique) et la \esp{protección de datos}. La maîtrise de l'espagnol ouvre donc l'accès à cet espace numérique régional en construction.
\end{savaistu}

\coursec{Lexique thématique des NTIC}

La lecture du texte \esp{Internet y las NTIC} nous a permis de comprendre les grandes idées : les nouvelles technologies transforment la communication, l'école et la vie quotidienne, de Madrid à Yaoundé. Mais pour \emph{parler} de ces sujets à l'écrit comme à l'oral — et surtout pour réussir la version et le thème au baccalauréat —, il faut maintenant posséder un vocabulaire précis et actif. Cette section organise ce lexique en quatre champs sémantiques clairs : les \cle{herramientas} (les outils), les \cle{acciones} (les actions), les \cle{problemas} (les problèmes et enjeux) et les \cle{ventajas} (les avantages). Apprendre un mot, c'est apprendre sa famille, son genre, sa prononciation et les verbes qui l'accompagnent : c'est exactement ce que nous allons faire.

\courssub{Champ 1 : las herramientas (les outils)}

\begin{definition}[Les outils numériques]
Une \esp{herramienta} est un outil. Dans le domaine des NTIC (\esp{nuevas tecnologías de la información y la comunicación}), ce sont les objets et services qui permettent d'accéder au monde numérique.
\end{definition}

Voici le vocabulaire de base, avec le genre grammatical — toujours à mémoriser avec l'article :

\begin{center}
\begin{tabularx}{\linewidth}{|l|X|X|}
\hline
\rowcolor{popblueL} Español & Français & Exemple et traduction \\
\hline
\esp{Internet} (m., souvent sans article) & Internet & \esp{Uso Internet todos los días.} « J'utilise Internet tous les jours. » \\
\hline
\esp{el teléfono inteligente} & le smartphone & \esp{Mi teléfono inteligente es nuevo.} « Mon smartphone est neuf. » \\
\hline
\esp{el ordenador} (Esp.) / \esp{la computadora} (Am.) & l'ordinateur & \esp{Trabajo con el ordenador.} « Je travaille avec l'ordinateur. » \\
\hline
\esp{la tableta} & la tablette & \esp{Los alumnos leen en la tableta.} « Les élèves lisent sur la tablette. » \\
\hline
\esp{la red social} & le réseau social & \esp{WhatsApp es una red social muy usada.} « WhatsApp est un réseau social très utilisé. » \\
\hline
\esp{la contraseña} & le mot de passe & \esp{No recuerdo mi contraseña.} « Je ne me souviens pas de mon mot de passe. » \\
\hline
\esp{el correo electrónico} & le courriel / l'e-mail & \esp{Envíame un correo electrónico.} « Envoie-moi un courriel. » \\
\hline
\esp{la pantalla} & l'écran & \esp{La pantalla de mi tableta es pequeña.} « L'écran de ma tablette est petit. » \\
\hline
\end{tabularx}
\end{center}

\begin{exemplebox}[Exemples en contexte]
\esp{Navego por Internet con mi teléfono inteligente.} « Je navigue sur Internet avec mon smartphone. »

\esp{Protege tu contraseña.} « Protège ton mot de passe. »

\esp{En el cibercafé de Douala hay diez ordenadores.} « Dans le cybercafé de Douala, il y a dix ordinateurs. »
\end{exemplebox}

\begin{attention}[Faux amis et pièges de genre]
\begin{itemize}
\item \esp{una red social} = un réseau social. Ne dites jamais \esp{una red sociale} : l'adjectif \esp{social} est invariable en genre en espagnol (\esp{la red social}, \esp{el problema social}).
\item \esp{la contraseña} : attention au tilde sur le \esp{ñ} ; ne pas écrire \esp{contrasena}. Le pluriel est \esp{las contraseñas}.
\item \esp{descargar} se prononce avec un [g] dur devant \esp{a} : \emph{des-car-gar}. Devant \esp{e}, le verbe devient \esp{descargue} (subjonctif présent) pour garder le même son.
\item \esp{el ordenador} est masculin ; en Amérique latine on dit plutôt \esp{la computadora} (ou \esp{el computador} en Colombie), féminin dans la plupart des pays.
\item \esp{la pantalla} = l'écran ; ne confondez pas avec le français « pantalon » !
\end{itemize}
\end{attention}

\begin{savaistu}[Móvil ou celular ?]
En Espagne, on dit couramment \esp{el móvil} pour « le téléphone portable » ; en Amérique latine et en Guinée équatoriale, on dit \esp{el celular}. Les deux sont corrects : \esp{teléfono inteligente} désigne précisément le smartphone. À la version du Bac, si vous rencontrez \esp{celular}, traduisez simplement « téléphone portable ». De même, \esp{el wifi} (ou \esp{la wifi} en Espagne) se traduit par « le wifi ».
\end{savaistu}

\courssub{Champ 2 : las acciones (les actions)}

Les outils ne servent à rien sans les verbes qui décrivent ce qu'on en fait. Observez d'abord ces phrases :

\begin{exemplebox}[Observer avant de mémoriser]
\esp{Cada mañana, Amina se conecta a Internet, busca información para sus deberes y chatea con sus amigas. Por la tarde, descarga una canción, sube fotos a su red social y comparte un vídeo con su hermano.}

« Chaque matin, Amina se connecte à Internet, cherche des informations pour ses devoirs et tchatche avec ses amies. L'après-midi, elle télécharge une chanson, met en ligne des photos sur son réseau social et partage une vidéo avec son frère. »
\end{exemplebox}

\begin{propriete}[Les verbes des NTIC : construction et régime]
\begin{itemize}
\item \esp{navegar por Internet} : naviguer sur Internet. La préposition est \esp{por}, jamais \esp{en} : \esp{navego por la red} « je navigue sur le Net ».
\item \esp{descargar} : télécharger (vers soi) ; \esp{subir} : téléverser, mettre en ligne. \esp{Descargo un documento} / \esp{Subo mis fotos}.
\item \esp{conectarse} : se connecter (verbe pronominal : \esp{me conecto, te conectas, se conecta...}) ; son contraire : \esp{desconectarse}.
\item \esp{chatear} : tchatter, discuter en ligne ; \esp{compartir} : partager ; \esp{publicar} : publier (\esp{publicar un mensaje}).
\item \esp{buscar información} : chercher des informations. Attention : \esp{buscar} se construit \emph{sans} préposition : \esp{busco información} « je cherche des informations » (et non \esp{busco por información}).
\end{itemize}
\end{propriete}

\begin{attention}[Verbes réguliers mais orthographe à surveiller]
Tous ces verbes sont du premier groupe (\esp{-ar}), sauf \esp{subir} et \esp{compartir} (\esp{-ir}), tous réguliers. Seule difficulté : \esp{buscar} (c $\to$ qu devant e : \esp{busqué}, \esp{busque}) et \esp{descargar} (g $\to$ gu devant e : \esp{descargué}, \esp{descargue}). N'écrivez pas \esp{chatter} avec deux t : l'espagnol a adapté le mot anglais en \esp{chatear} (\esp{chateo, chateas, chatea...}).
\end{attention}

\courssub{Champ 3 : los problemas y los desafíos (problèmes et enjeux)}

\begin{definition}[La brecha digital]
\esp{La brecha digital} désigne l'écart entre les populations qui ont accès aux nouvelles technologies (connexion, équipement, formation) et celles qui n'y ont pas accès. On parle de \esp{brecha digital} entre la ville et la campagne, entre les riches et les pauvres, entre les jeunes et les personnes âgées : \esp{En muchos países, la brecha digital separa a los jóvenes urbanos de los campesinos.} « Dans de nombreux pays, la fracture numérique sépare les jeunes des villes des paysans. »
\end{definition}

\begin{center}
\begin{tabularx}{\linewidth}{|l|X|X|}
\hline
\rowcolor{popblueL} Español & Français & Exemple et traduction \\
\hline
\esp{la brecha digital} & la fracture numérique & \esp{Hay que reducir la brecha digital.} « Il faut réduire la fracture numérique. » \\
\hline
\esp{la adicción (a Internet)} & l'addiction (à Internet) & \esp{Algunos jóvenes sufren adicción a las redes.} « Certains jeunes souffrent d'addiction aux réseaux. » \\
\hline
\esp{el robo de datos} & le vol de données & \esp{El robo de datos es un delito.} « Le vol de données est un délit. » \\
\hline
\esp{las noticias falsas} & les fausses informations (fake news) & \esp{No compartas noticias falsas.} « Ne partage pas de fausses informations. » \\
\hline
\esp{la privacidad} & la vie privée, la confidentialité & \esp{Las redes sociales amenazan nuestra privacidad.} « Les réseaux sociaux menacent notre vie privée. » \\
\hline
\esp{el correo basura} & le courrier indésirable (spam) & \esp{Mi bandeja está llena de correo basura.} « Ma boîte est pleine de courriers indésirables. » \\
\hline
\end{tabularx}
\end{center}

\begin{attention}[Pièges de traduction]
\begin{itemize}
\item \esp{privacidad} se traduit par « vie privée » ou « confidentialité », et non par « privacité » (mot qui n'existe pas en français standard).
\item « Fake news » se dit \esp{noticias falsas} ou \esp{bulos} ; le mot anglais est compris mais évitez-le dans une copie.
\item \esp{un delito} = un délit ; ne confondez pas avec \esp{delicioso} (délicieux) !
\item \esp{la adicción} prend deux \esp{c} : \esp{a-dic-ción} ; son adjectif est \esp{adicto/a} : \esp{Es adicto a las redes.} « Il est accro aux réseaux. »
\item \esp{los datos} = les données ; \esp{un dato} = une donnée. Ne traduisez pas par « dates » : « une date » (calendrier) se dit \esp{una fecha}.
\end{itemize}
\end{attention}

\courssub{Champ 4 : las ventajas (les avantages)}

Heureusement, les NTIC ne sont pas que des problèmes ! Pour un essai ou un débat équilibré (exercice fréquent au Bac), il faut aussi le lexique positif :

\begin{itemize}
\item \esp{comunicarse con los demás} : communiquer avec les autres — \esp{Internet nos permite comunicarnos con el mundo entero.} « Internet nous permet de communiquer avec le monde entier. »
\item \esp{aprender en línea} : apprendre en ligne — \esp{Durante las vacaciones, los alumnos pueden aprender en línea.} « Pendant les vacances, les élèves peuvent apprendre en ligne. »
\item \esp{el acceso a la información} : l'accès à l'information — \esp{Internet ofrece un acceso rápido a la información.} « Internet offre un accès rapide à l'information. »
\item \esp{un puente hacia el saber} : un pont vers le savoir — \esp{Para un estudiante de Yaoundé, Internet es un puente hacia el saber.} « Pour un élève de Yaoundé, Internet est un pont vers le savoir. »
\end{itemize}

\begin{propriete}[Expressions utiles à retenir par cœur]
\begin{itemize}
\item \esp{estar conectado} : être connecté — \esp{Los jóvenes están conectados día y noche.} « Les jeunes sont connectés jour et nuit. » (participe employé avec \esp{estar} : état résultant d'une action ; l'accord se fait avec le sujet : \esp{están conectado\textbf{s}}).
\item \esp{usar las redes sociales} : utiliser les réseaux sociaux — \esp{Mis primos usan las redes sociales para vender cacao.} « Mes cousins utilisent les réseaux sociaux pour vendre du cacao. »
\item \esp{proteger su contraseña} : protéger son mot de passe — \esp{Es importante proteger su contraseña.} « Il est important de protéger son mot de passe. »
\end{itemize}
\end{propriete}

\courssub{Carte mentale et synthèse des quatre champs}

\begin{popfigure}
\begin{tikzpicture}[mindmap, grow cyclic, every node/.style={concept, align=center, font=\small},
root concept/.append style={concept color=popblue, font=\bfseries},
level 1/.append style={level distance=3.4cm, sibling angle=90},
level 2/.append style={level distance=2.4cm, sibling angle=45, font=\scriptsize}]
\node[root concept]{Las NTIC}
  child[concept color=popgreen] { node {Herramientas}
    child { node[align=center]{teléfono\\inteligente} }
    child { node[align=center]{ordenador,\\tableta} }
    child { node[align=center]{red social,\\contraseña} }
  }
  child[concept color=poporange] { node {Acciones}
    child { node[align=center]{navegar,\\buscar} }
    child { node[align=center]{descargar,\\subir} }
    child { node[align=center]{chatear,\\compartir} }
  }
  child[concept color=poppink] { node {Problemas}
    child { node[align=center]{brecha\\digital} }
    child { node[align=center]{adicción,\\robo de datos} }
    child { node[align=center]{noticias\\falsas} }
  }
  child[concept color=poppurple] { node {Ventajas}
    child { node {comunicarse} }
    child { node[align=center]{aprender\\en línea} }
    child { node[align=center]{puente hacia\\el saber} }
  };
\end{tikzpicture}
\legende{Carte mentale du lexique des NTIC : quatre champs sémantiques à mémoriser ensemble.}
\end{popfigure}

\begin{popfigure}
\begin{tikzpicture}
\node[draw=popblue, fill=popblueL, rounded corners, font=\bfseries, minimum width=3.2cm, minimum height=0.9cm] (h) at (0,4) {Herramientas};
\node[draw=poporange, fill=poporangeL, rounded corners, font=\bfseries, minimum width=3.2cm, minimum height=0.9cm] (a) at (4.2,4) {Acciones};
\node[draw=poppink, fill=poppinkL, rounded corners, font=\bfseries, minimum width=3.2cm, minimum height=0.9cm] (p) at (8.4,4) {Problemas};
\node[draw=poppurple, fill=poppurpleL, rounded corners, font=\bfseries, minimum width=3.2cm, minimum height=0.9cm] (v) at (12.6,4) {Ventajas};
\node[align=left, font=\scriptsize] at (0,2) {Internet\\el teléfono inteligente\\el ordenador\\la tableta\\la red social\\la contraseña\\el correo electrónico};
\node[align=left, font=\scriptsize] at (4.2,2) {navegar por Internet\\descargar, subir\\chatear\\conectarse\\compartir, publicar\\buscar información};
\node[align=left, font=\scriptsize] at (8.4,2) {la brecha digital\\la adicción\\el robo de datos\\las noticias falsas\\la privacidad};
\node[align=left, font=\scriptsize] at (12.6,2) {comunicarse\\aprender en línea\\el acceso a la\\información\\el puente hacia\\el saber};
\draw[popline] (-1.8,3.5) -- (14.4,3.5);
\draw[popline] (-1.8,0.4) -- (14.4,0.4);
\end{tikzpicture}
\legende{Tableau de classement : Herramientas / Acciones / Problemas / Ventajas.}
\end{popfigure}

\courssub{Texte de réinvestissement et entraînement}

\begin{textebox}[Un debate en el instituto]
En la clase de español del Lycée de Nkol-Eton, la profesora propone un debate: «¿Internet: puente hacia el saber o máquina de perder el tiempo?». Moussa toma la palabra primero: «Para mí, Internet es una herramienta magnífica. Cada día navego por Internet con mi teléfono inteligente, busco información para mis trabajos y aprendo en línea con vídeos educativos.» Pero Nadège no está de acuerdo: «Es verdad, pero muchos jóvenes sufren adicción a las redes sociales. Chatean toda la noche, comparten noticias falsas y no protegen su contraseña. Además, el robo de datos es un peligro real.» Entonces interviene el pequeño Bruno: «En mi barrio, muchas familias no tienen ordenador ni tableta: la brecha digital es enorme. Los que tienen acceso a la información progresan; los otros se quedan atrás.» La profesora concluye: «Internet es como un cuchillo: todo depende del uso que hacemos de él. Usad las redes sociales con inteligencia, verificad las noticias y proteged vuestra privacidad. Así, Internet será de verdad un puente hacia el saber para todos los jóvenes de Camerún.»
\end{textebox}

\textbf{Glossaire.} \esp{el debate} : le débat ; \esp{perder el tiempo} : perdre son temps ; \esp{magnífica} : magnifique ; \esp{no está de acuerdo} : n'est pas d'accord ; \esp{el peligro} : le danger ; \esp{el barrio} : le quartier ; \esp{quedarse atrás} : rester à la traîne ; \esp{el cuchillo} : le couteau ; \esp{verificad} : vérifiez (impératif, 2\up{e} pers. du pluriel).

\textbf{Questions de compréhension.}
\begin{enumerate}
\item ¿Qué ventajas de Internet menciona Moussa?
\item Según Nadège, ¿qué problemas causan las redes sociales?
\item ¿Qué es la brecha digital según Bruno?
\item ¿Con qué compara la profesora a Internet? ¿Por qué?
\end{enumerate}

\begin{exoresolu}[Classement dans les quatre champs]
Classez les mots suivants dans les quatre colonnes : \esp{descargar, la privacidad, la tableta, chatear, la brecha digital, el correo electrónico, aprender en línea, compartir, el robo de datos, la contraseña, comunicarse, las noticias falsas}.
\tcblower
\textbf{Herramientas} : \esp{la tableta, el correo electrónico, la contraseña} (objets et services). \textbf{Acciones} : \esp{descargar, chatear, compartir} (verbes d'activité). \textbf{Problemas} : \esp{la privacidad} (enjeu menacé), \esp{la brecha digital, el robo de datos, las noticias falsas}. \textbf{Ventajas} : \esp{aprender en línea, comunicarse}. Astuce de méthode : un verbe à l'infinitif appartient presque toujours au champ des actions ou des avantages ; un nom abstrait négatif (\esp{robo, adicción, brecha}) appartient au champ des problèmes.
\end{exoresolu}

\begin{exercice}[À vous de jouer]
\begin{enumerate}
\item Traduisez : a) Je navigue sur Internet avec mon smartphone. b) Protège ton mot de passe. c) Les fausses informations menacent notre vie privée. d) Internet est un pont vers le savoir.
\item Complétez avec le mot qui convient : \esp{brecha digital / contraseña / descargar / noticias falsas} : a) No debes compartir tu \dots{} con nadie. b) Antes de compartir una noticia, hay que verificar si no es una de esas \dots{}. c) Quiero \dots{} esta canción camerunesa. d) En el campo, pocos alumnos tienen Internet: es la \dots{}.
\end{enumerate}
\end{exercice}

\begin{corrige}[Exercice « À vous de jouer »]
\textbf{1.} a) \esp{Navego por Internet con mi teléfono inteligente.} b) \esp{Protege tu contraseña.} c) \esp{Las noticias falsas amenazan nuestra privacidad.} d) \esp{Internet es un puente hacia el saber.}
\textbf{2.} a) \esp{contraseña} ; b) \esp{noticias falsas} ; c) \esp{descargar} ; d) \esp{brecha digital}.
\end{corrige}

\begin{aretenir}[L'essentiel du lexique NTIC]
Quatre champs à associer dans votre mémoire : les \cle{herramientas} (\esp{teléfono inteligente, ordenador, red social, contraseña, correo electrónico}), les \cle{acciones} (\esp{navegar por, descargar, subir, chatear, conectarse, compartir, buscar información}), les \cle{problemas} (\esp{brecha digital, adicción, robo de datos, noticias falsas, privacidad}) et les \cle{ventajas} (\esp{comunicarse, aprender en línea, acceso a la información, puente hacia el saber}). Trois réflexes : \esp{navegar \textbf{por} Internet}, \esp{buscar} sans préposition, et \esp{una red social} sans « e » final.
\end{aretenir}

\coursec{Grammaire 1 : l'expression de la préférence}

Après avoir exploré le lexique spécifique aux NTIC dans la section précédente, nous disposons maintenant du vocabulaire nécessaire (\cle{red social}, \cle{navegar}, \cle{descargar}, \cle{teléfono inteligente}, etc.) pour construire des énoncés complets. L'expression de la préférence est une compétence communicative fondamentale : elle permet de justifier ses choix numériques, de comparer des applications ou de débattre de l'impact de la technologie. En espagnol, deux structures principales dominent : le verbe \esp{preferir} (verbe lexical, transitif) et la construction \esp{gustar más} (verbe à construction inverse). Nous les observerons d'abord dans un dialogue authentique, nous les analyserons, puis nous les mettrons en parallèle avec le français avant de les mobiliser dans la traduction.

\courssub{Observation : préférences numériques dans un dialogue authentique}

Commençons par observer comment de jeunes locuteurs expriment leurs goûts en matière d'applications et d'usage d'Internet. Le dialogue suivant met en scène deux lycéens camerounais, \emph{Armand} et \emph{Blanche}, qui discutent de leurs habitudes de connexion au cybercafé du quartier Mvog-Ada à Yaoundé. Repérez, pendant la lecture, toutes les formes qui expriment un choix ou un goût.

\begin{textebox}[Dialogue : Choix d'applications et de navigation]
— \esp{¿Qué prefieres, Armand, WhatsApp o Telegram para hablar con tus amigos?}

— \esp{Prefiero WhatsApp a Telegram. Todos mis contactos están allí. Y a ti, ¿qué te gusta más?}

— \esp{A mí me gusta más Telegram por la privacidad. Pero para ver videos, prefiero YouTube a TikTok.}

— \esp{Entiendo. Yo suelo preferir leer noticias en el navegador antes que ver videos.}

— \esp{¿Prefieres la lectura a los videos? ¡Qué raro! A la mayoría le gusta más ver contenido visual.}

— \esp{Cada uno tiene sus costumbres. Prefiero informarme bien antes que perder el tiempo.}
\end{textebox}

\noindent
\textbf{Glossaire :} \esp{todos mis contactos} = tous mes contacts ; \esp{la privacidad} = la vie privée, la confidentialité ; \esp{el navegador} = le navigateur (web) ; \esp{perder el tiempo} = perdre son temps ; \esp{informarse} = s'informer ; \esp{cada uno tiene sus costumbres} = chacun a ses habitudes (« à chacun ses goûts »).

\noindent
\textbf{Relevé d'observation :} \esp{¿qué prefieres?}, \esp{prefiero WhatsApp a Telegram}, \esp{¿qué te gusta más?}, \esp{a mí me gusta más Telegram}, \esp{prefiero YouTube a TikTok}, \esp{prefiero informarme antes que perder el tiempo}. Deux verbes reviennent sans cesse : \cle{preferir} et \cle{gustar más}. C'est précisément ce que nous allons étudier.

\courssub{Le verbe \esp{preferir} : diphtongaison et conjugaison au présent}

Le verbe \cle{preferir} (préférer) est un verbe du troisième groupe (\cle{-ir}) qui présente une \textbf{diphtongaison} systématique à l'indicatif présent : la voyelle tonique \cle{e} du radical se transforme en \cle{ie} pour les trois personnes du singulier et la troisième personne du pluriel. Les personnes \cle{nosotros} et \cle{vosotros} échappent à cette alternance car l'accent tonique ne porte pas sur le radical mais sur la désinence.

\begin{propriete}[Diphtongaison \textit{e} $\rightarrow$ \textit{ie} dans \textit{preferir}]
Le verbe \esp{preferir} appartient au groupe des verbes à diphtongaison \textbf{e $\rightarrow$ ie} (comme \esp{querer}, \esp{pensar}, \esp{sentir}, \esp{cerrar}).
\begin{itemize}
    \item \textbf{Personnes affectées (accent tonique sur le radical) :} \esp{yo pref\underline{ie}ro}, \esp{tú pref\underline{ie}res}, \esp{él/ella/usted pref\underline{ie}re}, \esp{ellos/ellas/ustedes pref\underline{ie}ren}.
    \item \textbf{Personnes non affectées (accent tonique sur la désinence) :} \esp{nosotros preferimos}, \esp{vosotros preferís}.
\end{itemize}
Cette régularité phonétique est sans exception au présent de l'indicatif. \emph{Remarque :} au passé simple et au subjonctif imparfait, \esp{preferir} présente en outre l'alternance \textbf{e $\rightarrow$ i} à la 3e personne (\esp{prefirió}, \esp{prefirieron}), comme tous les verbes en \esp{-ir} à diphtongaison.
\end{propriete}

\begin{popfigure}
\begin{tikzpicture}[
    cell/.style={rectangle, draw=popline, minimum height=1.1cm, align=center, font=\small},
    header/.style={cell, fill=popblueL, font=\bfseries\small, text=popdark}
]
\matrix (conj) [matrix of nodes, nodes in empty cells, column sep=-\pgflinewidth, row sep=-\pgflinewidth,
    column 1/.style={nodes={cell, minimum width=3.2cm, fill=popcream}},
    column 2/.style={nodes={cell, minimum width=3.2cm}},
    column 3/.style={nodes={cell, minimum width=6.2cm, fill=popblue!10}}
]
{
|[header]| Persona & |[header]| Conjugación & |[header]| Observación \\
Yo & pref\textcolor{popink}{\textbf{ie}}ro & Diphtongaison \textcolor{popink}{\textbf{ie}} (accent sur le radical) \\
Tú & pref\textcolor{popink}{\textbf{ie}}res & Diphtongaison \textcolor{popink}{\textbf{ie}} \\
Él / Ella / Usted & pref\textcolor{popink}{\textbf{ie}}re & Diphtongaison \textcolor{popink}{\textbf{ie}} \\
Nosotros / Nosotras & preferimos & \textbf{Pas de diphtongaison} (radical \emph{prefer-}) \\
Vosotros / Vosotras & preferís & \textbf{Pas de diphtongaison} (accent sur \emph{-ís}) \\
Ellos / Ellas / Ustedes & pref\textcolor{popink}{\textbf{ie}}ren & Diphtongaison \textcolor{popink}{\textbf{ie}} \\
};
\end{tikzpicture}
\legende{Conjugaison de \textit{preferir} au présent de l'indicatif : la diphtongaison \textit{e $\rightarrow$ ie} est signalée en couleur.}
\end{popfigure}

\begin{exemplebox}[Exemples avec \textit{preferir} + infinitif / nom]
\esp{Prefiero navegar en mi teléfono inteligente.} « Je préfère naviguer sur mon smartphone. » (préférer + infinitif, sans préposition)

\esp{Prefiero la lectura a los videojuegos.} « Je préfère la lecture aux jeux vidéo. » (préférer + nom)

\esp{¿Qué prefieres, descargar la aplicación o usar la web?} « Que préfères-tu, télécharger l'application ou utiliser le site web ? »

\esp{Mi hermano prefiere Instagram a Facebook.} « Mon frère préfère Instagram à Facebook. »
\end{exemplebox}

\courssub{La structure comparative : \esp{preferir X a Y}}

En espagnol, le second terme de la comparaison s'introduit par la préposition \cle{a} (équivalent de « à » en français). Rappelons la contraction obligatoire devant l'article masculin singulier : \esp{a + el = al}.

\begin{propriete}[Structure \textit{preferir} + nom/infinitif + \textit{a} + nom/infinitif]
\begin{center}
\textbf{Sujet} + \textbf{preferir (conjugué)} + \textbf{élément préféré} + \textbf{a} + \textbf{élément rejeté}
\end{center}
\begin{itemize}
    \item \esp{Prefiero el correo electrónico \textbf{a} los mensajes de texto.} « Je préfère le courriel aux SMS. »
    \item \esp{Prefiero estudiar en línea \textbf{a} ir a la biblioteca.} « Je préfère étudier en ligne plutôt qu'aller à la bibliothèque. »
    \item \esp{Prefiero \textbf{esto} \textbf{a} \textbf{eso}.} « Je préfère ceci à cela. »
\end{itemize}
\emph{Attention :} devant un nom masculin singulier défini, contraction obligatoire : \esp{Prefiero el ordenador \textbf{al} móvil.} « Je préfère l'ordinateur au portable. »
\end{propriete}

\begin{center}
\begin{tabularx}{\linewidth}{|X|X|X|}
\hline
\rowcolor{popblueL} \textbf{Français} & \textbf{Español} & \textbf{Remarque} \\ \hline
Je préfère le café \textbf{au} thé. & \esp{Prefiero el café \textbf{al} té.} & \esp{a + el = al} (masc. sing. défini) \\ \hline
Je préfère les pommes \textbf{aux} poires. & \esp{Prefiero las manzanas \textbf{a} las peras.} & \esp{a + las} : pas de contraction au féminin ni au pluriel \\ \hline
Je préfère lire \textbf{plutôt qu}'écrire. & \esp{Prefiero leer \textbf{a} escribir.} & Entre deux infinitifs : \esp{a} invariable \\ \hline
Je préfère \textbf{ceci} \textbf{à} \textbf{cela}. & \esp{Prefiero \textbf{esto} \textbf{a} \textbf{eso}.} & Pronoms démonstratifs neutres \\ \hline
\end{tabularx}
\end{center}

\courssub{La construction \esp{gustar más} : accord et ordre des mots}

La structure \cle{gustar más} (aimer davantage, préférer) fonctionne comme \esp{gustar} : le verbe s'accorde avec \textbf{ce qui est aimé} (sujet grammatical), qui se place en général \textbf{après} le verbe. La personne qui éprouve le goût est exprimée par un \textbf{pronom d'objet indirect atone} (\esp{me, te, le, nos, os, les}) placé devant le verbe ; elle peut être précisée ou redoublée par un complément introduit par \esp{a} (\esp{a mí, a ti, a Blanche...}).

\begin{propriete}[Fonctionnement de \textit{gustar más}]
\begin{enumerate}
    \item \textbf{Identifiez le sujet réel} (la chose ou l'activité préférée) : il détermine la conjugaison (\esp{gusta} au singulier, \esp{gustan} au pluriel ; singulier obligatoire si le sujet est un infinitif).
    \item \textbf{Placez le pronom COI} (\esp{me, te, le, nos, os, les}) \textbf{devant} le verbe.
    \item \textbf{Ajoutez \esp{más}} juste après le verbe.
    \item \textbf{Optionnel :} précisez la personne avec \esp{a + pronom tonique ou nom} (\esp{a mí, a ti, a Blanche...}) pour insister ou lever l'ambiguïté (obligatoire en clarté à la 3e personne).
\end{enumerate}
\end{propriete}

\begin{exemplebox}[Variantes de \textit{gustar más}]
\esp{\textbf{Me gusta más} WhatsApp.} « Je préfère WhatsApp. » (sujet singulier : \esp{WhatsApp})

\esp{\textbf{Me gustan más} las redes sociales.} « Je préfère les réseaux sociaux. » (sujet pluriel : \esp{las redes sociales})

\esp{\textbf{Me gusta más} \textbf{estudiar} en línea.} « J'aime mieux étudier en ligne. » (sujet = infinitif $\rightarrow$ verbe au singulier)

\esp{\textbf{A Blanche le gusta más} Telegram.} « Blanche, elle, préfère Telegram. » (précision à la 3e personne)

\esp{\textbf{¿Te gusta más} navegar \textbf{o} chatear?} « Préfères-tu naviguer ou bavarder en ligne ? »
\end{exemplebox}

\courssub{Renforcement : \esp{antes que} et \esp{más que} (niveau Baccalauréat)}

Pour exprimer une préférence marquée dans un choix exclusif (« plutôt que »), on utilise deux structures avancées très prisées dans les sujets de thème et de version du Baccalauréat.

\begin{aretenir}[Structures renforcées de préférence]
\begin{itemize}
    \item \textbf{\esp{Preferir + infinitif + antes que + infinitif}} (même sujet pour les deux actions) :\\
    \esp{Prefiero leer \textbf{antes que} ver videos.} « Je préfère lire plutôt que de regarder des vidéos. »
    \item \textbf{\esp{Preferir + infinitif + antes de que + subjonctif}} (changement de sujet) :\\
    \esp{Prefiero hacerlo yo \textbf{antes de que} lo \textbf{hagas} tú.} « Je préfère le faire moi-même plutôt que tu ne le fasses. »
    \item \textbf{\esp{Gustar más + que} / \esp{preferir ... más que}} (comparatif de supériorité entre deux noms) :\\
    \esp{Me gusta más la privacidad \textbf{que} la popularidad.} « Je préfère la confidentialité à la popularité. »
\end{itemize}
\emph{Retenez :} \esp{antes que} + infinitif (même sujet) ; \esp{antes de que} + subjonctif (sujets différents) ; \esp{más ... que} entre deux noms.
\end{aretenir}

\courssub{Méthode : traduire « j'aime mieux / je préfère »}

\begin{methode}[Choisir entre \textit{preferir} et \textit{gustar más}]
Pour traduire une phrase française exprimant la préférence vers l'espagnol :
\begin{enumerate}
    \item \textbf{Analysez la nuance :} \esp{preferir} est plus direct et décisionnel (« je choisis X ») ; \esp{gustar más} est plus subjectif et affectif (« X me plaît davantage »). Les deux traduisent « préférer / aimer mieux ».
    \item \textbf{Si vous choisissez \esp{preferir} :} conjuguez selon le \textbf{sujet qui préfère} (\esp{yo} $\rightarrow$ \esp{prefiero}, avec diphtongaison), puis complétez : élément préféré + \esp{a} + élément rejeté (ou infinitif direct, sans préposition).
    \item \textbf{Si vous choisissez \esp{gustar más} :} identifiez \textbf{ce qui est préféré} (sujet grammatical), accordez le verbe (\esp{gusta} / \esp{gustan}), placez le pronom COI devant, puis \esp{más} + sujet. Pour comparer deux noms, ajoutez \esp{que} + second terme.
    \item \textbf{Vérifiez l'accord :} \esp{Me gustan más los libros} (pluriel) mais \esp{Me gusta más leer} (infinitif = singulier).
    \item \textbf{Pour « plutôt que » :} \esp{antes que} + infinitif si le sujet ne change pas ; \esp{antes de que} + subjonctif s'il change.
\end{enumerate}
\end{methode}

\courssub{Erreurs fréquentes des francophones : pièges à éviter}

\begin{attention}[Pièges classiques : \textit{preferir} et \textit{gustar más}]
\begin{itemize}
    \item \textbf{\esp{*Prefero} au lieu de \esp{Prefiero}} : oubli de la diphtongaison \cle{e $\rightarrow$ ie} aux trois personnes du singulier et à la 3e du pluriel. \emph{Règle : accent tonique sur le radical = diphtongue.}
    \item \textbf{\esp{*Me gusta más de WhatsApp}} : jamais de \esp{de} après \esp{gustar}. \esp{Gustar} construit la chose aimée comme sujet ; la préposition \esp{a} n'introduit que la \textbf{personne} qui éprouve le goût (\esp{A mí me gusta...}).
    \item \textbf{\esp{*Prefiero WhatsApp que Telegram}} : la préposition de comparaison avec \esp{preferir} est \cle{a}, pas \esp{que}. \esp{Que} s'emploie avec le comparatif \esp{más/menos} : \esp{Me gusta más WhatsApp \textbf{que} Telegram}.
    \item \textbf{Accord de \esp{gustar} :} \esp{*Me gusta más los videos} $\rightarrow$ \esp{Me \textbf{gustan} más los videos} (sujet pluriel).
    \item \textbf{\esp{Preferir} + \esp{de} :} \esp{*Prefiero de leer} $\rightarrow$ \esp{Prefiero leer} (infinitif direct, sans préposition, contrairement au français « je préfère de... » qui n'existe d'ailleurs pas non plus).
    \item \textbf{Faux ami :} \esp{preferir} ne signifie jamais « préférer » au sens de « privilégier officiellement » dans tous les contextes administratifs ; au lycée, retenez simplement qu'il traduit « préférer / aimer mieux » et qu'il exige la diphtongaison.
\end{itemize}
\end{attention}

\courssub{Exercice résolu : traduction guidée (thème)}

Appliquons la méthode aux quatre phrases suivantes, extraites de contextes numériques camerounais.

\begin{exoresolu}[Traduction : préférences et habitudes numériques]
\textbf{Énoncé :} Traduisez en espagnol.
\begin{enumerate}
    \item Je préfère utiliser mon téléphone intelligent pour naviguer sur Internet.
    \item Nous aimons davantage les réseaux sociaux que la télévision.
    \item Ma sœur préfère WhatsApp à Facebook Messenger.
    \item Ils préfèrent télécharger des vidéos plutôt que de les regarder en streaming.
\end{enumerate}

\tcblower

\textbf{Solution détaillée :}

\begin{enumerate}
    \item \esp{\textbf{Prefiero usar} mi teléfono inteligente \textbf{para} navegar \textbf{por} Internet.}
    \begin{itemize}
        \item \textit{Choix :} \esp{preferir} (action volontaire). \textit{Sujet :} \esp{yo} $\rightarrow$ \esp{prefiero} (diphtongaison).
        \item \textit{Infinitif :} \esp{usar}, directement après \esp{preferir}, sans préposition.
        \item \textit{Finalité :} \esp{para} + infinitif (« pour »).
        \item \textit{Vocabulaire :} \esp{teléfono inteligente} ; \esp{navegar por Internet} (la préposition \esp{por} est la plus courante avec \esp{navegar}).
    \end{itemize}

    \item \esp{\textbf{Nos gustan más} las redes sociales \textbf{que} la televisión.}
    \begin{itemize}
        \item \textit{Choix :} \esp{gustar más} + comparatif \esp{que}.
        \item \textit{Sujet grammatical :} \esp{las redes sociales} (fém. plur.) $\rightarrow$ \esp{gustan} (pluriel).
        \item \textit{Pronom COI 1re pers. plur. :} \esp{nos}.
        \item \textit{Comparatif :} \esp{más ... que}.
        \item \textit{Alternative avec \esp{preferir} :} \esp{Preferimos las redes sociales a la televisión.} (sans diphtongaison à \esp{nosotros}).
    \end{itemize}

    \item \esp{\textbf{A mi hermana le gusta más} WhatsApp \textbf{que} Facebook Messenger.} \quad \textbf{OU} \quad \esp{\textbf{Mi hermana prefiere} WhatsApp \textbf{a} Facebook Messenger.}
    \begin{itemize}
        \item \textit{Solution 1 (\esp{gustar}) :} sujet = \esp{WhatsApp} (sing.) $\rightarrow$ \esp{gusta} ; COI = \esp{le} (3e pers. sing.) ; la précision \esp{A mi hermana} lève l'ambiguïté du pronom \esp{le} ; comparatif \esp{que}.
        \item \textit{Solution 2 (\esp{preferir}) :} sujet = \esp{mi hermana} $\rightarrow$ \esp{prefiere} (diphtongaison) ; structure \esp{preferir X a Y}.
    \end{itemize}

    \item \esp{\textbf{Prefieren descargar} videos \textbf{antes que} verlos en \textit{streaming}.}
    \begin{itemize}
        \item \textit{Sujet :} \esp{ellos} $\rightarrow$ \esp{prefieren} (diphtongaison).
        \item \textit{Structure « plutôt que » :} \esp{antes que} + infinitif (même sujet pour les deux actions).
        \item \textit{Objet :} \esp{descargar videos} ; \esp{verlos} = \esp{ver} + \esp{los} (« les regarder »), pronom soudé à l'infinitif.
        \item \textit{Vocabulaire :} \esp{en streaming} (anglicisme admis) ou \esp{en línea}.
    \end{itemize}
\end{enumerate}
\end{exoresolu}

\courssub{Entraînement : exprimez vos préférences numériques}

\begin{exercice}[Production écrite : Mes choix numériques]
Complétez les phrases suivantes en conjuguant le verbe entre parenthèses et en choisissant la structure adaptée (\esp{preferir ... a} / \esp{gustar más ... que} / \esp{antes que}). Traduisez ensuite chaque phrase en français.
\begin{enumerate}
    \item Yo \_\_\_\_\_\_\_\_\_\_ (preferir) leer libros digitales \_\_\_\_\_\_\_\_\_\_ libros en papel.
    \item A mis amigos \_\_\_\_\_\_\_\_\_\_ (gustar más) TikTok \_\_\_\_\_\_\_\_\_\_ Instagram.
    \item Nosotros \_\_\_\_\_\_\_\_\_\_ (preferir) estudiar con la computadora \_\_\_\_\_\_\_\_\_\_ con el móvil.
    \item \_\_\_\_\_\_\_\_\_\_ (gustar más) a ti la privacidad \_\_\_\_\_\_\_\_\_\_ la popularidad en las redes ?
    \item Ellos \_\_\_\_\_\_\_\_\_\_ (preferir) la conexión wifi \_\_\_\_\_\_\_\_\_\_ los datos móviles.
\end{enumerate}
\end{exercice}

\begin{corrige}[Exercice : Mes choix numériques]
\begin{enumerate}
    \item \esp{Yo \textbf{prefiero} leer libros digitales \textbf{a} libros en papel.} « Je préfère lire des livres numériques plutôt que des livres papier. » (diphtongaison \esp{prefiero} ; \esp{a} pour la comparaison).
    \item \esp{A mis amigos \textbf{les gusta más} TikTok \textbf{que} Instagram.} « Mes amis préfèrent TikTok à Instagram. » (sujet \esp{TikTok} singulier $\rightarrow$ \esp{gusta} ; COI \esp{les} ; comparatif \esp{que}).
    \item \esp{Nosotros \textbf{preferimos} estudiar con la computadora \textbf{antes que} con el móvil.} « Nous préférons étudier avec l'ordinateur plutôt qu'avec le portable. » (\esp{preferimos} sans diphtongaison ; \esp{antes que} « plutôt que »).
    \item \esp{\textbf{¿Te gusta más} la privacidad \textbf{que} la popularidad en las redes?} « Préfères-tu la vie privée à la popularité sur les réseaux ? » (interrogative : \esp{¿Te gusta más...?} ; sujet \esp{la privacidad} singulier ; \esp{que} comparatif).
    \item \esp{Ellos \textbf{prefieren} la conexión wifi \textbf{a} los datos móviles.} « Ils préfèrent la connexion wifi aux données mobiles. » (diphtongaison \esp{prefieren} ; \esp{a} pour la comparaison).
\end{enumerate}
\end{corrige}

\begin{savaistu}[Le \textit{voseo} et la préférence en Amérique latine]
Dans une grande partie de l'Amérique latine (Argentine, Uruguay, Paraguay, Amérique centrale), le pronom \cle{vos} remplace \cle{tú} pour la deuxième personne du singulier. La conjugaison de \esp{preferir} change alors : \esp{vos preferís} (accent sur la désinence \emph{-ís}, donc \textbf{pas de diphtongaison}, radical \emph{prefer-}). On entendra donc : \esp{¿Qué preferís vos?} au lieu de \esp{¿Qué prefieres tú?}. Pour \esp{gustar}, le pronom COI reste \esp{te} : \esp{¿Te gusta más...?}. Au Cameroun, l'espagnol enseigné suit la norme péninsulaire (\esp{tú}), mais la connaissance du \esp{voseo} est un atout culturel majeur pour comprendre la diversité du monde hispanophone, y compris en Guinée équatoriale, seul pays hispanophone d'Afrique subsaharienne, voisin immédiat du Cameroun.
\end{savaistu}

\coursec{Grammaire 2 : l'expression de l'habitude}

Après avoir étudié les mécanismes de la préférence (\emph{preferir, gustar más}), nous abordons maintenant un pilier de la communication quotidienne et argumentative : l'expression de l'habitude. Que ce soit pour décrire votre routine numérique (\esp{« Suelo revisar el móvil al despertar »} « J'ai l'habitude de consulter mon portable au réveil ») ou pour contraster le passé et le présent dans un commentaire de texte (\esp{« Antes solía leer periódicos; ahora navego en la red »} « Avant, je lisais des journaux ; maintenant je navigue sur le net »), la maîtrise de \cle{soler}, \cle{acostumbrar}, \cle{tener la costumbre de} et de l'\cle{imparfait} est indispensable pour réussir l'épreuve de thème, de version et d'expression écrite au Baccalauréat.

\courssub{Observation dans le texte support}

Reprenons la phrase extraite de notre texte \emph{Internet y las NTIC} :
\begin{textebox}[Extrait du texte support]
Los expertos \textbf{acostumbran a aconsejar} prudencia en el uso de las redes sociales.
\end{textebox}
\noindent On identifie la structure \textbf{acostumbran a + infinitif} (3\textsuperscript{e} personne du pluriel de l'indicatif présent). Elle exprime une habitude générale, une pratique constante des experts. Remarquez la préposition \cle{a} qui lie le verbe conjugué à l'infinitif. C'est l'une des trois grandes façons de dire « avoir l'habitude de » en espagnol, aux côtés de \cle{soler} et de \cle{tener la costumbre de}.

\courssub{L'habitude présente : trois structures majeures}

\courssub{\texorpdfstring{\cle{soler} + infinitif}{soler + infinitif} : le marqueur par excellence de l'habitude}

\begin{propriete}[Conjugaison et emploi de \texorpdfstring{\emph{soler}}{soler} au présent]
\emph{Soler} est un verbe \textbf{défautif} : il ne s'utilise pratiquement qu'aux modes personnels (indicatif, subjonctif) et presque toujours suivi d'un \textbf{infinitif}. C'est un verbe à \textbf{diphtongaison} \cle{o $\rightarrow$ ue} aux trois personnes du singulier et à la 3\textsuperscript{e} personne du pluriel.
\begin{center}
\begin{tabular}{|l|l|l|}
\hline
\rowcolor{popblueL}
\textbf{Personne} & \textbf{Forme} & \textbf{Traduction} \\
\hline
Yo & \textbf{suelo} & j'ai l'habitude de / je tends à \\
Tú & \textbf{sueles} & tu as l'habitude de \\
Él / Ella / Usted & \textbf{suele} & il/elle a l'habitude de \\
\hline
Nosotros / -as & \textbf{solemos} & nous avons l'habitude de \\
Vosotros / -as & \textbf{soléis} & vous avez l'habitude de \\
Ellos / Ellas / Ustedes & \textbf{suelen} & ils/elles ont l'habitude de \\
\hline
\end{tabular}
\end{center}
\emph{Emploi} : \emph{Soler} + infinitif décrit une action répétée, une routine, une tendance générale dans le présent. Il ne se conjugue \textbf{jamais} au passé composé pour exprimer l'habitude passée (on utilisera l'imparfait \emph{solía}).
\end{propriete}

\begin{exemplebox}[Exemples avec \emph{soler} au présent]
\begin{itemize}
\item \esp{Los jóvenes \textbf{suelen} \textbf{pasar} horas en TikTok.} « Les jeunes ont l'habitude de passer des heures sur TikTok. »
\item \esp{\textbf{Suelo} \textbf{descargar} música los fines de semana.} « J'ai l'habitude de télécharger de la musique le week-end. »
\item \esp{En Yaoundé, \textbf{solemos} \textbf{navegar} con datos móviles.} « À Yaoundé, nous avons l'habitude de naviguer avec les données mobiles. »
\item \esp{¿\textbf{Soléis} \textbf{usar} contraseñas seguras?} « Avez-vous l'habitude d'utiliser des mots de passe sécurisés ? »
\end{itemize}
\end{exemplebox}

\courssub{\texorpdfstring{\emph{acostumbrar a} + infinitif}{acostumbrar a + infinitif} : registre plus formel}

\begin{propriete}[Acostumbrar a + infinitif]
\emph{Acostumbrar} (verbe régulier en \emph{-ar}) s'emploie avec la préposition \textbf{a} devant l'infinitif. Il est plus formel que \emph{soler} et fréquent à l'écrit (presse, essais, discours d'experts).
\begin{center}
\begin{tabularx}{\linewidth}{|X|X|X|}
\hline
\rowcolor{popblueL}
\textbf{Espagnol} & \textbf{Français} & \textbf{Remarque} \\
\hline
Los expertos \textbf{acostumbran a} \textbf{publicar} estudios anuales. & Les experts ont l'habitude de publier des études annuelles. & Habitude générale, fait établi. \\
\textbf{Acostumbro a} \textbf{estudiar} en la biblioteca. & J'ai l'habitude d'étudier à la bibliothèque. & Routine personnelle. \\
\hline
\end{tabularx}
\end{center}
\end{propriete}

\courssub{\texorpdfstring{\emph{acostumbrarse a} + infinitif / nom}{acostumbrarse a + infinitif / nom} : le processus d'adaptation}

\begin{attention}[Piège majeur : \texorpdfstring{\emph{acostumbrar}}{acostumbrar} vs \texorpdfstring{\emph{acostumbrarse}}{acostumbrarse} vs \texorpdfstring{\emph{acostarse}}{acostarse}]
Ne confondez surtout pas ces trois verbes, source d'erreurs graves en version et en thème :
\begin{itemize}
\item \textbf{Acostumbrar a + inf.} = \emph{Avoir l'habitude de} (état permanent). \esp{Acostumbro a leer el periódico.} (J'ai l'habitude de lire le journal.)
\item \textbf{Acostumbrarse a + inf./nom} = \emph{S'habituer à} (processus, changement d'état). \esp{Me \textbf{he acostumbrado a} \textbf{trabajar} desde casa.} (Je me suis habitué à travailler depuis chez moi.) $\rightarrow$ \cle{Prétérit / Passé composé} pour marquer l'aboutissement.
\item \textbf{Acostarse} (o $\rightarrow$ ue, pronominal) = \emph{Se coucher}. \esp{Me \textbf{acuesto} temprano.} (Je me couche tôt.)
\end{itemize}
\emph{Astuce mnémotechnique} : \cle{Acostumbr\underline{a}r} = H\underline{a}bitude (les deux ont un \emph{a}) ; \cle{Acost\underline{ar}se} = Se cou\underline{r}cher (les deux ont un \emph{r} fort).
\end{attention}

\courssub{\texorpdfstring{\emph{tener la costumbre de} + infinitif}{tener la costumbre de + infinitif} : la locution nominale}

\begin{propriete}[Tener la costumbre de + infinitif]
Structure très transparente pour les francophones (\emph{avoir l'habitude de}). Le verbe \emph{tener} se conjugue normalement (diphtongaison \emph{ie} : \emph{tengo, tienes, tiene, tenemos, tenéis, tienen}). L'article \emph{la} est obligatoire.
\end{propriete}

\begin{exemplebox}[Exemples]
\begin{itemize}
\item \esp{Tengo la costumbre de \textbf{revisar} mi correo cada mañana.} « J'ai l'habitude de vérifier mon courriel chaque matin. »
\item \esp{Mi abuelo \textbf{tenía} la costumbre de \textbf{leer} el periódico impreso.} « Mon grand-père avait l'habitude de lire le journal imprimé. » (Imparfait pour habitude passée).
\end{itemize}
\end{exemplebox}


\courssub{Tableau récapitulatif : \texorpdfstring{\og Avoir l'habitude de \fg{} en espagnol}{Avoir l'habitude de en español}}

\begin{center}
\begin{tabularx}{\linewidth}{|X|X|X|}
\hline
\rowcolor{popblueL}
\textbf{Structure} & \textbf{Niveau / Registre} & \textbf{Nuance} \\
\hline
\textbf{Soler} + inf. & Courant, oral et écrit & Habitude générale, tendance ; \emph{le plus idiomatique}. \\
\hline
\textbf{Acostumbrar a} + inf. & Soutenu, écrit, formel & Pratique constante, coutume établie (souvent 3\textsuperscript{e} pers.). \\
\hline
\textbf{Tener la costumbre de} + inf. & Standard, clair & Insiste sur la possession d'une habitude ; facile pour francophones. \\
\hline
\end{tabularx}
\end{center}

\courssub{Adverbes de fréquence associés}
Pour renforcer l'idée d'habitude, on associe souvent ces verbes à des adverbes :
\begin{center}
\begin{tabularx}{\linewidth}{|X|X|}
\hline
\rowcolor{poporangeL}
\textbf{Espagnol} & \textbf{Français} \\
\hline
\textbf{Normalmente} / \textbf{Por lo general} & \cle{D'habitude}, en général \\
\textbf{Habitualmente} / \textbf{Con frecuencia} & Habituellement, fréquemment \\
\textbf{Casi siempre} / \textbf{Siempre} & Presque toujours / Toujours \\
\textbf{A menudo} / \textbf{Muchas veces} & Souvent \\
\hline
\end{tabularx}
\end{center}
\esp{Normalmente, \textbf{suelo} \textbf{conectar} el Wi-Fi al llegar.} « D'habitude, je me connecte au Wi-Fi en arrivant. »

\courssub{L'habitude dans le passé : l'imparfait de l'indicatif}

En espagnol, comme en français, l'\textbf{imparfait de l'indicatif} (\emph{pretérito imperfecto}) est le temps roi pour exprimer une habitude passée, une action répétée dans le passé, sans limite temporelle précise.

\begin{propriete}[Imparfait d'habitude : formation et valeurs]
\textbf{Formation} (régulière pour tous les verbes, même \emph{ir} / \emph{ser} / \emph{ver} qui ont leur propre radical) :
\begin{itemize}
\item \emph{-AR} : \textbf{-aba, -abas, -aba, -ábamos, -abais, -aban} (\esp{hablaba})
\item \emph{-ER / -IR} : \textbf{-ía, -ías, -ía, -íamos, -íais, -ían} (\esp{comía, vivía})
\end{itemize}
\textbf{Valeurs principales (habitude passée)} :
\begin{enumerate}
\item Action répétée, routine passée : \esp{Cada día, mi abuelo \textbf{leía} el periódico.} « Chaque jour, mon grand-père lisait le journal. »
\item Description d'état passé : \esp{Antes, \textbf{teníamos} la costumbre de \textbf{escribir} cartas.} « Avant, nous avions l'habitude d'écrire des lettres. »
\item Contraste passé / présent (très fréquent en commentaire de texte) : \esp{\textbf{Antes navegaba} poco; \textbf{ahora paso} horas en línea.} « Avant je naviguais peu ; maintenant je passe des heures en ligne. »
\end{enumerate}
\end{propriete}

\courssub{\texorpdfstring{\emph{Soler} à l'imparfait : \emph{solía}}{Soler à l'imparfait : solía}}

C'est la forme spécifique pour dire « \emph{j'avais l'habitude de} / \emph{je faisais souvent} » avec le verbe \emph{soler}. C'est un \textbf{imparfait régulier en \emph{-ía}} (radical \emph{sol-}).

\begin{propriete}[Conjugaison de \texorpdfstring{\emph{soler}}{soler} à l'imparfait]
\begin{center}
\begin{tabular}{|l|l|l|}
\hline
\rowcolor{poppurpleL}
\textbf{Personne} & \textbf{Forme} & \textbf{Traduction} \\
\hline
Yo & \textbf{solía} & j'avais l'habitude de / je faisais souvent \\
Tú & \textbf{solías} & tu avais l'habitude de \\
Él / Ella / Usted & \textbf{solía} & il/elle avait l'habitude de \\
\hline
Nosotros / -as & \textbf{solíamos} & nous avions l'habitude de \\
Vosotros / -as & \textbf{solíais} & vous aviez l'habitude de \\
Ellos / Ellas / Ustedes & \textbf{solían} & ils/elles avaient l'habitude de \\
\hline
\end{tabular}
\end{center}
\emph{Attention} : Pas de diphtongaison à l'imparfait ! La diphtongaison \emph{o $\rightarrow$ ue} n'existe qu'au présent (et au subjonctif présent).
\end{propriete}

\begin{exemplebox}[Exemples traduits (habitude passée)]
\begin{itemize}
\item \esp{De niño, \textbf{solía} \textbf{jugar} al fútbol en la calle.} « Enfant, je jouais habituellement au football dans la rue / J'avais l'habitude de jouer... »
\item \esp{\textbf{Solíamos} \textbf{descargar} archivos lentos con el módem 56k.} « Nous avions l'habitude de télécharger des fichiers lents avec le modem 56k. »
\item \esp{Antes, la gente \textbf{tenía} la costumbre de \textbf{hablar} por teléfono fijo.} « Avant, les gens avaient l'habitude de parler au téléphone fixe. »
\item \esp{Los profesores \textbf{acostumbraban a} \textbf{dictar} las notas.} « Les professeurs avaient l'habitude de dicter les notes. »
\end{itemize}
\end{exemplebox}

\begin{popfigure}
\legende{Frise chronologique : de l'habitude passée (imparfait) à l'habitude présente (présent).}
\begin{tikzpicture}[
    node distance=1.2cm and 2cm,
    fleche/.style={->, >=Stealth, thick, popblue},
    temps/.style={rectangle, rounded corners, draw=popblue, fill=popblueL, thick, minimum width=3.5cm, minimum height=1.8cm, align=center, font=\small},
    label/.style={font=\footnotesize, text=popdark, align=center}
]
\node[temps, align=center] (passe){\textbf{HABITUDE PASSÉE}\\\emph{Imparfait de l'indicatif}\\\textcolor{poppurple}{\textbf{solía navegar}}\\navigaba / tenía la costumbre};
\node[temps, right=of passe, align=center] (present){\textbf{HABITUDE PRÉSENTE}\\\emph{Présent de l'indicatif}\\\textcolor{poporange}{\textbf{suelo navegar}}\\navego / acostumbro a navegar};
\draw[fleche] (passe.east) -- node[above, label] {\textbf{AVANT / AUTREFOIS}} (present.west);
\draw[fleche] (present.west) -- node[below, label] {\textbf{MAINTENANT / ACTUELLEMENT}} (passe.east);
\node[label, below=0.3cm of passe.south] {\textcolor{popmuted}{Cada día / Antes / Antaño}};
\node[label, below=0.3cm of present.south] {\textcolor{popmuted}{Normalmente / Ahora / Cada día}};
\node[label, left=0.5cm of passe.west, text width=3cm, align=right] {\esp{De niño, \textbf{solía} \textbf{jugar} fuera.}\\« Enfant, je jouais dehors. »};
\node[label, right=0.5cm of present.east, text width=3cm, align=left] {\esp{Hoy, los niños \textbf{suelen} \textbf{jugar} online.}\\« Aujourd'hui, les enfants jouent en ligne. »};
\end{tikzpicture}
\end{popfigure}

\begin{savaistu}[Le piège du Bac : Passé Composé vs Imparfait]
Au Baccalauréat (Thème et Version), la distinction entre \textbf{Passé Composé} (\emph{pretérito perfecto compuesto} : \esp{he navegado}) et \textbf{Imparfait} (\esp{navegaba / solía navegar}) est \textbf{systématiquement évaluée}.
\begin{itemize}
\item \textbf{Passé Composé} = Action \textbf{ponctuelle}, \textbf{accomplie}, \textbf{datée} dans un passé clos. \esp{Ayer \textbf{descargué} un archivo.} (Hier, j'ai téléchargé un fichier \emph{une fois}.)
\item \textbf{Imparfait} = Action \textbf{habituelle}, \textbf{répétée}, \textbf{non bornée} dans le passé. \esp{Antes \textbf{descargaba} archivos cada noche.} (Avant, je téléchargeais des fichiers \emph{toutes les nuits}.)
\end{itemize}
\cle{Consigne d'or} : Si le français dit « \emph{j'avais l'habitude de} / \emph{je faisais souvent} / \emph{chaque jour je...} » $\rightarrow$ \textbf{Imparfait} en espagnol (\emph{solía + inf.} ou \emph{verbe à l'imparfait}). Si le français dit « \emph{j'ai fait (une fois)} / \emph{hier j'ai...} » $\rightarrow$ \textbf{Passé Composé}.
\end{savaistu}

\courssub{Traduction guidée (Thème) : Exercice résolu}

\begin{exoresolu}[Thème : Exprimer l'habitude (présent et passé)]
\textbf{Énoncé :} Traduisez en espagnol les phrases suivantes en utilisant les structures étudiées (\emph{soler, acostumbrar a, tener la costumbre de, imparfait}).

\begin{enumerate}
\item D'habitude, je télécharge des applications le week-end.
\item Les experts ont l'habitude de recommander la prudence sur les réseaux.
\item Avant, nous avions l'habitude de lire le journal papier ; maintenant nous naviguons sur Internet.
\item Je me suis habitué à travailler avec mon téléphone intelligent.
\end{enumerate}

\tcblower

\textbf{Solution détaillée et commentée :}

\begin{enumerate}
\item \textbf{D'habitude, je télécharge des applications le week-end.}
    \begin{itemize}
    \item \emph{D'habitude} $\rightarrow$ \esp{Normalmente} / \esp{Por lo general}.
    \item \emph{J'ai l'habitude de télécharger} (présent) $\rightarrow$ \esp{Suelo descargar} (le plus naturel) ou \esp{Tengo la costumbre de descargar}.
    \item \emph{Le week-end} $\rightarrow$ \esp{los fines de semana} / \esp{el fin de semana}.
    \item \textbf{Proposition :} \esp{\textbf{Normalmente, suelo descargar aplicaciones los fines de semana.}}
    \end{itemize}

\item \textbf{Les experts ont l'habitude de recommander la prudence sur les réseaux.}
    \begin{itemize}
    \item \emph{Les experts} $\rightarrow$ \esp{Los expertos}.
    \item \emph{Ont l'habitude de recommander} (registre formel, vérité générale) $\rightarrow$ \esp{Acostumbran a recomendar} (idéal ici) ou \esp{Suelen recomendar}.
    \item \emph{La prudence} $\rightarrow$ \esp{la prudencia}.
    \item \emph{Sur les réseaux} $\rightarrow$ \esp{en las redes (sociales)}.
    \item \textbf{Proposition :} \esp{\textbf{Los expertos acostumbran a recomendar la prudencia en las redes sociales.}} (Reprise exacte du texte support !)
    \end{itemize}

\item \textbf{Avant, nous avions l'habitude de lire le journal papier ; maintenant nous naviguons sur Internet.}
    \begin{itemize}
    \item \emph{Avant} $\rightarrow$ \esp{Antes} / \esp{Antiguamente}.
    \item \emph{Nous avions l'habitude de lire} (habitude passée) $\rightarrow$ \esp{Solíamos leer} (meilleur choix) ou \esp{Leíamos} (imparfait simple) ou \esp{Teníamos la costumbre de leer}.
    \item \emph{Le journal papier} $\rightarrow$ \esp{el periódico impreso} / \esp{el periódico en papel}.
    \item \emph{Maintenant} $\rightarrow$ \esp{Ahora} / \esp{Actualmente}.
    \item \emph{Nous naviguons} (présent d'habitude ou action actuelle) $\rightarrow$ \esp{Navegamos} / \esp{Solemos navegar}.
    \item \emph{Sur Internet} $\rightarrow$ \esp{en Internet} / \esp{por Internet} / \esp{en la red}.
    \item \textbf{Proposition :} \esp{\textbf{Antes solíamos leer el periódico impreso; ahora navegamos en Internet.}}
    \end{itemize}

\item \textbf{Je me suis habitué à travailler avec mon téléphone intelligent.}
    \begin{itemize}
    \item \emph{Je me suis habitué} (processus achevé, résultat présent) $\rightarrow$ \emph{Me he acostumbrado} (Passé composé / \emph{Pretérito perfecto compuesto}) \textbf{NOT} \emph{Me acostumbré} (Prétérit, action passée coupée du présent, moins courant pour cet aspect).
    \item \emph{À travailler} $\rightarrow$ \esp{a trabajar} (préposition \textbf{a} obligatoire après \emph{acostumbrarse}).
    \item \emph{Avec mon téléphone intelligent} $\rightarrow$ \esp{con mi teléfono inteligente} / \esp{con mi smartphone}.
    \item \textbf{Proposition :} \esp{\textbf{Me he acostumbrado a trabajar con mi teléfono inteligente.}}
    \item \emph{Note} : Si le contexte était narratif passé pur (\emph{Il y a deux ans, je m'habituai à...}), on utiliserait le prétérit : \esp{Me acostumbré a trabajar...}.
    \end{itemize}
\end{enumerate}
\end{exoresolu}

\courssub{Entraînement : Lecture et production écrite}

Voici un texte original, dans la lignée du programme \emph{Médias, communication et culture de la paix}, pour réinvestir le lexique des NTIC et les structures d'habitude.

\begin{textebox}[La brecha digital en mi barrio]
En mi barrio de Douala, la conexión a Internet \textbf{suele} \textbf{ser} inestable. \textbf{Antes}, \textbf{solíamos} \textbf{ir} al cibercafé para revisar el correo o hacer trámites, pero \textbf{ahora} casi todos \textbf{tenemos} la costumbre de \textbf{usar} el teléfono inteligente. Los jóvenes \textbf{acostumbran a} \textbf{pasar} horas en WhatsApp o TikTok, mientras que los mayores \textbf{suelen} \textbf{preferir} las llamadas de voz. Sin embargo, \textbf{nos hemos acostumbrado a} \textbf{compartir} los datos móviles con los vecinos cuando la fibra óptica falla. \textbf{Antiguamente}, la \textbf{brecha digital} \textbf{era} un muro; hoy, \textbf{suele} \textbf{ser} un reto que superamos juntos con creatividad.
\end{textebox}

\noindent \textbf{Glossaire :}
\begin{itemize}
\item \esp{la conexión} : la connexion
\item \esp{inestable} : instable
\item \esp{el cibercafé} : le cybercafé
\item \esp{hacer trámites} : faire des démarches administratives
\item \esp{los mayores} : les aînés, les personnes âgées
\item \esp{las llamadas de voz} : les appels vocaux
\item \esp{compartir los datos móviles} : partager la connexion (hotspot)
\item \esp{la fibra óptica} : la fibre optique
\item \esp{fallar} : tomber en panne, dysfonctionner
\item \esp{la brecha digital} : la fracture numérique
\item \esp{un muro} : un mur
\item \esp{un reto} : un défi
\item \esp{superar} : surmonter
\end{itemize}

\noindent \textbf{Questions de compréhension et d'analyse grammaticale :}
\begin{enumerate}
\item Relevez \textbf{toutes} les formes verbales exprimant l'habitude (présent ou passé) dans le texte. Classez-les dans un tableau : \emph{Structure utilisée / Personne / Temps / Infinitif associé}.
\item Traduisez en français les deux phrases contenant \esp{solíamos} et \esp{nos hemos acostumbrado a}. Expliquez la différence de temps et d'aspect.
\item Réécrivez la première phrase en remplaçant \esp{suele ser} par une structure avec \esp{acostumbrar} (attention à la structure !) et par \esp{tener la costumbre de}. Laquelle sonne le plus naturel ? Pourquoi ?
\item \textbf{Production orale / Écrite :} Dans votre établissement (lycée de Yaoundé, Douala, Bafoussam...), comment \emph{soléis} / \emph{acostumbráis} / \emph{tenéis la costumbre} d'utiliser Internet ? Y a-t-il une \emph{brecha digital} entre les élèves ? Répondez en 80-100 mots en utilisant au moins 3 structures d'habitude différentes.
\end{enumerate}

\courssub{Exercice d'application : Transformation et choix}

\begin{exercice}[Mise en forme : Habitude présente vs passée]
Mettez le verbe entre parenthèses à la forme correcte (Présent d'habitude ou Imparfait d'habitude) et traduisez la phrase.

\begin{enumerate}
\item \esp{Antes, mi padre \_\_\_\_\_\_\_\_\_\_ (leer) el periódico cada mañana, pero ahora \_\_\_\_\_\_\_\_\_\_ (leer) las noticias en el móvil.}
\item \esp{Los estudiantes \_\_\_\_\_\_\_\_\_\_ (soler / presente) \_\_\_\_\_\_\_\_\_\_ (usar) Wikipedia para las tareas.}
\item \esp{Cuando era niño, yo \_\_\_\_\_\_\_\_\_\_ (soler / imperfecto) \_\_\_\_\_\_\_\_\_\_ (jugar) en la calle, no con la consola.}
\item \esp{En mi casa, siempre \_\_\_\_\_\_\_\_\_\_ (tener) la costumbre de \_\_\_\_\_\_\_\_\_\_ (cenar) juntos.}
\item \esp{Los profesores \_\_\_\_\_\_\_\_\_\_ (acostumbrar / presente) \_\_\_\_\_\_\_\_\_\_ (dar) deberes los viernes.}
\item \esp{¿Tú \_\_\_\_\_\_\_\_\_\_ (soler / presente) \_\_\_\_\_\_\_\_\_\_ (ver) series en streaming los fines de semana?}
\end{enumerate}
\end{exercice}

\begin{corrige}[Exercice d'application]
\begin{enumerate}
\item \esp{Antes, mi padre \textbf{leía} el periódico cada mañana, pero ahora \textbf{lee} / \textbf{suele leer} las noticias en el móvil.} (Imparfait d'habitude / Présent d'habitude). « Avant, mon père lisait le journal chaque matin, mais maintenant il lit / a l'habitude de lire les infos sur le mobile. »
\item \esp{Los estudiantes \textbf{suelen usar} Wikipedia para las tareas.} (Présent de \emph{soler} + inf.). « Les étudiants ont l'habitude d'utiliser Wikipedia pour les devoirs. »
\item \esp{Cuando era niño, yo \textbf{solía jugar} en la calle, no con la consola.} (Imparfait de \emph{soler} + inf.). « Quand j'étais enfant, j'avais l'habitude de jouer dans la rue, pas avec la console. »
\item \esp{En mi casa, siempre \textbf{teníamos} la costumbre de \textbf{cenar} juntos.} (Imparfait de \emph{tener} + locution). « Dans ma maison, nous avions toujours l'habitude de dîner ensemble. »
\item \esp{Los profesores \textbf{acostumbran a dar} deberes los viernes.} (Présent de \emph{acostumbrar a} + inf.). « Les professeurs ont l'habitude de donner des devoirs les vendredis. »
\item \esp{¿Tú \textbf{sueles ver} series en streaming los fines de semana?} (Présent de \emph{soler} + inf.). « As-tu l'habitude de regarder des séries en streaming le week-end ? »
\end{enumerate}
\end{corrige}

\coursec{Méthode du commentaire et production guidée}

Nous savons maintenant exprimer la préférence (\esp{preferir}, \esp{me gusta más...}) et l'habitude (\esp{soler + infinitivo}, \esp{acostumbrar a}, l'imparfait). Il reste à assembler tous ces outils dans l'exercice roi de la Terminale A : le \cle{commentaire de texte}, puis dans une production écrite personnelle. C'est précisément ce qui est demandé à l'épreuve d'espagnol du Baccalauréat.

\courssub{La méthode du commentaire de texte}

\begin{definition}[Le commentaire de texte]
Le \cle{commentaire de texte} (\esp{comentario de texto}) est un exercice d'analyse et d'argumentation : il ne s'agit pas de résumer le texte, mais d'expliquer \emph{comment} il est construit, \emph{ce que} l'auteur défend, et de donner son \emph{avis personnel argumenté}. Au Baccalauréat, il est rédigé en espagnol et évalué sur la compréhension, la langue et la qualité de l'argumentation.
\end{definition}

\begin{methode}[Comment réussir un commentaire au Bac]
\begin{enumerate}
\item \textbf{Lire deux fois le texte} : une première fois pour le sens général, une deuxième fois en soulignant les mots-clés, les connecteurs et les idées principales.
\item \textbf{Rédiger l'introduction} : présenter le texte (thème, type, idée principale de l'auteur) et formuler une \cle{problématique} (\esp{¿Internet une o separa a las personas?}).
\item \textbf{Rédiger le développement} : 2 ou 3 paragraphes. Chaque paragraphe = une idée + une \textbf{citation du texte} + une analyse + un exemple. Analyser, ne jamais paraphraser.
\item \textbf{Donner son avis argumenté} : \esp{En mi opinión...}, \esp{Estoy de acuerdo con el autor porque...}, en s'appuyant sur sa propre expérience (Cameroun, école, famille).
\item \textbf{Conclure} : bilan de l'analyse (\esp{En conclusión...}) puis \textbf{ouverture} vers une question plus large (\esp{¿Será Internet un derecho universal en el futuro?}).
\item \textbf{Relire} : accords, accents, diphtongaison, ponctuation espagnole ¿ ? ¡ !
\end{enumerate}
\end{methode}

\begin{popfigure}
\begin{tikzpicture}[
  bloc/.style={rectangle, rounded corners=4pt, draw=popblue, thick, fill=popblueL, text width=3.4cm, align=center, minimum height=1.6cm, font=\small},
  fleche/.style={-{Stealth[length=3mm]}, thick, poporange}
]
\node[bloc, align=center] (intro) at (0,0){\textbf{Introducción}\\ presentación +\\ problemática};
\node[bloc, fill=popgreenL, draw=popgreen, align=center] (des) at (5,0){\textbf{Desarrollo}\\ 2-3 párrafos:\\ argumentos + ejemplos};
\node[bloc, fill=poppinkL, draw=poppink, align=center] (conc) at (10,0){\textbf{Conclusión}\\ balance +\\ apertura};
\draw[fleche] (intro) -- (des);
\draw[fleche] (des) -- (conc);
\end{tikzpicture}
\legende{La structure en trois blocs du commentaire de texte}
\end{popfigure}

\courssub{Un commentaire modèle : début rédigé}

Observons comment un élève de Terminale peut commencer son commentaire sur le texte \esp{Internet: ¿puente o brecha?} étudié dans cette leçon.

\begin{textebox}[Début de commentaire modèle]
Este texto \textbf{trata sobre} Internet y las NTIC. El autor defiende la idea de que Internet es un puente hacia el saber, pero denuncia la brecha digital que separa a los jóvenes conectados de los que no tienen acceso a la red. ¿Es Internet una oportunidad para todos o un nuevo motivo de desigualdad?

En primer lugar, el autor muestra las ventajas de Internet: gracias a un teléfono inteligente, un estudiante de Yaoundé puede descargar libros y navegar por bibliotecas del mundo entero. Sin embargo, denuncia que muchas familias no pueden pagar una conexión. Además, subraya que las redes sociales pueden unir a los amigos, pero también aislar a las personas. Por eso, propone que la escuela enseñe un uso responsable de la red. En mi opinión, el autor tiene razón: en mi barrio, los chicos que suelen navegar con criterio aprenden más que los que solo miran vídeos.
\end{textebox}

\textbf{Glossaire} : \esp{el saber} = le savoir ; \esp{la brecha digital} = la fracture numérique ; \esp{subrayar} = souligner ; \esp{aislar} = isoler ; \esp{con criterio} = avec discernement ; \esp{el barrio} = le quartier.

\begin{exemplebox}[La conclusion : exemple traduit]
\esp{En conclusión, Internet es un puente hacia el saber si se usa con responsabilidad.}

« En conclusion, Internet est un pont vers le savoir s'il est utilisé avec responsabilité. »

Remarquez la tournure impersonnelle \esp{si se usa} (« si l'on utilise ») : l'espagnol emploie \esp{se} + verbe à la 3\up{e} personne là où le français utilise « on ».
\end{exemplebox}

\courssub{Les connecteurs à réutiliser}

Les \cle{connecteurs} (\esp{conectores}) sont la charpente du commentaire : sans eux, les idées s'entrechoquent. Voici le tableau à connaître par cœur.

\begin{center}
\begin{tabularx}{\linewidth}{|X|X|X|}
\hline
\rowcolor{popblueL} Español & Français & Fonction \\
\hline
\esp{en primer lugar} & en premier lieu & introduire le 1\up{er} argument \\
\hline
\esp{además} & de plus & ajouter un argument \\
\hline
\esp{sin embargo} & cependant & opposer, nuancer \\
\hline
\esp{por eso} & c'est pourquoi & conséquence \\
\hline
\esp{en mi opinión} & à mon avis & avis personnel \\
\hline
\esp{en conclusión} & en conclusion & bilan final \\
\hline
\end{tabularx}
\end{center}

\begin{attention}[Pièges des francophones]
\begin{itemize}
\item \esp{Sin embargo} seul signifie « cependant » : ne jamais écrire \esp{pero sin embargo} (pléonasme).
\item \esp{Por eso} = « c'est pourquoi » ; ne pas confondre avec \esp{porque} = « parce que ».
\item Ne traduisez pas « en fait » par \esp{en facto} : dites \esp{de hecho} ou \esp{en realidad}.
\end{itemize}
\end{attention}

\begin{exoresolu}[Transformer un résumé en analyse]
Énoncé : voici une phrase de paraphrase maladroite : \esp{El texto dice que Internet es bueno y también malo.} Réécrivez-la en véritable phrase d'analyse avec citation et connecteur.

\tcblower
Solution détaillée : il faut (1) un connecteur d'opposition, (2) une citation précise, (3) une interprétation.

\esp{Sin embargo, cuando el autor afirma que Internet es «un puente y una brecha a la vez», muestra que la red une a los que tienen acceso, pero excluye a los más pobres.}

« Cependant, lorsque l'auteur affirme qu'Internet est "à la fois un pont et une fracture", il montre que le réseau relie ceux qui y ont accès, mais exclut les plus pauvres. »

On est passé du résumé (« le texte dit que ») à l'analyse (« muestra que » + interprétation).
\end{exoresolu}

\courssub{Production guidée : « Mis hábitos en Internet »}

\begin{experience}[Rédiger un paragraphe personnel]
Rédigez un paragraphe de 6 à 8 lignes intitulé \esp{Mis hábitos en Internet}. Consignes obligatoires :
\begin{itemize}
\item au moins \textbf{2 structures de préférence} : \esp{prefiero... a...}, \esp{me gusta más...} ;
\item au moins \textbf{2 structures d'habitude} : \esp{suelo + infinitivo}, \esp{tengo la costumbre de...}, \esp{acostumbro a...} ;
\item au moins \textbf{3 mots du lexique NTIC} : \esp{red social, contraseña, descargar, navegar, teléfono inteligente} ;
\item au moins \textbf{2 connecteurs} : \esp{además, sin embargo, por eso}.
\end{itemize}
\textbf{Modèle de départ} : \esp{Por las tardes, suelo navegar en Internet con mi teléfono inteligente. Prefiero las redes sociales a los juegos en línea, pero me gusta más descargar música que ver vídeos...}
\end{experience}

\begin{attention}[Avant de rendre votre paragraphe]
\begin{itemize}
\item \textbf{Diphtongaison} : \esp{preferir} donne \esp{prefie\textbf{ro}} (e $\to$ ie) ; \esp{soler} donne \esp{sue\textbf{lo}} (o $\to$ ue). Jamais \esp{prefero} ni \esp{suelo}.
\item \textbf{Accord de gustar} : le verbe s'accorde avec la chose aimée : \esp{me gusta\textbf{n} las redes sociales} (pluriel), mais \esp{me gusta navegar} (infinitif = singulier).
\item \textbf{Préposition} : \esp{soler} et \esp{preferir} + infinitif direct ; mais \esp{acostumbrar \textbf{a}} et \esp{tener la costumbre \textbf{de}}.
\end{itemize}
\end{attention}

\courssub{Entraînement à la traduction : Thème et Version}

\begin{exercice}[Traduire préférence et habitude]
Traduisez les phrases suivantes en espagnol (Thème) ou en français (Version).
\begin{enumerate}
\item Je préfère lire des livres que regarder la télévision.
\item Nous avons l'habitude de naviguer sur Internet le soir.
\item D'habitude, il télécharge de la musique sur son téléphone.
\item \esp{Me gusta más el teléfono inteligente que la tableta.}
\item \esp{Los jóvenes acostumbran a chatear en las redes sociales.}
\item \esp{¿Prefieres navegar por Internet o jugar en línea?}
\end{enumerate}
\end{exercice}

\begin{corrige}[Exercice traduction]
\begin{enumerate}
\item \esp{Prefiero leer libros a ver la televisión.} (\esp{preferir} + infinitif, structure \esp{prefiero... a...})
\item \esp{Solemos navegar por Internet por la tarde.} (ou \esp{Tenemos la costumbre de navegar...} ; \esp{soler} + infinitif, diphtongue \esp{solemos})
\item \esp{Suele descargar música en su teléfono.} (\esp{soler} à la 3\up{e} pers. sg. \esp{suele} ; \esp{en} pour « sur » un support numérique)
\item « J'aime mieux le téléphone intelligent que la tablette. » (\esp{gustar} accordé au singulier \esp{gusta} ; comparatif \esp{más... que...})
\item « Les jeunes ont l'habitude de tchatter sur les réseaux sociaux. » (\esp{acostumbrar a} + infinitif ; \esp{chatear} anglicisme admis)
\item « Préfères-tu naviguer sur Internet ou jouer en ligne ? » (\esp{prefieres} diphtongue ; \esp{navegar por} / \esp{jugar en} ligne)
\end{enumerate}
\end{corrige}

\begin{savaistu}[Internet en zone francophone et hispanophone africaine]
La \cle{brecha digital} (fracture numérique) est une réalité majeure au Cameroun et en Guinée équatoriale. Si les grandes villes (Yaoundé, Douala, Malabo, Bata) bénéficient de la 4G/5G et de la fibre optique, les zones rurales restent souvent déconnectées. Le coût des données (\esp{megascosto}) et l'accès à l'électricité freinent l'usage éducatif des NTIC. Cependant, des initiatives comme les \esp{telecentros} communautaires ou les programmes \esp{« Un portátil por niño »} (inspirés du projet OLPC) tentent de réduire l'écart. En Guinée équatoriale, seul pays hispanophone d'Afrique, l'espagnol est la langue de l'administration et de l'école : maîtriser les NTIC en espagnol y est un atout professionnel décisif.
\end{savaistu}

\courssub{Grille d'auto-évaluation}

Après avoir rédigé, évaluez-vous honnêtement avec cette grille (à recopier dans le cahier) :

\begin{center}
\begin{tabularx}{\linewidth}{|X|l|l|}
\hline
\rowcolor{popblueL} Critère & Sí & No \\
\hline
J'ai utilisé au moins 4 mots du lexique NTIC correctement & & \\
\hline
J'ai écrit \esp{prefiero} et \esp{suelo} avec la diphtongue & & \\
\hline
J'ai accordé \esp{gustar} (gusta / gustan) & & \\
\hline
J'ai employé au moins 2 connecteurs variés & & \\
\hline
J'ai mis les accents et les signes ¿ ¡ & & \\
\hline
Mon paragraphe a une phrase de conclusion & & \\
\hline
\end{tabularx}
\end{center}

\begin{aretenir}[L'essentiel de la section]
Un bon commentaire = introduction (présentation + problématique) + développement (arguments, citations, exemples) + conclusion (bilan + ouverture). On \emph{analyse}, on ne paraphrase pas. Les connecteurs (\esp{en primer lugar, además, sin embargo, por eso, en conclusión}) structurent le propos. Dans toute production personnelle, vigilance absolue sur la diphtongaison (\esp{prefiero, suelo}) et l'accord de \esp{gustar}. Pour traduire « avoir l'habitude de » : \esp{soler} (présent/imparfait) ou \esp{acostumbrar a} / \esp{tener la costumbre de} + infinitif.
\end{aretenir}

\newpage

\coursec{Activité d'intégration}

\coursec{Leçon E28 : Lectura y comentario : Internet y las NTIC ; traduction : expression de la préférence et de l'habitude}

\courssub{Objectifs de l'activité}

Cette activité d'intégration vise à mobiliser l'ensemble des acquis de la leçon E28 dans une situation de communication authentique et collaborative. À l'issue de cette séance, l'élève doit être capable de :

\begin{itemize}
    \item Réinvestir le lexique thématique des NTIC (\esp{red social, contraseña, navegar, teléfono inteligente, brecha digital, descargar, conexión, ancho de banda, ciberseguridad, influenciador}) dans une production écrite cohérente.
    \item Manipuler correctement les structures d'expression de la préférence (\esp{preferir, me gusta más, prefiero... a...}) et de l'habitude (\esp{soler + infinitif, acostumbrar a + infinitif, tener la costumbre de + infinitif, imparfait de l'indicatif pour l'habitude passée}).
    \item Produire un message de forum structuré (8--10 lignes) respectant les conventions du genre (titre, salutations, paragraphes, marqueurs de relation logique, signature).
    \item Présenter oralement sa production en groupe, puis interagir en posant des questions de compréhension aux autres groupes.
    \item Identifier et corriger les erreurs fréquentes : diphtongaison radicale de \esp{preferir} (\emph{prefiero, prefieres, prefiere, preferimos, preferís, prefieren}), construction de \esp{gustar} (\esp{me gusta / me gustan} + infinitif ou nom), emploi de la préposition \esp{a} après \esp{preferir} pour la comparaison, accord du participe passé dans \esp{acostumbrado}, distinction \esp{soler} (habitude présente) / \esp{solía} (habitude passée).
\end{itemize}

\courssub{Situation}

\begin{textebox}[Contexte : Forum numérique inter-lycées Cameroun -- Espagne -- Guinée équatoriale]
Dans le cadre du projet \emph{« Puentes digitales »}, un forum en ligne réunit des élèves de Terminale A du Lycée Général Leclerc (Yaoundé), de l'Instituto Cervantes (Madrid) et du Colegio Nacional Rey Malabo (Malabo). Le thème du mois : \textbf{« Mis hábitos digitales : preferencias, rutinas y brecha digital »}. Chaque groupe (4 élèves) représente une « délégation » de son établissement et doit publier un message initial dans le fil de discussion. Les messages seront lus par tous les participants hispanophones du projet. La langue de travail est l'espagnol.
\end{textebox}

\begin{textebox}[Consignes officielles du forum (extrait du règlement)]
\textbf{BASES DE PARTICIPACIÓN EN EL FORO « PUENTES DIGITALES »}

1. Cada grupo redacta UN mensaje inicial (8--10 líneas, sin contar el saludo y la firma).
2. El mensaje DEBE incluir obligatoriamente:
   \begin{itemize}
       \item[a)] Un título claro: « Mis hábitos digitales ».
       \item[b)] Una presentación del grupo (país, ciudad, tipo de conexión).
       \item[c)] Las herramientas y redes preferidas con \emph{preferir} y \emph{gustar más}.
       \item[d)] Las rutinas actuales con \emph{soler}, \emph{acostumbrar a}, \emph{tener la costumbre de}.
       \item[e)] Una comparación con el pasado (antes vs. ahora) usando el imperfecto o \emph{solía}.
       \item[f)] Una reflexión breve sobre la \emph{brecha digital} en su entorno (Camerún, Guinea Ecuatorial, zonas rurales de España).
   \end{itemize}
3. Léxico mínimo obligatorio (5 ítems): \emph{red social, contraseña, navegar, teléfono inteligente, brecha digital, descargar, conexión, ancho de banda, ciberseguridad, influenciador}.
4. Tras la publicación, cada grupo lee su mensaje en clase. Los otros grupos formulan DOS preguntas de comprensión (escritas en el chat del foro).
5. El profesor recoge las producciones, corrige colectivamente la diphtongaison de \emph{preferir}, la construcción de \emph{gustar} y los acentos.
\end{textebox}

\begin{textebox}[Exemple de message modèle publié par le groupe « Madrid Centro » (texte support pour analyse)]
\textbf{Asunto: Mis hábitos digitales}

Hola a todos:

Somos el grupo « Madrid Centro », alumnos de 2º de Bachillerato del Instituto Cervantes. Tenemos conexión de fibra óptica (600 Mbps) y cada uno usa su teléfono inteligente.

En cuanto a preferencias, \textbf{preferimos Instagram a TikTok} porque \textbf{nos gusta más} la fotografía que los vídeos cortos. También \textbf{preferimos WhatsApp a Telegram} para la comunicación diaria.

Respecto a hábitos: \textbf{solíamos navegar} mucho por la tarde, pero ahora \textbf{solemos estudiar} primero. \textbf{Tenemos la costumbre de revisar} las notificaciones al despertar y \textbf{acostumbramos a descargar} aplicaciones educativas. Antes \textbf{pasábamos} horas en videojuegos; hoy \textbf{dedicamos} ese tiempo a proyectos colaborativos.

Sobre la \textbf{brecha digital}: en los barrios periféricos de Madrid hay familias sin ancho de banda suficiente. La \emph{ciberseguridad} y la gestión de \emph{contraseñas} siguen siendo retos pendientes.

¡Saludos desde la capital!

Grupo Madrid Centro
\end{textebox}

\begin{definition}[Glossaire du texte modèle (Madrid Centro)]
\begin{description}
    \item[\esp{fibra óptica}] fibre optique
    \item[\esp{en cuanto a}] en ce qui concerne, quant à
    \item[\esp{fotografía}] photographie
    \item[\esp{comunicación diaria}] communication quotidienne
    \item[\esp{solíamos}] nous avions l'habitude (imparfait de \esp{soler})
    \item[\esp{solemos}] nous avons l'habitude (présent de \esp{soler})
    \item[\esp{revisar}] vérifier, consulter
    \item[\esp{notificaciones}] notifications
    \item[\esp{al despertar}] au réveil / en se réveillant
    \item[\esp{aplicaciones educativas}] applications éducatives
    \item[\esp{barrios periféricos}] quartiers périphériques
    \item[\esp{retos pendientes}] défis en suspens / restants
\end{description}
\end{definition}

\courssub{Consignes}

\begin{enumerate}
    \item \textbf{Analyse du modèle (10 min).} En groupe, lisez le message du groupe « Madrid Centro ». Relevez et classez dans un tableau :
    \begin{itemize}
        \item les 5 items de lexique obligatoire utilisés ;
        \item toutes les occurrences de \esp{preferir} (noter la diphtongaison) et de \esp{gustar más} ;
        \item toutes les structures d'habitude présente (\esp{soler, acostumbrar a, tener la costumbre de}) et passée (\esp{solía, imperfecto}) ;
        \item les marqueurs de relation logique (\esp{en cuanto a, respecto a, antes / ahora, también}).
    \end{itemize}
    \item \textbf{Préparation de la production (25 min).} Rédigez le message de votre groupe en respectant scrupuleusement le cahier des charges (6 points obligatoires + 5 items lexicaux minimum). Répartissez les rôles : \emph{redactor principal} (écrit), \emph{corrector de léxico} (vérifie les 5 mots), \emph{corrector de gramática} (vérifie \esp{preferir, gustar, soler, acostumbra}), \emph{portavoz} (lira le message).
    \item \textbf{Relecture croisée (10 min).} Échangez votre brouillon avec un groupe voisin. Vérifiez avec une grille :
    \begin{itemize}
        \item Diphtongaison de \esp{preferir} : formes \textbf{avec} diphtongue (\esp{prefiero, prefieres, prefiere, prefieren}) ; formes \textbf{sans} diphtongue (\esp{preferimos, preferís}). Vérifier qu'on n'écrit pas \esp{*prefierimos, *prefierís}.
        \item Construction de \esp{gustar} : \esp{me gusta + infinitif / nom singulier} ; \esp{me gustan + nom pluriel}.
        \item \esp{Preferir algo A algo} (préposition \esp{a} pour la comparaison).
        \item \esp{Acostumbrar A + infinitif} (préposition \esp{a} obligatoire).
        \item Accents : \esp{contraseña, teléfono, conexión, hábito, también, más, sí}.
    \end{itemize}
    \item \textbf{Publication et lecture orale (15 min).} Chaque \emph{portavoz} lit le message final devant la classe. Les autres groupes notent deux questions de compréhension (ex. : \esp{¿Qué red social prefieren? ¿Cuál era su hábito antes?}).
    \item \textbf{Interaction écrite (10 min).} Les questions sont posées dans le « chat » du forum (ou au tableau). Les groupes concernés répondent oralement en espagnol.
    \item \textbf{Correction collective guidée (15 min).} L'enseignant projette 3--4 extraits anonymisés contenant des erreurs types. La classe corrige collectivement ; l'enseignant explicite la règle (diphtongaison, \esp{gustar}, \esp{preferir a}, \esp{acostumbrar a}, imparfait vs \esp{solía}).
\end{enumerate}

\courssub{Ressources à mobiliser}

\begin{propriete}[Rappel : Expression de la préférence]
\textbf{1. \esp{Preferir} (verbe à diphtongaison radicale \emph{e} $\to$ \emph{ie} aux trois personnes du singulier et à la 3e pers. pl. du présent de l'indicatif)} \\
\esp{Prefiero el correo electrónico a los mensajes de texto.} « Je préfère l'email aux SMS. » \\
\esp{¿Prefieres Instagram o TikTok?} « Tu préfères Instagram ou TikTok ? » \\
\emph{Construction :} \esp{preferir algo A algo} (préposition \esp{a} devant le second terme de comparaison). Ne pas confondre avec le français « préférer \textbf{à} » -- en espagnol, \esp{a} introduit le terme rejeté.

\textbf{2. \esp{Gustar más} (verbe à construction inverse : le sujet grammatical est la chose aimée)} \\
\esp{Me gusta más navegar que ver la tele.} « Je préfère naviguer que regarder la télé -- litt. "Me plaît plus naviguer..." » \\
\esp{Me gustan más las redes sociales que los foros.} (accord au pluriel avec \esp{redes sociales}). \\
\emph{Attention :} \esp{me gusta más} + infinitif ou nom singulier ; \esp{me gustan más} + nom pluriel. Pas de \esp{*prefiero más}.
\end{propriete}

\begin{propriete}[Rappel : Expression de l'habitude]
\begin{center}
\begin{tabularx}{\linewidth}{|X|X|X|}
\hline
\rowcolor{popblueL} Structure & Exemple & Nuance / Temps \\
\hline
\esp{soler + infinitif} & \esp{Suelo conectarme a las 20 h.} & Habitude \textbf{présente}, fréquence élevée, registre courant. \\
\hline
\esp{acostumbrar a + infinitif} & \esp{Acostumbro a revisar el correo al despertar.} & Habitude \textbf{présente}, bien ancrée, registre légèrement plus formel. \textbf{Préposition \esp{a} obligatoire}. \\
\hline
\esp{tener la costumbre de + infinitif} & \esp{Tengo la costumbre de cambiar la contraseña cada mes.} & Habitude \textbf{présente}, consciente, volontaire. \\
\hline
\esp{imperfecto de indicativo} & \esp{Antes navegaba sin cuidado; ahora uso VPN.} & Habitude \textbf{passée} (descriptive), sans marqueur aspectuel fort. \\
\hline
\esp{solía + infinitif} & \esp{Antes solía descargar apps sin leer permisos.} & Habitude \textbf{passée} révolue, contraste explicite avec le présent. \\
\hline
\end{tabularx}
\end{center}
\emph{Note :} \esp{soler} n'existe qu'aux temps où l'aspect habituel a du sens (présent, imparfait, conditionnel). Pas de *\esp{he solido}, *\esp{solere}.
\end{propriete}

\begin{definition}[Lexique obligatoire NTIC (minimum 5 items à intégrer)]
\begin{description}
    \item[\esp{red social}] réseau social (ex. \esp{Instagram, X, TikTok})
    \item[\esp{contraseña}] mot de passe
    \item[\esp{navegar}] naviguer (sur le web)
    \item[\esp{teléfono inteligente}] smartphone
    \item[\esp{brecha digital}] fracture numérique
    \item[\esp{descargar}] télécharger
    \item[\esp{conexión}] connexion (internet)
    \item[\esp{ancho de banda}] bande passante
    \item[\esp{ciberseguridad}] cybersécurité
    \item[\esp{influenciador}] influenceur / influenceuse (anglicisme fréquent : \esp{influencer})
\end{description}
\end{definition}

\begin{attention}[Erreurs fréquentes des francophones -- Points de vigilance pour la correction]
\begin{itemize}
    \item \textbf{Diphtongaison oubliée ou surgénéralisée :} \esp{*preferimos, *preferis} au lieu de \esp{preferimos, preferís} (pas de diphtongue aux 1re et 2e pers. pl.) MAIS \esp{prefiero, prefieres, prefiere, prefieren} (diphtongue). \emph{Astuce :} « \textbf{Noi} et \textbf{Voi} ne diphtonguent pas. »
    \item \textbf{Construction de \esp{gustar} :} \esp{*Me gusto más Instagram} (erreur : sujet = \esp{yo}). Correct : \esp{Me gusta más Instagram} (sujet = \esp{Instagram}). Pluriel : \esp{Me gustan más las redes}.
    \item \textbf{Préposition manquante :} \esp{*prefiero Instagram TikTok} $\to$ \esp{prefiero Instagram \textbf{a} TikTok}. \esp{*acostumbro revisar} $\to$ \esp{acostumbro \textbf{a} revisar}.
    \item \textbf{Accents :} \esp{contraseña, teléfono, conexión, hábito, también, más, sí} (vs \esp{si} conditionnel).
    \item \textbf{Faux ami :} \esp{actualidad} = actualité (pas « actualité » au sens de « présent » -- \esp{actualmente} = actuellement).
    \item \textbf{Anglicisme :} \esp{*descargar un app} $\to$ \esp{descargar una aplicación / una app} (féminin).
\end{itemize}
\end{attention}

\courssub{Proposition de production et corrigé commenté}

\begin{exoresolu}[Production modèle -- Groupe « Yaoundé Mokolo » (Lycée Général Leclerc)]
\textbf{Asunto: Mis hábitos digitales}

Hola a todos los participantes del foro « Puentes digitales »:

Somos el grupo « Yaoundé Mokolo », alumnos de Terminale A del Lycée Général Leclerc en Camerún. Nuestra \textbf{conexión} es por datos móviles (4G) en el \textbf{teléfono inteligente}; el \textbf{ancho de banda} varía según el barrio.

En cuanto a preferencias, \textbf{preferimos WhatsApp a Facebook} para comunicarnos porque \textbf{nos gusta más} la rapidez y el cifrado de extremo a extremo. También \textbf{preferimos TikTok a Instagram} para entretenimiento; \textbf{nos gustan más} los vídeos cortos.

Respecto a hábitos actuales: \textbf{solemos navegar} por las \textbf{redes sociales} dos horas al día. \textbf{Tenemos la costumbre de verificar} la \textbf{contraseña} con gestor de claves y \textbf{acostumbramos a activar} la verificación en dos pasos por \textbf{ciberseguridad}. Antes \textbf{solíamos descargar} aplicaciones sin leer permisos; ahora \textbf{leemos} las políticas de privacidad.

Sobre la \textbf{brecha digital}: en zonas rurales de Camerún muchas familias no tienen \textbf{conexión} estable ni \textbf{teléfono inteligente}. La electricidad irregular limita el acceso a la educación en línea. Urge invertir en infraestructura y formación.

¡Saludos desde Yaoundé!

Grupo Yaoundé Mokolo
\tcblower
\textbf{Commentaire détaillé des choix de langue}

\begin{enumerate}
    \item \textbf{Structure globale :} Respect du cahier des charges (titre, présentation, préférences, habitudes présentes, passé, brecha digital, 10 items lexicaux sur 10 possibles). Longueur : 9 lignes (hors salutation/signature).
    \item \textbf{Préférences :} \esp{preferimos WhatsApp a Facebook} (diphtongue \emph{preferimos} -- 1re pers. pl. sans diphtongue, correct) ; \esp{nos gusta más la rapidez} (singulier) ; \esp{nos gustan más los vídeos} (pluriel). Utilisation de \esp{en cuanto a} et \esp{también} pour lier les paragraphes.
    \item \textbf{Habitudes présentes :} \esp{solemos navegar} (habitude fréquente) ; \esp{tenemos la costumbre de verificar} (habitude consciente, sécurité) ; \esp{acostumbramos a activar} (préposition \esp{a} respectée, 1re pers. pl. sans diphtongue). Variété des trois structures exigées.
    \item \textbf{Habitude passée :} \esp{Antes solíamos descargar} (imparfait de \esp{soler} = \esp{solía} à toutes les personnes) + contraste \esp{ahora leemos} (présent). Alternative possible : \esp{Antes descargábamos} (imparfait simple).
    \item \textbf{Brecha digital contextualisée :} Référence concrète au Cameroun (zones rurales, électricité, éducation en ligne). Vocabulaire technique réinvesti : \esp{ancho de banda, ciberseguridad, verificación en dos pasos, gestor de claves}.
    \item \textbf{Orthographe / accents :} \esp{conexión, teléfono, contraseña, ciberseguridad, también, más, sí} (dans \esp{sí} -- absent ici mais rappelé). Pas d'anglicisme (\esp{app} évité, \esp{aplicación} implicite).
    \item \textbf{Registre :} Formel mais naturel pour un forum scolaire (\esp{ustedes} implicite dans \emph{saludos}, \esp{nuestra, nuestros}).
    \item \textbf{Glossaire complémentaire :} \esp{cifrado de extremo a extremo} (chiffrement de bout en bout), \esp{gestor de claves} (gestionnaire de mots de passe), \esp{verificación en dos pasos} (double authentification), \esp{políticas de privacidad} (politiques de confidentialité), \esp{infraestructura} (infrastructure).
\end{enumerate}
\end{exoresolu}

\begin{exoresolu}[Production modèle -- Groupe « Malabo Sur » (Colegio Nacional Rey Malabo)]
\textbf{Asunto: Mis hábitos digitales}

Estimados compañeros:

Somos el grupo « Malabo Sur », estudiantes de 2º de Bachillerato en Guinea Ecuatorial. Disponemos de \textbf{conexión} por fibra en el centro, pero en casa muchos usan datos móviles limitados.

Nuestras preferencias: \textbf{preferimos Facebook a X} porque \textbf{nos gusta más} mantener grupos de clase y familia. \textbf{Preferimos YouTube a TikTok} para tutoriales; \textbf{me gusta más} aprender viendo vídeos largos.

Hábitos: \textbf{suelo revisar} el correo académico cada mañana. \textbf{Acostumbro a subir} apuntes a la nube y \textbf{tengo la costumbre de cambiar} la \textbf{contraseña} cada trimestre. Antes \textbf{navegaba} sin antivirus; hoy \textbf{uso} protección y \textbf{descargo} solo de tiendas oficiales.

La \textbf{brecha digital} en Guinea Ecuatorial es marcada: en la zona insular (Annobón, Corisco) el \textbf{ancho de banda} es casi nulo. Falta formación en \textbf{ciberseguridad} para evitar estafas (\emph{phishing}). Los \textbf{influenciadores} locales empiezan a concienciar.

Un abrazo desde Malabo.

Grupo Malabo Sur
\tcblower
\textbf{Commentaire détaillé des choix de langue}

\begin{enumerate}
    \item \textbf{Variété des sujets :} Alternance \emph{nosotros} (\esp{preferimos, solemos}) et \emph{yo} (\esp{suelo, acostumbro, tengo, uso, descargo}) -- montre la maîtrise des personnes. Le \emph{portavoz} lit au nom du groupe mais le texte intègre des habitudes individuelles (\esp{suelo, acostumbro}) pour plus d'authenticité.
    \item \textbf{Préférences :} \esp{preferimos Facebook a X} (nouveau nom de Twitter, actualité) ; \esp{nos gusta más mantener} (infinitif sujet) ; \esp{me gusta más aprender} (1re pers. sg. de \esp{gustar} -- attention : \esp{me gusta}, pas \esp{*me gusto}).
    \item \textbf{Habitudes :} \esp{suelo revisar} (présent) ; \esp{acostumbro a subir} (préposition \esp{a}) ; \esp{tengo la costumbre de cambiar} (structure nominale). Trois structures distinctes bien employées.
    \item \textbf{Passé / présent :} \esp{Antes navegaba sin antivirus} (imparfait descriptif) ; \esp{hoy uso / descargo} (présent de rupture). Pas de \esp{solía} ici mais l'imparfait suffit ; l'alternance est acceptée.
    \item \textbf{Brecha digital :} Contexte guinéen spécifique (zone insulaire, \esp{ancho de banda} nul, \esp{phishing}, rôle des \esp{influenciadores}). Lexique : 9/10 items utilisés (\esp{red social} absent mais \esp{Facebook, YouTube, TikTok, X} cités).
    \item \textbf{Points de vigilance respectés :} Diphtongaison \esp{preferimos} (1re pl. sans diphtongue) ; \esp{acostumbro a} (préposition) ; accents \esp{conexión, contraseña, ciberseguridad, también, más}.
    \item \textbf{Glossaire complémentaire :} \esp{zona insular} (zone insulaire), \esp{casi nulo} (quasi nul), \esp{estafas} (escroqueries), \esp{concienciar} (sensibiliser), \esp{tiendas oficiales} (stores officiels : Play Store, App Store).
\end{enumerate}
\end{exoresolu}

\courssub{Bilan de l'activité}

Cette activité de forum numérique a permis une intégration effective des trois piliers de la leçon E28 :

\begin{itemize}
    \item \textbf{Lexique NTIC :} Les élèves ont mobilisé 8 à 10 items sur la liste de 10 dans un contexte signifiant (description réelle de leur environnement technique : 4G à Yaoundé, fibre à Malabo centre, zones blanches insulaires). L'ancrage camerounais et guinéen a donné du sens aux mots \esp{brecha digital, ancho de banda, conexión}.
    \item \textbf{Grammaire -- Préférence :} La diphtongaison de \esp{preferir} a été massivement respectée grâce à la relecture croisée ciblée. L'erreur résiduelle principale observée : \esp{*prefiero más} (calque de « je préfère plus ») -- corrigée en rappelant l'exclusivité de \esp{preferir a} ou \esp{gustar más}. La construction de \esp{gustar} (accord avec le sujet réel) a nécessité une remédiation orale pour deux groupes (\esp{*me gustan más WhatsApp} $\to$ \esp{me gusta más WhatsApp}).
    \item \textbf{Grammaire -- Habitude :} Les trois structures présentes (\esp{soler, acostumbrar a, tener la costumbre de}) ont été utilisées par tous les groupes. La préposition \esp{a} après \esp{acostumbrar} a été acquise (0 erreur en version finale). L'opposition passé/présent (\esp{antes solía / navegaba} vs \esp{ahora suelo / uso}) a structuré le paragraphe « évolution ».
    \item \textbf{Compétences transverses :} Production écrite collaborative (négociation de sens, répartition des rôles), expression orale (lecture à voix haute, prosodie, intonation interrogative pour les questions), compréhension orale/écrite (questions des pairs), interaction authentique (forum simulé mais règles réelles).
    \item \textbf{Évaluation formative :} La grille de relecture croisée (diphtongaison, \esp{gustar}, \esp{a}, accents) est un outil réutilisable pour toute production écrite future. La correction collective au tableau a ancré les règles par l'erreur analysée (pédagogie de l'erreur).
\end{itemize}

\begin{aretenir}[Synthèse pour l'élève -- Fiche mémo à coller dans le cahier]
\textbf{FORUM NTIC -- CHECK-LIST AVANT DE POSTER}
\begin{enumerate}
    \item Titre : \esp{Mis hábitos digitales}
    \item Présentation : pays, ville, type de \esp{conexión}.
    \item Préférences : \esp{prefiero A a B} (diphtongue !) + \esp{me gusta(n) más}.
    \item Habitudes NOW : \esp{suelo / acostumbro a / tengo la costumbre de} + INF.
    \item Habitudes BEFORE : \esp{antes solía / imperfecto} + contraste \esp{ahora}.
    \item \esp{Brecha digital} : exemple concret (Cameroun / Guinée / Espagne rural).
    \item 5+ mots obligatoires : \esp{red social, contraseña, navegar, teléfono inteligente, brecha digital, descargar, conexión, ancho de banda, ciberseguridad, influenciador}.
    \item Accents : \esp{contraseña, teléfono, conexión, también, más, sí}.
    \item Relecture : diphtongaison \esp{preferir}, \esp{gustar} accord, \esp{a} après \esp{preferir/acostumbrar}.
\end{enumerate}
\end{aretenir}

\begin{savaistu}[Le saviez-vous ? -- La brecha digital en Afrique centrale]
Selon les dernières estimations de l'UIT (Union internationale des télécommunications), le taux de pénétration d'Internet au Cameroun tourne autour de 40 \% (contre plus de 90 \% en Espagne). En Guinée équatoriale, la fracture est forte entre l'île de Bioko (Malabo, fibre) et la zone continentale (Bata, 3G limitée) et l'île d'Annobón (quasi pas de connexion). Le projet \emph{« África Digital »} de l'Union africaine vise à connecter 100\% des écoles à l'horizon 2030. Au Lycée Général Leclerc, le club \emph{« Jóvenes Conectados »} forme les élèves à la \esp{ciberseguridad} : création de \esp{contraseñas} robustes, détection de \emph{phishing}, gestion de l'empreinte numérique. Votre message de forum peut être lu par des pairs de Madrid ou Malabo -- soyez des ambassadeurs numériques responsables !
\end{savaistu}

\newpage

\coursec{Méthodes et exercices résolus}

\coursec{Leçon E28 : Lectura y comentario : Internet y las NTIC ; traduction : expression de la préférence et de l'habitude}

\courssub{Texte support : Internet y las NTIC}

\begin{textebox}[Article d'opinion original — Contexte camerounais et africain]
En muchas ciudades africanas, los jóvenes navegan por Internet con su \textbf{teléfono inteligente} para estudiar, informarse y comunicarse. Sin embargo, la \textbf{brecha digital} sigue siendo grande: en los pueblos, muchas familias no tienen ni ordenador ni conexión. Por eso, algunos alumnos deben ir al \textbf{cibercafé} para \textbf{descargar} sus documentos. Las \textbf{redes sociales} como WhatsApp o Facebook facilitan el contacto con la diáspora, pero exigen una \textbf{contraseña} segura. Navegar con responsabilidad es el reto de hoy.
\end{textebox}

\emph{Traduction :} « Dans de nombreuses villes africaines, les jeunes naviguent sur Internet avec leur téléphone intelligent pour étudier, s'informer et communiquer. Cependant, la fracture numérique reste importante : dans les villages, de nombreuses familles n'ont ni ordinateur ni connexion. C'est pourquoi certains élèves doivent aller au cybercafé pour télécharger leurs documents. Les réseaux sociaux comme WhatsApp ou Facebook facilitent le contact avec la diaspora, mais exigent un mot de passe sûr. Naviguer avec responsabilité est le défi d'aujourd'hui. »

\courssub{Compréhension du texte}

\begin{exoresolu}[Compréhension de texte : type Bac]
Énoncé. Lis le texte ci-dessus puis réponds en espagnol aux questions.

Questions : 
1) ¿Para qué usan Internet los jóvenes? 
2) ¿Qué es la brecha digital según el texto? 
3) ¿Por qué van algunos alumnos al cibercafé? 
4) ¿Y tú, para qué usas Internet?
\tcblower
\textbf{Raisonnement} : chaque réponse doit être justifiée par une citation courte du texte, puis reformulée avec ses propres mots.

1) \esp{Los jóvenes usan Internet «para estudiar, informarse y comunicarse», es decir, para su educación y su vida social.} — On reprend les trois verbes du texte et on ajoute une reformulation (\esp{es decir}) qui montre la compréhension.

2) \esp{La brecha digital es la desigualdad de acceso a las nuevas tecnologías: el texto dice que «en los pueblos, muchas familias no tienen ni ordenador ni conexión».} — On définit le terme avec ses propres mots (\esp{desigualdad de acceso}) et on cite la preuve.

3) \esp{Algunos alumnos van al cibercafé «para descargar sus documentos» porque no tienen conexión en casa.} — Le connecteur \esp{porque} explicite la cause sous-entendue par \esp{por eso}.

4) Réponse personnelle possible : \esp{Yo uso Internet sobre todo para hacer mis deberes, ver vídeos educativos y hablar con mis amigos por WhatsApp. También consulto las noticias de mi país.} — Vocabulaire NTIC mobilisé, première personne, deux activités justifiées.

\textbf{Conclusion méthodologique} : une bonne réponse de compréhension = citation + reformulation + (pour l'opinion) exemple personnel concret.
\end{exoresolu}

\courssub{Lexique thématique des NTIC}

\begin{definition}[Vocabulaire essentiel — Module 5]
\begin{center}
\begin{tabularx}{\linewidth}{|X|X|X|}
\hline
\rowcolor{popblueL} \textbf{Español} & \textbf{Français} & \textbf{Exemple / Contexte} \\
\hline
\esp{la red social} & le réseau social & \esp{Uso las redes sociales para hablar con mis amigos.} \\
\hline
\esp{la contraseña} & le mot de passe & \esp{Olvidé mi contraseña de correo electrónico.} \\
\hline
\esp{descargar} & télécharger (depuis Internet) & \esp{Voy a descargar la aplicación del banco.} \\
\hline
\esp{navegar (por Internet)} & naviguer (sur Internet) & \esp{Me gusta navegar por la red por la noche.} \\
\hline
\esp{el teléfono inteligente / el smartphone} & le smartphone & \esp{Mi teléfono inteligente tiene buena cámara.} \\
\hline
\esp{la brecha digital} & la fracture numérique & \esp{La brecha digital afecta a las zonas rurales.} \\
\hline
\esp{el cibercafé} & le cybercafé & \esp{Fui al cibercafé a imprimir mi CV.} \\
\hline
\esp{la conexión (a Internet)} & la connexion (Internet) & \esp{La conexión es lenta en mi barrio.} \\
\hline
\esp{el ordenador / la computadora} & l'ordinateur & \esp{Necesito un ordenador para programar.} \\
\hline
\esp{subir / colgar (archivos)} & mettre en ligne / uploader & \esp{Subí las fotos a la nube.} \\
\hline
\end{tabularx}
\end{center}
\emph{Note :} \esp{ordenador} (Espagne) / \esp{computadora} (Amérique latine, Guinée équatoriale) ; \esp{móvil} (Espagne) / \esp{celular} (Amérique). Au Cameroun, l'usage suit la norme péninsulaire (\esp{ordenador, móvil}).
\end{definition}

\courssub{Grammaire 1 : Expression de la préférence}

\begin{propriete}[Règle : Exprimer la préférence]
\begin{itemize}
\item \textbf{Preferir (e $\rightarrow$ ie) + nom / infinitif} : \esp{Prefiero el té.} / \esp{Prefiero navegar.}
\item \textbf{Preferir X a Y} (comparaison de noms ou d'infinitifs) : \esp{Prefiero el fútbol al baloncesto.} \esp{Prefiero leer a ver la tele.} (Pour les verbes, \esp{antes que} + subjonctif est plus élégant : \esp{Prefiero leer antes que ver la tele.}).
\item \textbf{Gustar más X que Y} : \esp{Me gusta más el té que el café.} \esp{Me gusta más leer que escribir.}
\item \textbf{Conjugaison de \textit{preferir} (présent indicatif)} : \esp{prefiero, prefieres, prefiere, preferimos, preferís, prefieren.} (Pas de diphtongue à \esp{nosotros/as} et \esp{vosotros/as}).
\item \textbf{Devant infinitif} : pas de préposition (\esp{Prefiero estudiar}, non *\esp{prefiero a estudiar}).
\end{itemize}
\end{propriete}

\courssub{Fiche méthode 2 : traduire l'expression de la préférence (thème)}

\begin{methode}[Traduire « préférer, aimer mieux » en espagnol]
\begin{enumerate}
\item \textbf{Identifier la structure française} : « préférer X à Y », « aimer mieux X que Y », « j'aime plus... ».
\item \textbf{Choisir l'équivalent espagnol} :
\begin{itemize}
\item \esp{preferir (e $\rightarrow$ ie) + nom/infinitif} : \esp{Prefiero el té.} « Je préfère le thé. »
\item \esp{preferir X a Y} : \esp{Prefiero el fútbol al baloncesto.} « Je préfère le football au basket. » (Attention : \esp{a + el = al})
\item \esp{me gusta más X que Y} : \esp{Me gusta más leer que ver la tele.} « J'aime mieux lire que regarder la télé. »
\end{itemize}
\item \textbf{Conjuguer correctement preferir} : verbe à changement radical e $\rightarrow$ ie : \esp{prefiero, prefieres, prefiere, preferimos, preferís, prefieren}.
\item \textbf{Devant un infinitif}, pas de préposition en espagnol : \esp{Prefiero navegar por Internet.} « Je préfère naviguer sur Internet. »
\item \textbf{Vérifier} : article devant le nom (\esp{el té}), contraction \esp{al}, accent absent sur \esp{prefiero} mais présent au passé simple \esp{preferí}.
\end{enumerate}
\end{methode}

\begin{exoresolu}[Thème : la préférence]
Énoncé. Traduis en espagnol : 
1) Je préfère les réseaux sociaux aux journaux. 
2) Mon frère aime mieux télécharger de la musique que regarder la télévision. 
3) Nous préférons étudier à la bibliothèque. 
4) Tu préfères le téléphone portable, n'est-ce pas ?
\tcblower
1) \esp{Prefiero las redes sociales a los periódicos.} — Structure \esp{preferir X a Y} ; « réseaux sociaux » = \esp{redes sociales} (féminin pluriel) ; \esp{a + los = a los} (pas de contraction avec \esp{los}).

2) \esp{A mi hermano le gusta más descargar música que ver la televisión.} (Recommandé pour « aimer mieux » + infinitifs). 
   \emph{Alternative avec \textit{preferir}} : \esp{Mi hermano prefiere descargar música a ver la televisión.} (Préposition \esp{a}, pas \esp{que}). 
   — Avec \esp{gustar}, ne pas oublier le pronom \esp{le} et la personne introduite par \esp{a}.

3) \esp{Preferimos estudiar en la biblioteca.} — Première personne du pluriel : pas de changement radical (\esp{preferimos}, jamais *\esp{prefierimos}) ; « à la bibliothèque » = \esp{en la biblioteca} (lieu où l'on est).

4) \esp{Prefieres el teléfono móvil, ¿verdad? / ¿no?} — Deuxième personne : \esp{prefieres} avec diphtongue \esp{ie} ; « n'est-ce pas » = \esp{¿verdad?} ou \esp{¿no?} entre signes d'interrogation.

\textbf{Conclusion} : la préférence exige trois contrôles : changement radical e $\rightarrow$ ie (sauf nous/vous), préposition \esp{a} devant le terme rejeté avec \esp{preferir}, article devant chaque nom.
\end{exoresolu}

\courssub{Grammaire 2 : Expression de l'habitude}

\begin{propriete}[Règle : Exprimer l'habitude (présent et passé)]
\begin{itemize}
\item \textbf{Habitude présente} :
\begin{itemize}
\item \esp{soler (o $\rightarrow$ ue) + infinitif} : \esp{Suelo navegar por Internet por la noche.} « J'ai l'habitude de naviguer le soir. »
\item \esp{acostumbrar a + infinitif} : \esp{Acostumbro a leer antes de dormir.} « J'ai l'habitude de lire avant de dormir. »
\item \esp{tener la costumbre de + infinitif} : \esp{Tengo la costumbre de apagar el móvil en clase.} « J'ai l'habitude d'éteindre mon portable en classe. »
\end{itemize}
\item \textbf{Habitude passée} : \textbf{Imparfait de l'indicatif} (valeur itérative) : \esp{Cuando era niño, jugaba al fútbol todos los días.} « Quand j'étais enfant, je jouais au foot tous les jours. »
\item \textbf{Soler à l'imparfait} : \esp{solía, solías, solía, solíamos, solíais, solían} + infinitif = « j'avais l'habitude de... ».
\item \textbf{Marqueurs temporels} : \esp{siempre, todos los días, a menudo, antes, antaño, de niño, normalmente, por lo general}.
\item \textbf{« D'habitude » (adverbe)} = \esp{normalmente, por lo general, de costumbre, generalmente} (sans verbe \esp{soler}).
\end{itemize}
\end{propriete}

\begin{center}
\begin{tabular}{|l|l|l|l|l|l|l|}
\hline
\rowcolor{popblueL} \textbf{Persona} & \textbf{Soler (Pres.)} & \textbf{Soler (Pret. Imp.)} & \textbf{Acostumbrar (Pres.)} & \textbf{Tener la costumbre (Pres.)} & \textbf{Verbe principal (Imp.)} \\
\hline
yo & suelo & solía & acostumbro & tengo la costumbre & jugaba / iba / hacía \\
tú & sueles & solías & acostumbras & tienes la costumbre & jugabas / ibas / hacías \\
él/ella/usted & suele & solía & acostumbra & tiene la costumbre & jugaba / iba / hacía \\
nosotros/as & solemos & solíamos & acostumbramos & tenemos la costumbre & jugábamos / íbamos / hacíamos \\
vosotros/as & soléis & solíais & acostumbráis & tenéis la costumbre & jugabais / ibais / hacíais \\
ellos/as/ustedes & suelen & solían & acostumbran & tienen la costumbre & jugaban / iban / hacían \\
\hline
\end{tabular}
\end{center}

\courssub{Fiche méthode 3 : traduire l'expression de l'habitude (version)}

\begin{methode}[Traduire « avoir l'habitude de, d'habitude » en espagnol vers le français]
\begin{enumerate}
\item \textbf{Repérer le type d'habitude} : présente (\esp{suelo, acostumbro, tengo la costumbre, normalmente}) ou passée (\esp{solía, imparfait simple}).
\item \textbf{Habitude présente} $\rightarrow$ français présent + « avoir l'habitude de » / « d'habitude » / « généralement » : \esp{Suelo navegar} $\rightarrow$ « J'ai l'habitude de naviguer » / « D'habitude, je navigue ».
\item \textbf{Habitude passée} ($\esp{solía}$ ou imparfait \esp{-aba/-ía}) $\rightarrow$ français imparfait : \esp{Solía estudiar} / \esp{Estudiaba} $\rightarrow$ « J'avais l'habitude d'étudier » / « J'étudiais ».
\item \textbf{Attention aux marqueurs} : \esp{antes, de niño, todos los días, una vez por semana} signalent le passé $\rightarrow$ imparfait en français.
\item \textbf{\esp{Normalmente / Por lo general}} au début de phrase $\rightarrow$ « D'habitude / Normalement » + présent (si contexte présent) ou imparfait (si récit passé).
\end{enumerate}
\end{methode}

\begin{exoresolu}[Version : l'habitude]
Énoncé. Traduis en français : 
\esp{Antes, los alumnos solían estudiar sin Internet. Mi abuelo tiene la costumbre de escuchar la radio todas las mañanas. Cuando vivíamos en el pueblo, íbamos al cibercafé una vez por semana. Normalmente, descargo mis lecciones por la noche.}
\tcblower
« Autrefois, les élèves \textbf{avaient l'habitude d'étudier} sans Internet. Mon grand-père \textbf{a l'habitude d'écouter} la radio tous les matins. Quand nous vivions au village, \textbf{nous allions} au cybercafé une fois par semaine. \textbf{D'habitude}, je télécharge mes leçons le soir. »

\textbf{Justifications} :
\begin{itemize}
\item \esp{solían estudiar} : imparfait de \esp{soler} + infinitif = habitude révolue « avaient l'habitude de » ; \esp{antes} confirme le passé.
\item \esp{tiene la costumbre de} : présent = habitude actuelle « a l'habitude de » ; \esp{todas las mañanas} = « tous les matins ».
\item \esp{íbamos} : imparfait de \esp{ir}, habitude répétée dans le passé (« nous allions »), signalée par \esp{una vez por semana}.
\item \esp{Normalmente} : adverbe = « d'habitude / normalement » ; \esp{descargo} au présent car habitude présente ; \esp{por la noche} = « le soir ».
\end{itemize}
\textbf{Conclusion} : en version, tout indice de répétition (\esp{siempre, todos los días, una vez por semana}) doit déclencher le réflexe « habitude » : imparfait en français si le contexte est passé.
\end{exoresolu}

\courssub{Entraînement : Transformation de phrases (préférence et habitude)}

\begin{methode}[Transformer une phrase sans changer le sens]
\begin{enumerate}
\item \textbf{Repérer la structure de départ} : verbe de préférence (\esp{preferir}, \esp{gustar más}) ou d'habitude (\esp{soler}, imparfait).
\item \textbf{Choisir la structure cible synonyme} : \esp{prefiero X a Y} $\leftrightarrow$ \esp{me gusta más X que Y} ; \esp{suelo + inf.} $\leftrightarrow$ \esp{acostumbro a + inf.} $\leftrightarrow$ \esp{normalmente + présent}.
\item \textbf{Adapter la conjugaison} : avec \esp{gustar}, le verbe s'accorde avec la chose aimée et la personne devient pronom (\esp{me, te, le...}) : \esp{Prefiero el café.} $\rightarrow$ \esp{Me gusta más el café.}
\item \textbf{Conserver le temps} : présent $\rightarrow$ présent ; imparfait $\rightarrow$ imparfait (\esp{solía} $\leftrightarrow$ imparfait simple du verbe).
\item \textbf{Relire} : vérifier accents, articles et contractions (\esp{al}, \esp{del}).
\end{enumerate}
\end{methode}

\begin{exoresolu}[Transformation : type Bac]
Énoncé. Réécris chaque phrase avec l'expression entre parenthèses, sans changer le sens. 
1) \esp{Prefiero el fútbol al baloncesto.} (\esp{gustar más}) 
2) \esp{Suelo navegar por Internet después de clase.} (\esp{acostumbrar}) 
3) \esp{De niña, María leía muchos cuentos.} (\esp{solía}) 
4) \esp{Me gusta más el té que el café.} (\esp{preferir})
\tcblower
1) \esp{Me gusta más el fútbol que el baloncesto.} — \esp{gustar} à la 3e personne du singulier car ce qui plaît est \esp{el fútbol} ; la préposition \esp{a} de \esp{preferir} devient \esp{que} devant le terme comparé.

2) \esp{Acostumbro a navegar por Internet después de clase.} — \esp{acostumbrar a + infinitif} ; verbe régulier en -ar, première personne \esp{acostumbro} ; l'infinitif \esp{navegar} reste inchangé.

3) \esp{De niña, María solía leer muchos cuentos.} — \esp{solía} : imparfait de \esp{soler} (3e personne, sans changement radical à l'imparfait) + infinitif \esp{leer} ; le sens d'habitude passée est conservé.

4) \esp{Prefiero el té al café.} — \esp{a + el = al} (contraction obligatoire) ; changement radical : \esp{prefiero} avec \esp{ie}.

\textbf{Conclusion} : en transformation, seule la structure change ; le temps, la personne et le sens doivent rester strictement identiques.
\end{exoresolu}

\courssub{Méthode du commentaire et production guidée}

\begin{methode}[Commenter un texte sur les NTIC en 5 étapes]
\begin{enumerate}
\item \textbf{Présenter le document} : type de texte (article de presse, essai, blog), thème (Internet, réseaux sociaux, fracture numérique), situation d'énonciation. Formule utile : \esp{El texto es un artículo periodístico que trata de...}
\item \textbf{Dégager l'idée générale} en une phrase : \esp{La idea principal es que Internet transforma la vida cotidiana, pero no todos tienen acceso.}
\item \textbf{Repérer les arguments} : avantages (\esp{ventajas}: comunicación rápida, acceso a la información) et inconvénients (\esp{desventajas}: adicción, brecha digital, desinformación). Souligner les connecteurs : \esp{sin embargo}, \esp{además}, \esp{por eso}, \esp{en cambio}.
\item \textbf{Analyser le lexique des NTIC} : relever les mots-clés (\esp{red social, contraseña, descargar, navegar, teléfono inteligente, brecha digital}) et montrer comment ils construisent le thème.
\item \textbf{Donner un avis personnel argumenté} : \esp{En mi opinión... / Estoy de acuerdo en que... porque...} Toujours justifier avec un exemple concret (Cameroun, Guinée équatoriale, ton quartier, ton lycée).
\end{enumerate}
Réflexes : citer le texte entre guillemets pour justifier chaque réponse ; ne jamais recopier tout un paragraphe ; répondre en phrases complètes en espagnol.
\end{methode}

\begin{methode}[Rédiger un paragraphe d'opinion (8 à 12 lignes)]
\begin{enumerate}
\item \textbf{Introduction du sujet} : \esp{Hoy en día, Internet forma parte de nuestra vida cotidiana.}
\item \textbf{Annonce de ton point de vue} : \esp{En mi opinión... / A mi juicio... / Pienso que...}
\item \textbf{Deux arguments} reliés par des connecteurs : \esp{En primer lugar... Además...} Utilise le lexique du thème : \esp{red social, navegar, descargar, brecha digital}.
\item \textbf{Une nuance} : \esp{Sin embargo, no hay que olvidar que...}
\item \textbf{Conclusion personnelle} avec préférence ou habitude : \esp{Por eso, prefiero... / Por eso, suelo...}
\item \textbf{Relecture} : accords (genre/nombre), accents, ponctuation ¿ ¡, verbes conjugués à la bonne personne.
\end{enumerate}
\end{methode}

\begin{exoresolu}[Rédaction guidée : type Bac]
Énoncé. Rédige un paragraphe de 8 à 10 lignes : « ¿Es Internet una ventaja o un problema para los jóvenes? Da tu opinión con ejemplos de tu vida. »
\tcblower
\textbf{Plan} : opinion (avantage) $\rightarrow$ argument 1 (études) $\rightarrow$ argument 2 (communication) $\rightarrow$ nuance (addiction, brecha digital) $\rightarrow$ conclusion avec habitude/préférence.

\textbf{Production corrigée} :

\esp{Hoy en día, Internet forma parte de la vida de los jóvenes. En mi opinión, es una gran ventaja, aunque tiene algunos peligros. En primer lugar, gracias a mi \textbf{teléfono inteligente}, puedo \textbf{navegar} por Internet para hacer mis deberes y \textbf{descargar} documentos para el \textbf{instituto}. Además, las \textbf{redes sociales} me permiten comunicarme con mis amigos y con mi familia que vive lejos. \textbf{Sin embargo}, no hay que olvidar que algunos jóvenes pasan demasiado tiempo en línea y que la \textbf{brecha digital} deja fuera a los alumnos de los pueblos. \textbf{Por eso}, \textbf{suelo} usar Internet con moderación: \textbf{prefiero} estudiar primero y conectarme después. \textbf{En conclusión}, Internet es una herramienta útil si sabemos usarla bien.}

\textbf{Justifications} : connecteurs variés (\esp{en primer lugar, además, sin embargo, por eso, en conclusión}) ; lexique NTIC correct (\esp{teléfono inteligente, redes sociales, brecha digital, descargar, navegar}) ; structures de la leçon employées (\esp{suelo usar}, \esp{prefiero estudiar}) ; indicatif après \esp{si sabemos} (fait général, certitude).

\textbf{Conclusion} : une rédaction réussie combine opinion claire, connecteurs, lexique du thème et au moins une structure grammaticale de la leçon.
\end{exoresolu}

\courssub{Point de vigilance : Erreurs à éviter au Bac}

\begin{attention}[Les 5 erreurs qui coûtent des points]
\begin{enumerate}
\item \textbf{Oublier le changement radical} : *\esp{prefero} au lieu de \esp{prefiero} ; *\esp{solo navegar} pour « j'ai l'habitude de naviguer » au lieu de \esp{suelo navegar}. Mais jamais de changement à \esp{nosotros} : \esp{preferimos}, \esp{solemos}.
\item \textbf{Calquer le français « préférer... que »} : on dit \esp{Prefiero el té al café} (préposition \esp{a}), et non *\esp{prefiero el té que el café}. En revanche, avec \esp{gustar más}, c'est bien \esp{que} : \esp{Me gusta más el té que el café}.
\item \textbf{Faux amis du lexique NTIC} : \esp{descargar} = « télécharger (depuis Internet) », pas « décharger un camion » ; \esp{la contraseña} = « le mot de passe » (et non *\esp{la contraseñia}) ; « une application » = \esp{una aplicación}, pas *\esp{una aplicación de la ley} au sens juridique.
\item \textbf{Fautes d'accents et de ponctuation} : \esp{teléfono, móvil, música, lección, conexión, opinión, cibercafé} prennent un accent ; toute question exige ¿...? et toute exclamation ¡...! : \esp{¿Prefieres el té?}
\item \textbf{Mauvais choix de temps pour l'habitude} : une habitude passée (\esp{antes, de niño, todos los días}) exige l'imparfait (\esp{jugaba, íbamos, solía}), jamais le passé simple ; réserve le passé simple (\esp{preferí, descargué}) à une action unique et terminée.
\end{enumerate}
\end{attention}

\begin{savaistu}[Culture numérique : Cameroun et Guinée équatoriale]
La \textbf{brecha digital} est une réalité majeure en zone CEMAC. Au Cameroun, le taux de pénétration d'Internet dépasse 40\% en ville (Yaoundé, Douala) mais chute sous les 10\% en zone rurale. La Guinée équatoriale, unique pays hispanophone d'Afrique subsaharienne, investit dans la fibre optique (projet \esp{GETESA}) mais le coût des données reste élevé. Les jeunes camerounais utilisent massivement \esp{WhatsApp} et \esp{Facebook} (appelés \esp{« las redes »}) pour l'information, le commerce informel (\esp{« business en ligne »}) et la préparation du Bac. Le \esp{cibercafé} reste un lieu social essentiel pour ceux qui n'ont pas de \esp{smartphone} personnel ou de \esp{conexión} fixe à domicile. \esp{¡Ojo!} : en Guinée équatoriale, on dit souvent \esp{ordenador} et \esp{móvil} (norme espagnole), mais l'influence française fait apparaître \esp{ordinateur} et \esp{portable} dans le parler mixte local.
\end{savaistu}

\newpage

\coursec{Exercices}

\courssub{Je vérifie mes connaissances}

\begin{exercice}[QCM : préférence et habitude]
Pour chaque phrase, choisis la bonne réponse (a, b ou c).

\begin{enumerate}
\item \esp{Yo \ldots\ WhatsApp a Facebook.} \\
a) \esp{prefiero} \quad b) \esp{prefiere} \quad c) \esp{prefero}
\item \esp{Mi hermano \ldots\ navegar por Internet cada noche.} \\
a) \esp{suele} \quad b) \esp{sole} \quad c) \esp{suelo}
\item \esp{Antes, nosotros \ldots\ menos tiempo en las redes sociales.} \\
a) \esp{pasamos} \quad b) \esp{pasábamos} \quad c) \esp{paseamos}
\item \esp{¿\ldots\ te gusta más, el teléfono inteligente o la tableta?} \\
a) \esp{Qué} \quad b) \esp{Cómo} \quad c) \esp{Cuánto}
\item \esp{Ella \ldots\ la costumbre de descargar música los fines de semana.} \\
a) \esp{tiene} \quad b) \esp{hace} \quad c) \esp{está}
\end{enumerate}
\end{exercice}

\begin{exercice}[Vrai ou faux ? Justifie]
Réponds par \esp{Verdadero} ou \esp{Falso} et justifie ta réponse par une phrase complète en espagnol.

\begin{enumerate}
\item \esp{« Preferir »} est un verbe régulier au présent de l'indicatif.
\item \esp{« Soler »} se conjugue comme \esp{volver} (changement \esp{o > ue}).
\item Pour exprimer une habitude au passé, on utilise l'imparfait de l'indicatif.
\item \esp{« Me gusta más el fútbol »} signifie « je n'aime pas le football ».
\item \esp{« Acostumbrar a »} est toujours suivi d'un infinitif.
\end{enumerate}
\end{exercice}

\begin{exercice}[Vocabulaire des NTIC]
Complète chaque phrase avec le mot qui convient, choisi dans la liste : \esp{contraseña}, \esp{red social}, \esp{descargar}, \esp{navegar}, \esp{brecha digital}, \esp{teléfono inteligente}.

\begin{enumerate}
\item Para entrar en mi cuenta, necesito mi \ldots\ldots\ldots
\item Con mi \ldots\ldots\ldots\ puedo llamar, sacar fotos y conectarme a Internet.
\item Me gusta \ldots\ldots\ldots\ canciones de artistas cameruneses.
\item Facebook es una \ldots\ldots\ldots\ muy popular entre los jóvenes.
\item En el cibercafé de Douala, los estudiantes \ldots\ldots\ldots\ por Internet para hacer sus deberes.
\item La \ldots\ldots\ldots\ separa a los que tienen acceso a Internet de los que no lo tienen.
\end{enumerate}
\end{exercice}

\begin{exercice}[Conjugaison à compléter]
Complète avec le verbe \esp{preferir} ou \esp{soler} au présent ou à l'imparfait, selon le sens.

\begin{enumerate}
\item Yo \ldots\ldots\ldots\ (preferir, présent) leer las noticias en el periódico digital.
\item Antes, nosotros \ldots\ldots\ldots\ (soler, imparfait) ir al cibercafé del barrio.
\item ¿Tú \ldots\ldots\ldots\ (preferir, présent) llamar o enviar mensajes?
\item De niño, él \ldots\ldots\ldots\ (soler, imparfait) jugar al fútbol después de clase.
\item Mis padres \ldots\ldots\ldots\ (preferir, présent) la radio a la televisión.
\item Cuando vivía en Yaoundé, yo \ldots\ldots\ldots\ (soler, imparfait) conectarme por la mañana.
\item Ellas \ldots\ldots\ldots\ (preferir, présent) estudiar con el ordenador.
\item Vosotros \ldots\ldots\ldots\ (soler, imparfait) ver películas en español, ¿verdad?
\item Mi profesor \ldots\ldots\ldots\ (preferir, présent) corregir en papel.
\item Antes, tú \ldots\ldots\ldots\ (soler, imparfait) descargar mucha música.
\end{enumerate}
\end{exercice}

\courssub{Je m'entraîne}

\begin{exercice}[Appariements : mots et définitions]
Associe chaque mot espagnol (1 à 10) à sa définition française (a à j), puis réemploie cinq de ces mots dans des phrases personnelles en espagnol.

\begin{enumerate}
\item \esp{la contraseña} \hspace{1cm} a) téléphone portable multifonction
\item \esp{la red social} \hspace{1cm} b) écart entre les personnes connectées et non connectées
\item \esp{descargar} \hspace{1cm} c) mot de passe
\item \esp{navegar} \hspace{1cm} d) transférer un fichier d'Internet vers son appareil
\item \esp{el teléfono inteligente} \hspace{1cm} e) plateforme de communication en ligne
\item \esp{la brecha digital} \hspace{1cm} f) se déplacer de page en page sur Internet
\item \esp{el usuario} \hspace{1cm} g) personne qui utilise un service
\item \esp{la pantalla} \hspace{1cm} h) surface où s'affichent les images
\item \esp{conectarse} \hspace{1cm} i) établir une liaison avec Internet
\item \esp{el mensaje} \hspace{1cm} j) texte court envoyé à quelqu'un
\end{enumerate}
\end{exercice}

\begin{exercice}[Texte à trous]
Complète le texte suivant avec les mots de la liste : \esp{suelo}, \esp{prefiero}, \esp{acostumbra}, \esp{solía}, \esp{me gusta más}, \esp{contraseña}, \esp{descargar}, \esp{navegar}.

\begin{textebox}[Mi vida digital]
Todos los días, después de las clases, yo \ldots\ldots\ldots\ conectarme a Internet con mi teléfono inteligente. \ldots\ldots\ldots\ WhatsApp a las llamadas, porque es más barato. Mi hermana, en cambio, \ldots\ldots\ldots\ a llamar directamente. Antes, cuando no teníamos teléfono, ella \ldots\ldots\ldots\ escribir cartas a sus amigas. Ahora, para entrar en mi cuenta, escribo mi \ldots\ldots\ldots\ y puedo \ldots\ldots\ldots\ por las redes sociales o \ldots\ldots\ldots\ música. A mí \ldots\ldots\ldots\ escuchar música que ver vídeos.
\end{textebox}
\end{exercice}

\begin{exercice}[Transformations]
Transforme chaque phrase selon la consigne, sans changer le sens.

\begin{enumerate}
\item \esp{Me gusta más el chocolate que el café.} $\rightarrow$ Reformule avec \esp{preferir\ldots a\ldots}.
\item \esp{Normalmente estudio por la noche.} $\rightarrow$ Reformule avec \esp{soler}.
\item \esp{Mi abuelo tiene la costumbre de leer el periódico.} $\rightarrow$ Reformule avec \esp{acostumbrar a}.
\item \esp{Antes íbamos al cibercafé todos los sábados.} $\rightarrow$ Reformule avec \esp{soler} à l'imparfait.
\item \esp{Prefiero el té al café.} $\rightarrow$ Reformule avec \esp{me gusta más}.
\end{enumerate}
\end{exercice}

\begin{exercice}[Phrases à traduire]
Traduis les phrases suivantes en espagnol.

\begin{enumerate}
\item Je préfère les réseaux sociaux aux journaux.
\item Mon frère a l'habitude de télécharger des films le week-end.
\item Avant, nous naviguions moins sur Internet.
\item Elle préfère appeler à envoyer des messages.
\item J'aime mieux étudier avec mon téléphone intelligent qu'avec les livres.
\item D'habitude, ils se connectent au cybercafé de Yaoundé.
\end{enumerate}
\end{exercice}

\courssub{Je me prépare au Bac}

\begin{exercice}[Compréhension écrite : texte et questions]
Lis le texte suivant, puis réponds aux questions en espagnol, avec des phrases complètes.

\begin{textebox}[Los jóvenes españoles y las redes sociales]
Hoy en día, la mayoría de los jóvenes españoles tiene un teléfono inteligente y pasa varias horas al día conectado a Internet. Según los profesores, los alumnos suelen consultar las redes sociales incluso durante las clases, lo que preocupa a las familias. Muchos jóvenes prefieren comunicarse por mensajes a hablar cara a cara, y acostumbran a compartir fotos y vídeos de su vida diaria. Antes, los adolescentes solían reunirse en las plazas para charlar; ahora, en cambio, se reúnen en grupos virtuales. Sin embargo, no todos tienen las mismas oportunidades : la brecha digital todavía existe entre las ciudades grandes y los pueblos pequeños, donde la conexión es lenta o cara. Algunas asociaciones organizan cursos gratuitos para enseñar a los jóvenes a usar Internet de forma responsable y proteger su contraseña y sus datos personales.
\end{textebox}

Questions :
\begin{enumerate}
\item Vrai ou faux ? Justifie par une citation du texte : \esp{Los jóvenes españoles nunca usan las redes sociales en clase.}
\item Vrai ou faux ? Justifie par une citation du texte : \esp{La brecha digital ha desaparecido en España.}
\item Cite du texte deux expressions qui montrent une habitude des jeunes.
\item \esp{¿Cómo se reunían los adolescentes antes, según el texto?}
\item Question personnelle : \esp{¿Prefieres comunicarte por mensajes o cara a cara? ¿Por qué?} (3 à 4 phrases)
\end{enumerate}
\end{exercice}

\begin{exercice}[Thème et version type Bac]
\textbf{A. Thème.} Traduis en espagnol le texte suivant :

« Je préfère WhatsApp à Facebook, mais j'ai l'habitude d'utiliser les deux. Avant, je naviguais moins sur Internet. Maintenant, mon frère et moi nous connectons tous les soirs avec notre téléphone intelligent. Ma mère préfère la radio aux réseaux sociaux. »

\textbf{B. Version.} Traduis en français le texte suivant :

\begin{textebox}[La abuela conectada]
Mi abuela Rosa tiene setenta años, pero le encanta la tecnología. Cuando era joven, solía escribir cartas a su familia y esperaba semanas para recibir una respuesta. Ahora acostumbra a enviar mensajes por WhatsApp todas las mañanas. Prefiere las videollamadas a las cartas, porque así puede ver a sus nietos que viven lejos. Dice que Internet le ha cambiado la vida.
\end{textebox}
\end{exercice}

\begin{exercice}[Expression écrite type Bac]
Sujet : \esp{La brecha digital en mi país}.

Tu es journaliste pour le journal de ton lycée. Rédige un court article en espagnol (environ 15 lignes, soit 120 à 150 mots) sur l'accès à Internet au Cameroun : qui a accès à Internet, où, avec quels appareils, et quelles différences existent entre la ville et le village.

Consignes obligatoires :
\begin{itemize}
\item Utilise au moins \textbf{trois structures de la leçon} : une expression de la préférence (\esp{preferir}, \esp{me gusta más}, \esp{prefiero\ldots a\ldots}) et deux expressions de l'habitude (\esp{soler}, \esp{acostumbrar a}, \esp{tener la costumbre de}, imparfait).
\item Emploie au moins \textbf{cinq mots du lexique des NTIC} étudiés dans la leçon.
\item Soigne la ponctuation espagnole (¿ ? ¡ !) et les accents.
\end{itemize}
\end{exercice}

\newpage

\coursec{Corrigés des exercices}

\begin{corrige}[Exercice 1 — QCM : préférence et habitude]
\begin{enumerate}
\item \textbf{a)} \esp{Yo prefiero WhatsApp a Facebook.} — Première personne du singulier de \esp{preferir} au présent : changement de voyelle \esp{e > ie} (\esp{prefiero}). « Je préfère WhatsApp à Facebook. »
\item \textbf{a)} \esp{Mi hermano suele navegar por Internet cada noche.} — Troisième personne du singulier de \esp{soler} : \esp{o > ue} donne \esp{suele}. « Mon frère a l'habitude de naviguer sur Internet chaque soir. »
\item \textbf{b)} \esp{Antes, nosotros pasábamos menos tiempo en las redes sociales.} — Habitude au passé (\esp{antes}) : imparfait \esp{pasábamos}. « Avant, nous passions moins de temps sur les réseaux sociaux. »
\item \textbf{a)} \esp{¿Qué te gusta más, el teléfono inteligente o la tableta?} — Pour demander un choix entre deux choses, on utilise \esp{¿Qué\ldots?} avec l'accent. « Qu'est-ce que tu aimes le plus, le téléphone intelligent ou la tablette ? »
\item \textbf{a)} \esp{Ella tiene la costumbre de descargar música los fines de semana.} — Expression figée : \esp{tener la costumbre de + infinitif}. « Elle a l'habitude de télécharger de la musique le week-end. »
\end{enumerate}
\end{corrige}

\begin{corrige}[Exercice 2 — Vrai ou faux ? Justifie]
\begin{enumerate}
\item \textbf{Falso.} \esp{« Preferir » no es un verbo regular : tiene un cambio de vocal e > ie en el presente (prefiero, prefieres, prefiere, prefieren).} — \esp{preferir} est irrégulier au présent (et \esp{e > i} au prétérit : \esp{prefirió}).
\item \textbf{Verdadero.} \esp{« Soler » se conjuga como « volver » : o > ue (suelo, sueles, suele, solemos, soléis, suelen).} — Même alternance que \esp{volver : vuelvo}.
\item \textbf{Verdadero.} \esp{Para expresar una costumbre en el pasado se usa el imperfecto : « Antes navegábamos menos ».} — L'imparfait exprime l'habitude passée, comme en français.
\item \textbf{Falso.} \esp{« Me gusta más el fútbol » significa « j'aime davantage le football » (comparatif de préférence), no una negación.} — \esp{más} marque ici la préférence, pas la négation.
\item \textbf{Verdadero.} \esp{« Acostumbrar a » va siempre seguido de un infinitivo : « Acostumbro a estudiar por la noche ».} — La préposition \esp{a} est obligatoire devant l'infinitif.
\end{enumerate}
\end{corrige}

\begin{corrige}[Exercice 3 — Vocabulaire des NTIC]
\begin{enumerate}
\item \esp{Para entrar en mi cuenta, necesito mi \textbf{contraseña}.} — « Pour entrer dans mon compte, j'ai besoin de mon mot de passe. »
\item \esp{Con mi \textbf{teléfono inteligente} puedo llamar, sacar fotos y conectarme a Internet.} — « Avec mon téléphone intelligent, je peux appeler, prendre des photos et me connecter à Internet. »
\item \esp{Me gusta \textbf{descargar} canciones de artistas cameruneses.} — « J'aime télécharger des chansons d'artistes camerounais. »
\item \esp{Facebook es una \textbf{red social} muy popular entre los jóvenes.} — « Facebook est un réseau social très populaire chez les jeunes. »
\item \esp{En el cibercafé de Douala, los estudiantes \textbf{navegan} por Internet para hacer sus deberes.} — « Au cybercafé de Douala, les élèves naviguent sur Internet pour faire leurs devoirs. »
\item \esp{La \textbf{brecha digital} separa a los que tienen acceso a Internet de los que no lo tienen.} — « La fracture numérique sépare ceux qui ont accès à Internet de ceux qui ne l'ont pas. »
\end{enumerate}
\end{corrige}

\begin{corrige}[Exercice 4 — Conjugaison à compléter]
\begin{enumerate}
\item \esp{Yo \textbf{prefiero} leer las noticias en el periódico digital.} — présent, 1\iere{} pers. sing. : \esp{e > ie}.
\item \esp{Antes, nosotros \textbf{solíamos} ir al cibercafé del barrio.} — imparfait de \esp{soler}, régulier : \esp{solía, solías, solía, solíamos, solíais, solían}.
\item \esp{¿Tú \textbf{prefieres} llamar o enviar mensajes?}
\item \esp{De niño, él \textbf{solía} jugar al fútbol después de clase.}
\item \esp{Mis padres \textbf{prefieren} la radio a la televisión.}
\item \esp{Cuando vivía en Yaoundé, yo \textbf{solía} conectarme por la mañana.}
\item \esp{Ellas \textbf{prefieren} estudiar con el ordenador.}
\item \esp{Vosotros \textbf{solíais} ver películas en español, ¿verdad?} — à l'imparfait, \esp{soler} n'a pas de changement de voyelle : \esp{solíais}.
\item \esp{Mi profesor \textbf{prefiere} corregir en papel.}
\item \esp{Antes, tú \textbf{solías} descargar mucha música.}
\end{enumerate}
\textbf{Point de vigilance :} le changement \esp{o > ue} de \esp{soler} n'existe qu'au présent ; à l'imparfait, la conjugaison est entièrement régulière.
\end{corrige}

\begin{corrige}[Exercice 5 — Appariements : mots et définitions]
\begin{center}
\begin{tabularx}{\linewidth}{|l|X|}\hline
\rowcolor{popblueL} Mot & Définition \\ \hline
1. \esp{la contraseña} & c) mot de passe \\ \hline
2. \esp{la red social} & e) plateforme de communication en ligne \\ \hline
3. \esp{descargar} & d) transférer un fichier d'Internet vers son appareil \\ \hline
4. \esp{navegar} & f) se déplacer de page en page sur Internet \\ \hline
5. \esp{el teléfono inteligente} & a) téléphone portable multifonction \\ \hline
6. \esp{la brecha digital} & b) écart entre les personnes connectées et non connectées \\ \hline
7. \esp{el usuario} & g) personne qui utilise un service \\ \hline
8. \esp{la pantalla} & h) surface où s'affichent les images \\ \hline
9. \esp{conectarse} & i) établir une liaison avec Internet \\ \hline
10. \esp{el mensaje} & j) texte court envoyé à quelqu'un \\ \hline
\end{tabularx}
\end{center}
\textbf{Exemples de phrases personnelles (réemploi) :}
\begin{itemize}
\item \esp{Nunca comparto mi contraseña con nadie.} « Je ne partage jamais mon mot de passe avec personne. »
\item \esp{Mi red social favorita es WhatsApp.} « Mon réseau social préféré est WhatsApp. »
\item \esp{Suelo descargar música los sábados.} « J'ai l'habitude de télécharger de la musique le samedi. »
\item \esp{Los usuarios del cibercafé pagan por hora.} « Les usagers du cybercafé paient à l'heure. »
\item \esp{Me conecto a Internet todas las noches.} « Je me connecte à Internet tous les soirs. »
\end{itemize}
\end{corrige}

\begin{corrige}[Exercice 6 — Texte à trous]
\begin{textebox}[Mi vida digital — texte complété]
Todos los días, después de las clases, yo \textbf{suelo} conectarme a Internet con mi teléfono inteligente. \textbf{Prefiero} WhatsApp a las llamadas, porque es más barato. Mi hermana, en cambio, \textbf{acostumbra} a llamar directamente. Antes, cuando no teníamos teléfono, ella \textbf{solía} escribir cartas a sus amigas. Ahora, para entrar en mi cuenta, escribo mi \textbf{contraseña} y puedo \textbf{navegar} por las redes sociales o \textbf{descargar} música. A mí \textbf{me gusta más} escuchar música que ver vídeos.
\end{textebox}
\textbf{Traduction :} « Tous les jours, après les cours, j'ai l'habitude de me connecter à Internet avec mon téléphone intelligent. Je préfère WhatsApp aux appels, parce que c'est moins cher. Ma sœur, en revanche, a l'habitude d'appeler directement. Avant, quand nous n'avions pas de téléphone, elle avait l'habitude d'écrire des lettres à ses amies. Maintenant, pour entrer dans mon compte, j'écris mon mot de passe et je peux naviguer sur les réseaux sociaux ou télécharger de la musique. Moi, j'aime mieux écouter de la musique que regarder des vidéos. »
\end{corrige}

\begin{corrige}[Exercice 7 — Transformations]
\begin{enumerate}
\item \esp{Me gusta más el chocolate que el café.} $\rightarrow$ \esp{\textbf{Prefiero el chocolate al café.}} — structure \esp{preferir X a Y} ; attention à la contraction \esp{a + el = al}.
\item \esp{Normalmente estudio por la noche.} $\rightarrow$ \esp{\textbf{Suelo estudiar por la noche.}} — \esp{soler + infinitif}, 1\iere{} pers. sing. : \esp{suelo}.
\item \esp{Mi abuelo tiene la costumbre de leer el periódico.} $\rightarrow$ \esp{\textbf{Mi abuelo acostumbra a leer el periódico.}} — \esp{acostumbrar a + infinitif}.
\item \esp{Antes íbamos al cibercafé todos los sábados.} $\rightarrow$ \esp{\textbf{Antes solíamos ir al cibercafé todos los sábados.}} — \esp{soler} à l'imparfait + infinitif.
\item \esp{Prefiero el té al café.} $\rightarrow$ \esp{\textbf{Me gusta más el té que el café.}} — \esp{gustar} + comparatif \esp{más\ldots que}.
\end{enumerate}
\textbf{Point de vigilance :} avec \esp{preferir}, la comparaison se fait avec \esp{a} (\esp{prefiero el té \emph{al} café}) ; avec \esp{gustar más}, elle se fait avec \esp{que} (\esp{me gusta más el té \emph{que} el café}). Ne pas mélanger les deux structures.
\end{corrige}

\begin{corrige}[Exercice 8 — Phrases à traduire]
\begin{enumerate}
\item « Je préfère les réseaux sociaux aux journaux. » $\rightarrow$ \esp{\textbf{Prefiero las redes sociales a los periódicos.}}
\item « Mon frère a l'habitude de télécharger des films le week-end. » $\rightarrow$ \esp{\textbf{Mi hermano tiene la costumbre de descargar películas los fines de semana.}} (ou : \esp{Mi hermano acostumbra a descargar\ldots} / \esp{suele descargar\ldots})
\item « Avant, nous naviguions moins sur Internet. » $\rightarrow$ \esp{\textbf{Antes, navegábamos menos por Internet.}} — habitude passée : imparfait.
\item « Elle préfère appeler à envoyer des messages. » $\rightarrow$ \esp{\textbf{Ella prefiere llamar a enviar mensajes.}} — devant deux infinitifs, la préposition \esp{a} se maintient.
\item « J'aime mieux étudier avec mon téléphone intelligent qu'avec les livres. » $\rightarrow$ \esp{\textbf{Me gusta más estudiar con mi teléfono inteligente que con los libros.}}
\item « D'habitude, ils se connectent au cybercafé de Yaoundé. » $\rightarrow$ \esp{\textbf{Normalmente, se conectan en el cibercafé de Yaoundé.}} (ou : \esp{Suelen conectarse\ldots})
\end{enumerate}
\textbf{Points de vigilance :} « les week-ends » se dit \esp{los fines de semana} ; « naviguer sur Internet » = \esp{navegar \emph{por} Internet} ; « au cybercafé » = \esp{en el cibercafé} (lieu : \esp{en}).
\end{corrige}

\begin{corrige}[Exercice 9 — Compréhension écrite : texte et questions]
\begin{enumerate}
\item \textbf{Falso.} \esp{El texto dice : « los alumnos suelen consultar las redes sociales incluso durante las clases ».} — La phrase affirme le contraire : les élèves consultent les réseaux sociaux même pendant les cours.
\item \textbf{Falso.} \esp{El texto dice : « la brecha digital todavía existe entre las ciudades grandes y los pueblos pequeños ».} — \esp{todavía existe} = « elle existe encore » ; la fracture numérique n'a donc pas disparu.
\item Deux expressions d'habitude : \esp{« los alumnos suelen consultar las redes sociales »} (\esp{soler + infinitif}) et \esp{« acostumbran a compartir fotos y vídeos »} (\esp{acostumbrar a + infinitif}). On peut aussi citer \esp{« solían reunirse en las plazas »} (habitude passée).
\item \esp{Antes, los adolescentes se reunían en las plazas para charlar.} — « Avant, les adolescents se retrouvaient sur les places pour bavarder. »
\item \textbf{Réponse personnelle (modèle) :} \esp{Yo prefiero comunicarme por mensajes porque es más rápido y más barato. Además, puedo escribir a mis amigos en cualquier momento del día. Sin embargo, para hablar de cosas importantes, me gusta más hablar cara a cara, porque así entiendo mejor los sentimientos de la otra persona.} — « Moi, je préfère communiquer par messages parce que c'est plus rapide et moins cher. De plus, je peux écrire à mes amis à tout moment de la journée. Cependant, pour parler de choses importantes, j'aime mieux parler face à face, car ainsi je comprends mieux les sentiments de l'autre personne. »
\end{enumerate}
\textbf{Barème indicatif (10 pts) :} questions 1 et 2 : 2 pts chacune (1 pt réponse + 1 pt citation exacte) ; question 3 : 2 pts ; question 4 : 2 pts ; question 5 : 2 pts (structures de préférence, correction, 3-4 phrases).
\end{corrige}

\begin{corrige}[Exercice 10 — Thème et version type Bac]
\textbf{A. Thème — proposition de traduction :}

\begin{textebox}[Traduction du thème]
Prefiero WhatsApp a Facebook, pero tengo la costumbre de usar los dos. Antes, navegaba menos por Internet. Ahora, mi hermano y yo nos conectamos todas las noches con nuestro teléfono inteligente. Mi madre prefiere la radio a las redes sociales.
\end{textebox}

\textbf{Points de vigilance (thème) :}
\begin{itemize}
\item « j'ai l'habitude d'utiliser » : \esp{tengo la costumbre de usar} ou \esp{suelo usar} ou \esp{acostumbro a usar}.
\item « je naviguais » : imparfait \esp{navegaba} (habitude passée), jamais le prétérit.
\item « nous nous connectons » : verbe pronominal \esp{nos conectamos} ; ne pas oublier le pronom réfléchi.
\item \esp{preferir X a Y} : \esp{prefiere la radio \emph{a} las redes sociales}.
\end{itemize}

\textbf{B. Version — proposition de traduction :}

« Ma grand-mère Rosa a soixante-dix ans, mais elle adore la technologie. Quand elle était jeune, elle avait l'habitude d'écrire des lettres à sa famille et elle attendait des semaines pour recevoir une réponse. Maintenant, elle a l'habitude d'envoyer des messages par WhatsApp tous les matins. Elle préfère les appels vidéo aux lettres, parce qu'ainsi elle peut voir ses petits-enfants qui vivent loin. Elle dit qu'Internet a changé sa vie. »

\textbf{Points de vigilance (version) :}
\begin{itemize}
\item \esp{le encanta la tecnología} : « elle adore la technologie » (\esp{encantar} se construit comme \esp{gustar}).
\item \esp{solía escribir} : imparfait = habitude passée, « elle avait l'habitude d'écrire ».
\item \esp{las videollamadas} : « les appels vidéo » ; \esp{le ha cambiado la vida} : « a changé sa vie » (possessif rendu par l'article défini).
\end{itemize}

\textbf{Barème indicatif (10 pts) :} thème : 5 pts (1 pt par phrase, moins 0,5 par faute de structure ou d'accord) ; version : 5 pts (fidélité du sens, rendu des temps et des structures d'habitude).
\end{corrige}

\begin{corrige}[Exercice 11 — Expression écrite type Bac]
\textbf{Proposition de rédaction modèle :}

\begin{textebox}[La brecha digital en mi país]
En Camerún, el acceso a Internet ha cambiado mucho en los últimos años. En las ciudades grandes como Yaoundé y Douala, la mayoría de los jóvenes tiene un teléfono inteligente y suele conectarse todos los días. Los estudiantes acostumbran a navegar por Internet para hacer sus deberes y descargar documentos. Muchos prefieren las redes sociales a los periódicos para informarse, porque es más rápido y más barato.

Sin embargo, en los pueblos la situación es diferente : la conexión es lenta o cara, y muchas familias no tienen ordenador. Antes, los habitantes de los pueblos solían ir al cibercafé de la ciudad más cercana, pero ahora algunos se conectan con el teléfono móvil. Esta diferencia entre la ciudad y el campo se llama brecha digital.

Para reducir esta brecha, el gobierno y las asociaciones deben instalar más cibercafés y enseñar a los jóvenes a usar Internet de forma responsable, protegiendo su contraseña y sus datos personales. Internet es una herramienta magnífica : ¡todos los cameruneses deben poder aprovecharla!
\end{textebox}

\textbf{Traduction du modèle :} « Au Cameroun, l'accès à Internet a beaucoup changé ces dernières années. Dans les grandes villes comme Yaoundé et Douala, la plupart des jeunes ont un téléphone intelligent et ont l'habitude de se connecter tous les jours. Les élèves ont l'habitude de naviguer sur Internet pour faire leurs devoirs et télécharger des documents. Beaucoup préfèrent les réseaux sociaux aux journaux pour s'informer, parce que c'est plus rapide et moins cher. Cependant, dans les villages, la situation est différente : la connexion est lente ou chère, et beaucoup de familles n'ont pas d'ordinateur. Avant, les habitants des villages avaient l'habitude d'aller au cybercafé de la ville la plus proche, mais maintenant certains se connectent avec le téléphone mobile. Cette différence entre la ville et la campagne s'appelle la fracture numérique. Pour réduire cette fracture, le gouvernement et les associations doivent installer plus de cybercafés et apprendre aux jeunes à utiliser Internet de façon responsable, en protégeant leur mot de passe et leurs données personnelles. Internet est un outil magnifique : tous les Camerounais doivent pouvoir en profiter ! »

\textbf{Structures de la leçon employées :} \esp{suele conectarse} (\esp{soler}), \esp{acostumbran a navegar} (\esp{acostumbrar a}), \esp{prefieren las redes sociales a los periódicos} (\esp{preferir\ldots a\ldots}), \esp{solían ir} (imparfait d'habitude).

\textbf{Lexique NTIC employé :} \esp{teléfono inteligente}, \esp{navegar}, \esp{descargar}, \esp{redes sociales}, \esp{cibercafé}, \esp{brecha digital}, \esp{contraseña}.

\textbf{Barème indicatif (10 pts) :} respect du sujet et de la longueur : 2 pts ; trois structures de la leçon correctement employées : 3 pts ; cinq mots du lexique NTIC : 2 pts ; correction grammaticale (accords, conjugaison, accents) : 2 pts ; ponctuation espagnole et présentation : 1 pt.

\textbf{Points de vigilance :} vérifier le changement de voyelle de \esp{preferir} (\esp{prefieren}, pas \esp{preferen}) ; ne pas oublier la préposition \esp{a} après \esp{acostumbrar} et dans \esp{preferir X a Y} ; employer l'imparfait (\esp{solían}) pour les habitudes passées ; soigner les accents (\esp{conexión}, \esp{información}, \esp{rápido}).
\end{corrige}

\newpage

\coursec{Fiche bilan}

\begin{popfigure}
\begin{tikzpicture}[mindmap, grow cyclic, every node/.style={concept, font=\small},
  root concept/.append style={concept color=popblue, font=\bfseries, minimum size=3.2cm},
  level 1/.append style={level distance=3.6cm, sibling angle=51.5},
  level 1 concept/.append style={minimum size=2.4cm, font=\footnotesize}]
\node[root concept]{Internet y las NTIC}
  child[concept color=poporange]{ node[align=center]{Léxico NTIC\\ \esp{red social, contraseña, descargar}} }
  child[concept color=popgreen]{ node[align=center]{Preferencia\\ \esp{preferir, me gusta más}} }
  child[concept color=poppurple]{ node[align=center]{Hábito\\ \esp{soler + inf., acostumbrar}} }
  child[concept color=poppink]{ node[align=center]{Imperfecto\\ descripción del pasado} }
  child[concept color=popgold]{ node[align=center]{Comentario\\ de texto} }
  child[concept color=popblue!60]{ node[align=center]{Traducción\\ thème / version} }
  child[concept color=popmuted]{ node[align=center]{Brecha digital\\ sociedad} };
\end{tikzpicture}
\legende{Carte mentale de la leçon : Internet et les NTIC, préférence et habitude.}
\end{popfigure}

\begin{bilan}
\textbf{Lexique clé des NTIC}

\begin{center}
\begin{tabularx}{\linewidth}{|X|X|}\hline
\rowcolor{popblueL} Español & Français \\ \hline
\esp{la red social} & le réseau social \\ \hline
\esp{la contraseña} & le mot de passe \\ \hline
\esp{descargar / subir} & télécharger (recevoir / envoyer) \\ \hline
\esp{navegar (por Internet)} & naviguer (sur Internet) \\ \hline
\esp{el teléfono inteligente} & le smartphone \\ \hline
\esp{la brecha digital} & la fracture numérique \\ \hline
\esp{estar conectado / desconectado} & être connecté / déconnecté \\ \hline
\esp{el usuario / la usuaria} & l'utilisateur / l'utilisatrice \\ \hline
\esp{enviar un mensaje} & envoyer un message \\ \hline
\esp{la pantalla / el teclado} & l'écran / le clavier \\ \hline
\end{tabularx}
\end{center}

\textbf{Grammaire 1 : l'expression de la préférence}

\begin{center}
\begin{tabularx}{\linewidth}{|X|X|}\hline
\rowcolor{popgreenL} Structure & Exemple \\ \hline
\esp{preferir (e $\to$ ie) + inf. / nom} & \esp{Prefiero leer a ver la tele.} (Je préfère lire plutôt que regarder la télé.) \\ \hline
\esp{preferir X a Y} & \esp{Prefiero el fútbol al baloncesto.} (Je préfère le foot au basket.) \\ \hline
\esp{me gusta(n) más X (que Y)} & \esp{Me gusta más WhatsApp que Facebook.} \\ \hline
\esp{antes que / más bien} & \esp{Antes que chatear, prefiero llamar.} \\ \hline
\end{tabularx}
\end{center}

\textbf{Grammaire 2 : l'expression de l'habitude}

\begin{center}
\begin{tabularx}{\linewidth}{|X|X|}\hline
\rowcolor{poppurpleL} Structure & Exemple \\ \hline
\esp{soler (o $\to$ ue) + infinitivo} & \esp{Suelo navegar por la noche.} (J'ai l'habitude de naviguer le soir.) \\ \hline
\esp{acostumbrar (a) + inf.} & \esp{Acostumbro a revisar mis mensajes.} \\ \hline
\esp{tener la costumbre de + inf.} & \esp{Tengo la costumbre de apagar el móvil en clase.} \\ \hline
Imparfait (habitude passée) & \esp{Antes, usaba el cibercafé todos los días.} \\ \hline
\end{tabularx}
\end{center}

\textbf{Méthodes express}
\begin{itemize}
\item \textbf{Commentaire de texte} : présenter (auteur, source, thème) $\to$ résumer les idées $\to$ analyser lexique et arguments $\to$ donner son avis argumenté (\esp{en mi opinión, me parece que...}).
\item \textbf{Traduction} : repérer la structure française (préférence ou habitude), choisir l'équivalent espagnol, vérifier la conjugaison (\esp{prefiero, sueles, acostumbran}) et l'accord.
\item \textbf{Verbe \esp{soler}} : toujours suivi de l'\cle{infinitif} ; jamais de subjonctif après \esp{preferir} quand le sujet est identique.
\end{itemize}
\end{bilan}

\begin{aretenir}[Checklist avant le Bac]
\begin{itemize}
\item[$\square$] Je connais le lexique des NTIC : \esp{red social, contraseña, descargar, navegar, teléfono inteligente, brecha digital}.
\item[$\square$] Je sais conjuguer \esp{preferir} (e $\to$ ie : \esp{prefiero, prefieres, prefiere...}).
\item[$\square$] J'emploie \esp{preferir X a Y} avec la préposition \esp{a}, jamais \esp{que}.
\item[$\square$] Je sais dire « j'aime mieux » : \esp{me gusta más...} avec accord singulier/pluriel (\esp{me gustan más las redes sociales}).
\item[$\square$] Je maîtrise \esp{soler + infinitif} (o $\to$ ue : \esp{suelo, sueles, suele...}).
\item[$\square$] Je connais les synonymes : \esp{acostumbrar a + inf.} et \esp{tener la costumbre de + inf.}.
\item[$\square$] J'utilise l'\cle{imparfait} pour les habitudes passées : \esp{antes navegaba en el cibercafé}.
\item[$\square$] Je n'oublie pas les accents : \esp{contraseña, teléfono, opinión, conexión}.
\item[$\square$] Je sais structurer un commentaire : présentation, résumé, analyse, opinion personnelle.
\item[$\square$] Je sais donner mon avis : \esp{en mi opinión, desde mi punto de vista, me parece que + indicativo}.
\item[$\square$] Je vérifie mes traductions : conjugaison, prépositions, accords.
\item[$\square$] Je peux parler de la \esp{brecha digital} au Cameroun et dans le monde hispanophone.
\end{itemize}
\end{aretenir}

\findecours

\end{document}
